% La diversité des ensembles biogéographiques — Géographie Tle Tech — Cours signé OnBuch+
% Programme officiel MINESEC. Compiler avec : tectonic N01-diversite-ensembles-biogeographiques.tex
\newcommand{\DOCMATIERE}{Géographie}
\newcommand{\DOCNIVEAU}{Tle Tech}
\newcommand{\DOCMODULE}{Géographie – classe de Terminale technique}
\newcommand{\DOCLECON}{N01}
\newcommand{\DOCTITRE}{La diversité des ensembles biogéographiques}
% OnBuch+ — préambule « cours signés » ENRICHI (compilé avec tectonic / XeLaTeX)
% Système visuel « Pop » d'OnBuch+ : fond crème (jamais blanc), bordures encre
% épaisses, ombres « dures » décalées, titres Archivo Black en capitales.
% Version MATHÉMATIQUES : graphes (pgfplots/tikz), boîtes pédagogiques
% (définition, propriété, méthode, attention, le savais-tu, exercice résolu,
% exercice, TP, bilan), page de garde et signature OnBuch+ ancrée en bas.
%
% Macros attendues dans main.tex AVANT \input{preamble} :
%   \DOCMATIERE  \DOCNIVEAU  \DOCMODULE  \DOCLECON (numéro)  \DOCTITRE
\documentclass[11pt]{article}

\usepackage{fontspec}
\usepackage[french]{babel}
\usepackage[a4paper,margin=2cm,top=3.4cm,bottom=2.8cm,headheight=1.4cm,footskip=1.3cm]{geometry}
\usepackage{amsmath,amssymb,amsfonts}
\usepackage{mathtools} % \xmapsto, \coloneqq... souvent utilisés par les modèles
\usepackage{cancel} % simplifications barrées (bilans de forces)
% \abs{z} (module, valeur absolue) et \norm{u} (norme), utilisés par les modèles
\providecommand{\abs}{}\let\abs\relax\DeclarePairedDelimiter{\abs}{\lvert}{\rvert}
\providecommand{\norm}{}\let\norm\relax\DeclarePairedDelimiter{\norm}{\lVert}{\rVert}
\usepackage[table]{xcolor} % [table] : fournit aussi \rowcolors (alternance de lignes)
\usepackage{pagecolor}
\usepackage{tikz}
\usetikzlibrary{arrows.meta,positioning,calc,shapes.geometric,shapes.misc,shapes.arrows,decorations.pathmorphing,decorations.pathreplacing,decorations.markings,patterns,shadows,mindmap,trees,fit,backgrounds,matrix,intersections,angles,quotes}
\usepackage{pgfplots}
\usepackage[version=4]{mhchem} % équations nucléaires \ce{^{235}_{92}U}
\usepackage[european,siunitx]{circuitikz} % schémas électriques
\pgfplotsset{compat=1.18}
\usepgfplotslibrary{fillbetween,groupplots} % aires entre courbes (calcul intégral)
\usepackage{enumitem}
\usepackage{fancyhdr}
% [most] charge toutes les bibliothèques tcolorbox (listings, minted, raster,
% documentation...) et gonfle inutilement la mémoire TeX sur un document long
% (~300 boîtes) : on ne charge que ce qui est réellement utilisé.
\usepackage[skins,breakable]{tcolorbox}
\usepackage{graphicx}
\usepackage{colortbl}
\usepackage{booktabs}
\usepackage{tabularx}
\usepackage{multirow}
\usepackage{diagbox} % case d'en-tête diagonale des tableaux à double entrée
% Colonnes extensibles centrée / gauche / droite (C, L, R), souvent utilisées par les modèles
\newcolumntype{C}{>{\centering\arraybackslash}X}
\newcolumntype{L}{>{\raggedright\arraybackslash}X}
\newcolumntype{R}{>{\raggedleft\arraybackslash}X}
\usepackage{array}
\usepackage{multicol}
\usepackage{siunitx}
% parse-numbers=false : accepte aussi \SI{2,5 \times 10^{-3}}{...}
\sisetup{locale=FR,output-decimal-marker={,},group-minimum-digits=4,parse-numbers=false}
\DeclareSIUnit\ohm{\ensuremath{\symup{\Omega}}} % U+2126 absent de la police
\DeclareSIUnit\division{div} % réglages d'oscilloscope (ms/div, V/div)
\DeclareSIUnit\tr{tr} % tours (vitesse de rotation, ex. \SI{1500}{\tr\per\minute})
\DeclareSIUnit\molar{M} % concentration molaire (ex. \SI{3}{\molar}, \SI{1}{\micro\molar})
\DeclareSIUnit\an{an} % année (le modèle confond parfois avec \year, un compteur TeX) — ex. \SI{5}{\kilo\meter\per\an}
\DeclareSIUnit\FCFA{FCFA} % franc CFA (ex. \SI{500}{\FCFA\per\kilo\gram})
\DeclareSIUnit\degreeBrix{\textdegree Brix} % teneur en sucre d'un sirop/jus (réfractométrie)
\DeclareSIUnit\month{mois} % le modèle utilise parfois \month, un compteur TeX, comme unité de durée
\DeclareSIUnit\microliter{\micro\liter} % le modèle écrit parfois \microliter en un seul token
\DeclareSIUnit\million{millions} % dénombrement (ex. \SI{15}{\million\per\mL} spermatozoïdes)
\usepackage{qrcode}
% \degree : commande fréquemment utilisée par les modèles pour un angle ou une
% température en dehors de \SI{}{} (non fournie par les paquets chargés).
\providecommand{\degree}{^{\circ}}
% \qed (fin de démonstration), utilisé par les modèles sans amsthm
\providecommand{\qed}{\hfill$\square$}
% \barre : double barre des valeurs interdites dans un tableau de variations (tkz-tab)
\providecommand{\barre}{\ensuremath{\|}}
% environnement proof (amsthm non chargé) utilisé par les modèles
\ifdefined\proof\else\newenvironment{proof}[1][Démonstration]{\par\noindent\textit{#1.}\ }{\qed\par}\fi
\usepackage{lastpage}
\usepackage{hyperref}
\hypersetup{hidelinks}

% ============================================================
% Palette « Pop » OnBuch+ — mêmes hex que la classe P (plus_ui.dart)
% ============================================================
\definecolor{poporange}{HTML}{F59321}
\definecolor{popdark}{HTML}{E07A0C}
\definecolor{popcream}{HTML}{FFF6E8}
\definecolor{popink}{HTML}{1C1714}
\definecolor{poppurple}{HTML}{7A5AE0}
\definecolor{popgreen}{HTML}{2E9E6A}
\definecolor{poppink}{HTML}{E0457B}
\definecolor{popblue}{HTML}{2F7DE1}
\definecolor{popgold}{HTML}{C98A12}
\definecolor{popmuted}{HTML}{8A7F76}
\definecolor{popline}{HTML}{EADFD2}
\definecolor{popgray}{HTML}{8A7F76}
\colorlet{popgrey}{popgray}
% Alias tolérants (le modèle nomme parfois les couleurs différemment)
\colorlet{popwhite}{white}
\colorlet{popblack}{popink}
\colorlet{popred}{poppink}
\colorlet{popyellow}{popgold}
% Teintes claires (fonds de tableaux/encadrés) — le modèle nomme parfois
% les couleurs avec un suffixe L (ex. popblueL) sans les avoir définies.
\colorlet{poporangeL}{poporange!20}
\colorlet{popdarkL}{popdark!20}
\colorlet{poppurpleL}{poppurple!20}
\colorlet{popgreenL}{popgreen!20}
\colorlet{poppinkL}{poppink!20}
\colorlet{popblueL}{popblue!20}
\colorlet{popgoldL}{popgold!20}
\colorlet{popredL}{popred!20}
\colorlet{popgrayL}{popgray!20}
% Teintes claires pour fonds de tableaux / zones
\colorlet{popblueL}{popblue!10!white}
\colorlet{poporangeL}{poporange!14!white}
\colorlet{popgreenL}{popgreen!12!white}
\colorlet{poppurpleL}{poppurple!12!white}
\colorlet{poppinkL}{poppink!12!white}
\colorlet{popgoldL}{popgold!14!white}

\pagecolor{popcream}

% ============================================================
% Typographie
% ============================================================
\defaultfontfeatures{Ligatures=TeX}
\setmainfont{PlusJakartaSans-Medium.ttf}[Path=../../../fonts/,BoldFont=PlusJakartaSans-Bold.ttf]
% Maths en sans-serif (Fira Math) avec chiffres et lettres droites de Plus
% Jakarta, pour que les formules se fondent dans le texte.
\usepackage[math-style=ISO,bold-style=ISO]{unicode-math}
\setmathfont{FiraMath-Regular.otf}[Path=../../../fonts/]
\setmathfont{PlusJakartaSans-Medium.ttf}[Path=../../../fonts/,range={up/num,up/{Latin,latin},\mathup}]
\usepackage{icomma} % virgule décimale sans espace en mode math
% Caractères mathématiques Unicode absents des polices de texte (Plus Jakarta,
% Archivo Black) : rendus par leur équivalent mathématique, en texte comme en
% formule — sinon ils s'affichent en carré vide (ex. « Arithmétique dans ℤ »).
\usepackage{newunicodechar}
\newunicodechar{ℝ}{\ensuremath{\mathbb{R}}}
\newunicodechar{ℤ}{\ensuremath{\mathbb{Z}}}
\newunicodechar{ℂ}{\ensuremath{\mathbb{C}}}
\newunicodechar{ℕ}{\ensuremath{\mathbb{N}}}
\newunicodechar{ℚ}{\ensuremath{\mathbb{Q}}}
\newunicodechar{𝔻}{\ensuremath{\mathbb{D}}}
\newunicodechar{α}{\ensuremath{\alpha}}
\newunicodechar{β}{\ensuremath{\beta}}
\newunicodechar{γ}{\ensuremath{\gamma}}
\newunicodechar{δ}{\ensuremath{\delta}}
\newunicodechar{ε}{\ensuremath{\varepsilon}}
\newunicodechar{θ}{\ensuremath{\theta}}
\newunicodechar{λ}{\ensuremath{\lambda}}
\newunicodechar{μ}{\ensuremath{\mu}}
\newunicodechar{σ}{\ensuremath{\sigma}}
\newunicodechar{τ}{\ensuremath{\tau}}
\newunicodechar{φ}{\ensuremath{\varphi}}
\newunicodechar{ψ}{\ensuremath{\psi}}
\newunicodechar{ω}{\ensuremath{\omega}}
\newunicodechar{Σ}{\ensuremath{\Sigma}}
% majuscules produites par \MakeUppercase dans les titres (α-aminés -> Α)
\newunicodechar{Α}{\ensuremath{\alpha}}
\newunicodechar{Β}{\ensuremath{\beta}}
\newunicodechar{Γ}{\ensuremath{\Gamma}}
\newunicodechar{Δ}{\ensuremath{\Delta}}
\newunicodechar{Ε}{\ensuremath{\varepsilon}}
\newunicodechar{Θ}{\ensuremath{\Theta}}
\newunicodechar{Λ}{\ensuremath{\Lambda}}
\newunicodechar{Μ}{\ensuremath{\mu}}
\newunicodechar{Τ}{\ensuremath{\tau}}
\newunicodechar{Φ}{\ensuremath{\Phi}}
\newunicodechar{Ψ}{\ensuremath{\Psi}}
\newunicodechar{Ω}{\ensuremath{\Omega}}
\newunicodechar{↦}{\ensuremath{\mapsto}}
\newunicodechar{∈}{\ensuremath{\in}}
\newunicodechar{∉}{\ensuremath{\notin}}
\newunicodechar{∧}{\ensuremath{\wedge}}
\newunicodechar{⇔}{\ensuremath{\Leftrightarrow}}
\newunicodechar{⟺}{\ensuremath{\Longleftrightarrow}}
\newunicodechar{⇒}{\ensuremath{\Rightarrow}}
\newunicodechar{⟹}{\ensuremath{\Longrightarrow}}
\newunicodechar{∘}{\ensuremath{\circ}}
\newunicodechar{∪}{\ensuremath{\cup}}
\newunicodechar{∩}{\ensuremath{\cap}}
\newunicodechar{≡}{\ensuremath{\equiv}}
\newunicodechar{⊂}{\ensuremath{\subset}}
\newunicodechar{⊥}{\ensuremath{\perp}}
\newunicodechar{∃}{\ensuremath{\exists}}
\newunicodechar{∀}{\ensuremath{\forall}}
\newunicodechar{∛}{\ensuremath{\sqrt[3]{\,}}}
\newunicodechar{∜}{\ensuremath{\sqrt[4]{\,}}}
\newunicodechar{⃗}{\ensuremath{{}^{\to}}}
\newunicodechar{ⁿ}{\ifmmode^{n}\else\textsuperscript{\ensuremath{n}}\fi}
\newunicodechar{ˣ}{\ifmmode^{x}\else\textsuperscript{\ensuremath{x}}\fi}
\newunicodechar{ᵇ}{\ifmmode^{b}\else\textsuperscript{\ensuremath{b}}\fi}
\newunicodechar{ʳ}{\ifmmode^{r}\else\textsuperscript{\ensuremath{r}}\fi}
\newunicodechar{ᵘ}{\ifmmode^{u}\else\textsuperscript{\ensuremath{u}}\fi}
\newunicodechar{ʸ}{\ifmmode^{y}\else\textsuperscript{\ensuremath{y}}\fi}
\newunicodechar{ᵃ}{\ifmmode^{a}\else\textsuperscript{\ensuremath{a}}\fi}
\newunicodechar{ᵉ}{\ifmmode^{e}\else\textsuperscript{\ensuremath{e}}\fi}
\newunicodechar{ᵏ}{\ifmmode^{k}\else\textsuperscript{\ensuremath{k}}\fi}
\newunicodechar{ᵖ}{\ifmmode^{p}\else\textsuperscript{\ensuremath{p}}\fi}
\newunicodechar{ᴺ}{\ifmmode^{N}\else\textsuperscript{\ensuremath{N}}\fi}
\newunicodechar{ᶜ}{\ifmmode^{c}\else\textsuperscript{\ensuremath{c}}\fi}
\newunicodechar{ʰ}{\ifmmode^{h}\else\textsuperscript{\ensuremath{h}}\fi}
\newunicodechar{ᵠ}{\ifmmode^{q}\else\textsuperscript{\ensuremath{q}}\fi}
\newunicodechar{ₙ}{\ifmmode_{n}\else\textsubscript{\ensuremath{n}}\fi}
\newunicodechar{ₐ}{\ifmmode_{a}\else\textsubscript{\ensuremath{a}}\fi}
\newunicodechar{ᵢ}{\ifmmode_{i}\else\textsubscript{\ensuremath{i}}\fi}
\newunicodechar{ᵣ}{\ifmmode_{r}\else\textsubscript{\ensuremath{r}}\fi}
\newunicodechar{ₖ}{\ifmmode_{k}\else\textsubscript{\ensuremath{k}}\fi}
\newunicodechar{ᵦ}{\ifmmode_{\beta}\else\textsubscript{\ensuremath{\beta}}\fi}
\newunicodechar{ₓ}{\ifmmode_{x}\else\textsubscript{\ensuremath{x}}\fi}
\newfontfamily\popdisplay{ArchivoBlack-Regular.ttf}[Path=../../../fonts/]
\newfontfamily\poplogo{SpaceGrotesk-Bold.ttf}[Path=../../../fonts/]
\newfontfamily\popsemi{PlusJakartaSans-SemiBold.ttf}[Path=../../../fonts/]
\newfontfamily\popxbold{PlusJakartaSans-ExtraBold.ttf}[Path=../../../fonts/]
\newfontfamily\popxboldls{PlusJakartaSans-ExtraBold.ttf}[Path=../../../fonts/,LetterSpace=3.0]

\setlist[enumerate]{leftmargin=1.7em,itemsep=5pt,topsep=4pt}
\setlist[itemize]{leftmargin=1.7em,itemsep=4pt,topsep=4pt,label={\color{poporange}\textbullet}}
\setlength{\parindent}{0pt}
\setlength{\parskip}{5pt}
\renewcommand{\arraystretch}{1.25}

% pgfplots : style Pop par défaut
\pgfplotsset{
  every axis/.append style={axis line style={popink,line width=1pt},tick style={popink},
    grid style={popline,line width=0.5pt},label style={font=\small},tick label style={font=\footnotesize}},
  popcurve/.style={poporange,line width=1.8pt,no markers,smooth},
}
% Même style disponible dans un \draw TikZ simple (hors pgfplots)
\tikzset{popcurve/.style={poporange,line width=1.8pt}}
\tikzset{popgrid/.style={popline,very thin}}
\colorlet{popcurve}{poporange} % le nom du style est parfois utilisé comme couleur

% ============================================================
% Pill / squiggle
% ============================================================
\newcommand{\poppill}[3]{%
  \tikz[baseline=-0.55ex]\node[draw=popink,line width=1.2pt,rounded corners=2.6pt,%
    fill=#2,inner xsep=6.5pt,inner ysep=3pt]{%
    \popxboldls\fontsize{7}{7}\selectfont\color{#3}\MakeUppercase{#1}};%
}
\newcommand{\popsquiggle}[1]{%
  \tikz[baseline=-0.4ex]{
    \draw[#1,line width=2.6pt,line cap=round]
      plot [smooth] coordinates {(0,0.09) (0.34,0.01) (0.68,0.21) (1.02,0.01) (1.36,0.10)};
  }%
}
% Mot-clé surligné (marqueur orange sous le texte)
\newcommand{\cle}[1]{{\popxbold\color{popdark}#1}}

% ============================================================
% En-tête / pied
% ============================================================
\pagestyle{fancy}
\fancyhf{}
\renewcommand{\headrulewidth}{0pt}
\renewcommand{\footrulewidth}{0pt}
\lhead{\raisebox{-0.15ex}{\poplogo\fontsize{13}{13}\selectfont\color{popink}OnBuch\color{poporange}+}}
\rhead{\poppill{\DOCMATIERE\ \textbullet\ \DOCNIVEAU\ \textbullet\ Leçon \DOCLECON}{white}{popink}}
\lfoot{\popxbold\fontsize{7.6}{7.6}\selectfont\color{popmuted}\MakeUppercase{Cours signé OnBuch+}\ \textbullet\ {\popsemi\DOCTITRE}}
\rfoot{\tikz[baseline=-0.6ex]\node[rounded corners=8.5pt,fill=popink,minimum height=17pt,minimum width=17pt,inner xsep=5pt,inner sep=0]{\popxbold\fontsize{8}{8}\selectfont\color{popcream}\thepage\,/\,\pageref*{LastPage}};}

% ============================================================
% Titres
% ============================================================
\newcommand{\chaptitre}[2]{%
  \begin{tcolorbox}[enhanced,colback=poporange,colframe=popink,boxrule=2.6pt,arc=11pt,
      drop shadow=popink,left=15pt,right=15pt,top=12pt,bottom=11pt,before skip=3pt,after skip=18pt]
    \poppill{#2}{popink}{popcream}\par\vspace{9pt}
    {\raggedright\hyphenpenalty=10000\exhyphenpenalty=10000\popdisplay\fontsize{22}{24}\selectfont\color{popcream}\MakeUppercase{#1}\par}
  \end{tcolorbox}
}
% Section numérotée : pastille orange avec le numéro + titre Archivo
\newcounter{popsec}
\newcommand{\coursec}[1]{%
  \stepcounter{popsec}\setcounter{popsub}{0}%
  \par\vspace{16pt}\noindent
  \begin{minipage}{\linewidth}
  \tikz[baseline=-0.7ex]\node[draw=popink,line width=1.6pt,fill=poporange,rounded corners=4pt,minimum size=22pt,inner sep=0]{\popdisplay\fontsize{12}{12}\selectfont\color{popcream}\thepopsec};%
  \hspace{8pt}{\popdisplay\fontsize{15}{17}\selectfont\color{popink}\MakeUppercase{#1}}\par
  \vspace{4pt}\noindent{\color{poporange}\rule{36pt}{3.6pt}}
  \end{minipage}\par\nopagebreak\vspace{8pt}
}
\newcounter{popsub}
\newcommand{\courssub}[1]{%
  \stepcounter{popsub}%
  \par\vspace{9pt}\noindent{\popxbold\fontsize{11.5}{13}\selectfont\color{popink}{\color{poporange}\thepopsec.\thepopsub}\hspace{6pt}#1}\par\nopagebreak\vspace{4pt}
}

% ============================================================
% Boîtes pédagogiques — même recette : fond blanc, bordure épaisse colorée,
% ombre dure encre, badge plein. Titre optionnel [#1].
% ============================================================
\tcbset{popbase/.style={breakable,enhanced,colback=white,boxrule=2.3pt,arc=9pt,
  shadow={3pt}{-3pt}{0pt}{popink},left=14pt,right=14pt,top=11pt,bottom=11pt,before skip=11pt,after skip=15pt,
  fontupper=\popsemi\fontsize{10.6}{14.5}\selectfont\color{popink},
  fontlower=\popsemi\fontsize{10.6}{14.5}\selectfont\color{popink}}}
% Coche dessinée (le glyphe ✓ n'existe pas dans les polices du document)
\newcommand{\popcheck}[1]{\tikz[baseline=-0.5ex]\draw[#1,line width=1.6pt,line cap=round,line join=round](-0.13,0.0)--(-0.04,-0.09)--(0.13,0.1);}
\providecommand{\coche}{\popcheck{popgreen}}
\providecommand{\resultat}[1]{\par\smallskip\noindent\fcolorbox{popgreen}{popgreenL}{\parbox{\dimexpr\linewidth-2\fboxsep-2\fboxrule\relax}{\textbf{Résultat.} #1}}\par\smallskip}
\newcommand{\popboxhead}[3]{\poppill{#1}{#2}{white}\ifx\relax#3\relax\else\hspace{6pt}{\popxbold\fontsize{10.6}{12}\selectfont\color{popink}#3}\fi\par\vspace{7pt}}

\newtcolorbox{aretenir}[1][]{popbase,colframe=popgreen,before upper={\popboxhead{À retenir}{popgreen}{#1}}}
\newtcolorbox{exemplebox}[1][]{popbase,colframe=poporange,before upper={\popboxhead{Exemple}{poporange}{#1}}}
\newtcolorbox{definition}[1][]{popbase,colframe=popblue,before upper={\popboxhead{Définition}{popblue}{#1}}}
\newtcolorbox{propriete}[1][]{popbase,colframe=poppurple,before upper={\popboxhead{Propriété}{poppurple}{#1}}}
\newtcolorbox{methode}[1][]{popbase,colframe=popgold,before upper={\popboxhead{Méthode}{popgold}{#1}}}
\newtcolorbox{attention}[1][]{popbase,colframe=poppink,before upper={\popboxhead{Attention}{poppink}{#1}}}
\newtcolorbox{experience}[1][]{popbase,colframe=popblue,colback=popblueL,before upper={\popboxhead{Expérience}{popblue}{#1}}}
% « Le savais-tu ? » : carte violette pleine (contexte, culture, Cameroun)
\newtcolorbox{savaistu}[1][]{popbase,colback=poppurple,colframe=popink,
  fontupper=\popsemi\fontsize{10.4}{14.2}\selectfont\color{white},
  before upper={\poppill{Le savais-tu\,?}{white}{poppurple}\ifx\relax#1\relax\else\hspace{6pt}{\popxbold\color{white}#1}\fi\par\vspace{7pt}}}
% Exercice résolu : énoncé en haut, \tcblower puis la solution sur fond crème
\newcounter{popexo}\newcounter{popexr}
\newtcolorbox{exoresolu}[1][]{popbase,colframe=poporange,colbacklower=popcream,
  segmentation style={popink,line width=1.2pt,dashed},
  before upper={\stepcounter{popexr}\popboxhead{Exercice résolu \thepopexr}{poporange}{#1}},
  before lower={\poppill{Solution}{popink}{popcream}\par\vspace{6pt}}}
% Exercice (non résolu en place) — la correction est dans la partie Corrigés
\newtcolorbox{exercice}[1][]{popbase,colframe=popink,
  before upper={\stepcounter{popexo}\popboxhead{Exercice \thepopexo}{popink}{#1}}}
% Corrigé
\newtcolorbox{corrige}[1][]{popbase,colframe=popgreen,colback=popgreenL,
  before upper={\popboxhead{Corrigé}{popgreen}{#1}}}
% Objectifs / prérequis / bilan
\newtcolorbox{objectifs}{popbase,colframe=popink,colback=poporangeL,
  before upper={\popboxhead{Objectifs de la leçon}{popink}{}}}
\newtcolorbox{prerequis}{popbase,colframe=popink,colback=popblueL,
  before upper={\popboxhead{Prérequis}{popblue}{}}}
\newtcolorbox{bilan}{popbase,colframe=popink,colback=white,boxrule=2.8pt,
  before upper={\popboxhead{Fiche bilan}{poporange}{L'essentiel à connaître pour l'examen}}}
% Figure encadrée avec légende
\newtcolorbox{popfigure}[1][]{enhanced,breakable=false,colback=white,colframe=popline,boxrule=1.4pt,arc=8pt,
  left=10pt,right=10pt,top=10pt,bottom=8pt,before skip=10pt,after skip=12pt,halign=center,
  lower separated=false,halign lower=center,
  fontlower=\popsemi\fontsize{9}{11.5}\selectfont\color{popmuted},
  #1}
\newcommand{\legende}[1]{\par\vspace{6pt}{\popsemi\fontsize{9}{11.5}\selectfont\color{popmuted}#1}}

% ============================================================
% Signature OnBuch+
% ============================================================
\newcommand{\poplogoinline}[1]{{\poplogo\fontsize{#1}{#1}\selectfont\color{popink}OnBuch\color{poporange}+}}
\newcommand{\signatureonbuch}{%
  \begin{tcolorbox}[enhanced,colback=popink,colframe=popink,boxrule=2.2pt,arc=10pt,
      drop shadow=poporange,left=14pt,right=14pt,top=10pt,bottom=10pt,before skip=0pt,after skip=0pt]
    \tikz[baseline=-0.7ex]\node[circle,fill=popgreen,draw=popcream,line width=1.3pt,minimum size=22pt,inner sep=0]{\popcheck{white}};%
    \hspace{9pt}\begin{minipage}[c]{0.62\linewidth}
      {\popxbold\fontsize{12}{13}\selectfont\color{popcream}Signé\ \textbullet\ {\poplogo OnBuch\color{poporange}+}}\par\vspace{2pt}
      {\raggedright\popsemi\fontsize{8}{10.5}\selectfont\color{popline}\MakeUppercase{Équipe pédagogique OnBuch+}\\\MakeUppercase{Conforme au programme officiel MINESEC}\par}
    \end{minipage}\hfill
    {\poplogo\fontsize{22}{22}\selectfont\color{popcream}OnBuch\color{poporange}+}
  \end{tcolorbox}%
}

% ============================================================
% Page de garde
% #1 titre · #2 matière · #3 niveau · #4 module · #5 n° leçon · #6 durée ·
% #7 description / objectifs (texte ou itemize)
% ============================================================
\newcommand{\pagegarde}[7]{%
  \thispagestyle{empty}
  \newgeometry{a4paper,margin=1.7cm,top=1.6cm,bottom=1.4cm}%
  \begin{tikzpicture}[remember picture,overlay]
    \fill[poporange!18] ([xshift=-3.2cm,yshift=-3.6cm]current page.north east) circle (4.6cm);
    \fill[poppurple!14] ([xshift=2.4cm,yshift=7.6cm]current page.south west) circle (3.8cm);
  \end{tikzpicture}%
  {\poplogo\fontsize{30}{30}\selectfont\color{popink}OnBuch\color{poporange}+}\hfill
  \raisebox{0.6ex}{\poppill{Cours signé \textbullet\ Premium}{popink}{popcream}}\par
  \vspace{1pt}\popsquiggle{poppurple}\par
  \vspace{8pt}
  \begin{tcolorbox}[enhanced,colback=poporange,colframe=popink,boxrule=2.8pt,arc=13pt,
      drop shadow=popink,left=18pt,right=18pt,top=12pt,bottom=13pt,after skip=10pt]
    \poppill{#2 \textbullet\ #3}{popink}{popcream}\hspace{5pt}\poppill{#4}{white}{popink}\par\vspace{8pt}
    {\popdisplay\fontsize{14}{14}\selectfont\color{popink}LEÇON #5}\par\vspace{4pt}
    {\raggedright\hyphenpenalty=10000\exhyphenpenalty=10000\popdisplay\fontsize{25}{27}\selectfont\color{popcream}\MakeUppercase{#1}\par}
    \vspace{8pt}
    {\popxbold\fontsize{9.5}{9.5}\selectfont\color{popink}\MakeUppercase{Durée conseillée : #6}}
  \end{tcolorbox}
  \begin{tcolorbox}[enhanced,colback=white,colframe=popink,boxrule=2.2pt,arc=10pt,
      drop shadow=popink,left=16pt,right=16pt,top=10pt,bottom=10pt,after skip=8pt]
    \poppill{Dans cette leçon}{poppurple}{white}\par\vspace{5pt}
    \popsemi\fontsize{9.3}{12.6}\selectfont\color{popink}#7
  \end{tcolorbox}
  \noindent\begin{minipage}[t]{0.487\linewidth}
    \begin{tcolorbox}[enhanced,colback=popgreen,colframe=popink,boxrule=2.2pt,arc=10pt,
        drop shadow=popink,left=11pt,right=11pt,top=10pt,bottom=10pt,height=3.1cm,valign=center]
      \tcbox[colback=white,colframe=popink,boxrule=1.3pt,arc=5pt,boxsep=0pt,nobeforeafter,
        left=3pt,right=3pt,top=3pt,bottom=3pt,valign=center]{\qrcode[height=1.9cm]{https://play.google.com/store/apps/details?id=com.luvvix.onbuch&hl=fr}}%
      \hspace{8pt}\begin{minipage}[c]{0.5\linewidth}
        {\popdisplay\fontsize{11}{12.5}\selectfont\color{white}\MakeUppercase{Télécharge OnBuch}}\par\vspace{4pt}
        {\raggedright\popsemi\fontsize{8.4}{11}\selectfont\color{white}Gratuit \textbullet\ Google Play\\Tuteur IA, annales, concours, résultats\par}
      \end{minipage}
    \end{tcolorbox}
  \end{minipage}\hfill
  \begin{minipage}[t]{0.487\linewidth}
    \begin{tcolorbox}[enhanced,colback=poppurple,colframe=popink,boxrule=2.2pt,arc=10pt,
        drop shadow=popink,left=11pt,right=11pt,top=10pt,bottom=10pt,height=3.1cm,valign=center]
      \tikz[baseline=-0.7ex]\node[circle,fill=white,draw=popink,line width=1.3pt,minimum size=1.6cm,inner sep=0]{\popdisplay\fontsize{24}{24}\selectfont\color{poppurple}+};%
      \hspace{8pt}\begin{minipage}[c]{0.6\linewidth}
        {\popdisplay\fontsize{11}{12.5}\selectfont\color{white}\MakeUppercase{Passe à OnBuch+}}\par\vspace{4pt}
        {\raggedright\popsemi\fontsize{8.4}{11}\selectfont\color{white}Tous les cours signés, Tuteur IA illimité, fiches PDF — onglet OnBuch+ de l'app\par}
      \end{minipage}
    \end{tcolorbox}
  \end{minipage}
  \vspace{8pt}
  \popcontenu
  \vfill
  \signatureonbuch
  \restoregeometry
  \newpage
}

% Rangée « Ce cours contient » (page de garde)
\newcommand{\poptile}[3]{% #1 couleur #2 gros texte #3 légende
  \begin{tcolorbox}[enhanced,colback=#1,colframe=popink,boxrule=2pt,arc=9pt,shadow={2.5pt}{-2.5pt}{0pt}{popink},
      left=6pt,right=6pt,top=9pt,bottom=9pt,halign=center,valign=center,height=1.75cm,width=\linewidth,nobeforeafter]
    {\popdisplay\fontsize{12.5}{13}\selectfont\color{white}\MakeUppercase{#2}}\par\vspace{3pt}
    {\centering\popsemi\fontsize{7.8}{9.5}\selectfont\color{white}#3\par}
  \end{tcolorbox}}
\newcommand{\popcontenu}{%
  \par\noindent{\popxboldls\fontsize{7.5}{7.5}\selectfont\color{popmuted}CE COURS SIGNÉ CONTIENT}\par\vspace{7pt}
  \noindent\begin{minipage}[t]{0.235\linewidth}\poptile{popblue}{Cours}{complet, illustré, pas à pas}\end{minipage}\hfill
  \begin{minipage}[t]{0.235\linewidth}\poptile{poporange}{Méthodes}{et exercices résolus}\end{minipage}\hfill
  \begin{minipage}[t]{0.235\linewidth}\poptile{poppink}{Exercices}{type examen + corrigés}\end{minipage}\hfill
  \begin{minipage}[t]{0.235\linewidth}\poptile{popgreen}{Fiche bilan}{l'essentiel à retenir}\end{minipage}\par}

% Sommaire
\newtcolorbox{sommaire}{enhanced,colback=white,colframe=popink,boxrule=2.2pt,arc=10pt,
  drop shadow=popink,left=18pt,right=18pt,top=13pt,bottom=13pt,before skip=6pt,after skip=20pt,
  before upper={\poppill{Sommaire}{poppurple}{white}\par\vspace{9pt}\popsemi\fontsize{10.6}{14.5}\selectfont\color{popink}}}

% Signature de fin de cours — poussée tout en bas de la dernière page
\newcommand{\findecours}{%
  \par\vfill\nopagebreak
  \begin{center}\popsquiggle{poporange}\end{center}\vspace{4pt}
  \signatureonbuch
}
\AtBeginDocument{\def\llbracket{\mathopen{[\mkern-3.5mu[}}\def\rrbracket{\mathclose{]\mkern-3.5mu]}}} % intervalles d'entiers (glyphe ⟦ absent de la police)
% Glyphes absents de Fira Math : redéfinis à partir de glyphes disponibles
\AtBeginDocument{%
  \def\setminus{\mathbin{\backslash}}\let\smallsetminus\setminus
  \def\vdots{\mathord{\vbox{\baselineskip3.2pt\lineskiplimit0pt\kern4pt\hbox{.}\hbox{.}\hbox{.}}}}%
  \def\ddots{\mathinner{\mkern1mu\raise6.5pt\hbox{.}\mkern2mu\raise3.6pt\hbox{.}\mkern2mu\raise0.7pt\hbox{.}\mkern1mu}}%
  \def\triangle{\mathord{\scalebox{0.9}{$\symup{\Delta}$}}}%
  \def\nabla{\mathord{\raisebox{\depth}{\scalebox{1}[-1]{$\symup{\Delta}$}}}}%
  \def\checkmark{\popcheck{popdark}}%
  \def\star{\mathbin{\ast}}\def\bigstar{\mathord{\ast}}%
  \def\diamond{\mathbin{\scalebox{0.6}{$\Diamond$}}}%
  \def\bigcup{\mathop{\mathchoice{\raisebox{-0.3em}{\scalebox{1.7}{$\cup$}}}{\scalebox{1.25}{$\cup$}}{\cup}{\cup}}\displaylimits}%
  \def\bigcap{\mathop{\mathchoice{\raisebox{-0.3em}{\scalebox{1.7}{$\cap$}}}{\scalebox{1.25}{$\cap$}}{\cap}{\cap}}\displaylimits}%
  \def\lesseqqgtr{\mathrel{\lessgtr}}\def\bigoplus{\mathop{\oplus}\displaylimits}%
}
\providecommand{\ssi}{si et seulement si} % abréviation fréquente
\AtBeginDocument{\providecommand{\wideparen}{\overparen}} % arc de cercle

\begin{document}

\pagegarde{La diversité des ensembles biogéographiques}{Géographie}{Tle Tech}{Géographie – classe de Terminale technique}{N01}{2 séances (fiche de progression harmonisée)}{Cette leçon présente la diversité physique du Cameroun (reliefs, climats, végétations, sols) et les grands groupes humains qui peuplent le territoire. Elle montre comment cette double diversité, loin d'être un handicap, constitue la richesse et le fondement de l'unité nationale camerounaise.
\vspace{4pt}
\begin{multicols}{2}\small
\begin{itemize}
\item À la fin de cette leçon, je sais décrire les grands types de relief du Cameroun à partir d'une carte et de photographies
\item À la fin de cette leçon, je sais distinguer les grands domaines climatiques et végétaux du Cameroun et les localiser sur une carte
\item À la fin de cette leçon, je sais expliquer les relations entre climat, végétation, sols et activités humaines dans chaque ensemble biogéographique
\item À la fin de cette leçon, je sais identifier et localiser les grands groupes humains du Cameroun (soudaniens, peuls, bantous, semi-bantous, pygmées)
\item À la fin de cette leçon, je sais construire un croquis géographique et un diagramme à partir de données chiffrées
\item À la fin de cette leçon, je sais appliquer la notion d'unité dans la diversité à des situations concrètes de la vie courante (migrations, échanges, brassage culturel)
\end{itemize}\end{multicols}}

\begin{sommaire}
\begin{multicols}{2}
\begin{enumerate}
\item Introduction : le Cameroun, « Afrique en miniature »
\item La diversité des reliefs et des climats
\item La diversité des végétations et des sols
\item Les grands ensembles biogéographiques du Cameroun
\item Les grands groupes humains du Cameroun
\item L'unité dans la diversité : atouts et défis
\item Activité d'intégration
\item Méthodes et exercices résolus
\item Exercices
\item Corrigés des exercices
\item Fiche bilan
\end{enumerate}
\end{multicols}
\end{sommaire}

\begin{objectifs}
\begin{itemize}
\item À la fin de cette leçon, je sais décrire les grands types de relief du Cameroun à partir d'une carte et de photographies
\item À la fin de cette leçon, je sais distinguer les grands domaines climatiques et végétaux du Cameroun et les localiser sur une carte
\item À la fin de cette leçon, je sais expliquer les relations entre climat, végétation, sols et activités humaines dans chaque ensemble biogéographique
\item À la fin de cette leçon, je sais identifier et localiser les grands groupes humains du Cameroun (soudaniens, peuls, bantous, semi-bantous, pygmées)
\item À la fin de cette leçon, je sais construire un croquis géographique et un diagramme à partir de données chiffrées
\item À la fin de cette leçon, je sais appliquer la notion d'unité dans la diversité à des situations concrètes de la vie courante (migrations, échanges, brassage culturel)
\item À la fin de cette leçon, je sais rédiger et justifier une démarche, une réponse ou une conclusion en géographie
\item À la fin de cette leçon, je sais valoriser la diversité biogéographique et humaine comme atout de développement national
\end{itemize}
\end{objectifs}

\begin{prerequis}
\begin{itemize}
\item Les grandes lignes du relief et du climat du Cameroun vues en classe de 4ème
\item La localisation des dix régions du Cameroun et de leurs chefs-lieux
\item La lecture d'une carte géographique (légende, échelle, orientation)
\item Les notions de population, d'ethnie et de densité vues en Première
\item La construction d'un croquis géographique simple
\end{itemize}
\end{prerequis}

\newpage

\coursec{Introduction : le Cameroun, « Afrique en miniature »}

\courssub{Une situation d'accroche : le voyage de Maroua à Kribi}

Imaginez un voyageur qui quitte \textbf{Maroua}, au cœur du Grand Nord, par une matinée sèche et brûlante. L'horizon est immense, le ciel d'un bleu lavé, le sol sableux parsemé d'acacias épineux et de baobabs majestueux. Il croise des troupeaux de zébus conduits par des éleveurs \textbf{Peuls} aux vêtements amples, des femmes \textbf{Kotoko} revenant de la pêche sur le Logone, des champs de mil et de sorgho qui attendent les premières pluies. L'air est chaud, le vent porte la poussière du Sahel.

Quelques jours plus tard, le même voyageur arrive à \textbf{Kribi}, sur la côte atlantique, au sud du pays. Tout a changé : l'air est lourd, saturé d'humidité, la température est douce mais étouffante. Devant lui s'étend la \textbf{forêt dense humide}, un mur vert impénétrable où s'élancent des arbres géants -- iroko, sapelli, ayous -- hauts de 40 à 50 mètres. Le sol est recouvert d'une épaisse litière, les lianes s'entrelacent, la lumière filtre à peine. Sur la plage, des pirogues des pêcheurs \textbf{Batanga} et \textbf{Mabi} rentrent chargées de poissons ; dans l'arrière-pays, des planteurs \textbf{Bassa} et \textbf{Bulu} entretiennent des champs de cacao, de bananiers et d'hévéas. Les chants, les langues, les odeurs, les gestes quotidiens n'ont plus rien de commun avec ceux du Nord.

\medskip
\noindent
\textbf{Comment un seul pays peut-il offrir une telle diversité de paysages, de climats, de végétations et de peuples en l'espace de quelques centaines de kilomètres ?} Pourquoi l'expression « \cle{Afrique en miniature} » s'applique-t-elle si parfaitement au Cameroun ? Cette diversité est-elle une fragilité ou une richesse ? C'est à ces questions que cette leçon répondra en explorant les \cle{grands ensembles biogéographiques} du Cameroun et les \cle{groupes humains} qui les occupent, pour comprendre les fondements de l'unité nationale dans la diversité.

\courssub{La notion d'ensemble biogéographique : définition et portée}

Avant d'aller plus loin, il faut préciser l'outil conceptuel qui va structurer notre étude. Observer un paysage, ce n'est pas seulement voir des couleurs et des formes ; c'est lire une \textbf{combinaison originale} de facteurs naturels et humains.

\begin{definition}[Ensemble biogéographique]
Un \cle{ensemble biogéographique} (ou région naturelle) est une \textbf{portion de l'espace terrestre} caractérisée par une \textbf{combinaison originale et relativement homogène} de quatre composantes majeures :
\begin{enumerate}
    \item le \cle{climat} (températures, pluviométrie, saisonnalité) ;
    \item la \cle{végétation} (formation végétale climax, strates, espèces dominantes) ;
    \item les \cle{sols} (type pédologique, fertilité, épaisseur, drainage) ;
    \item l'\cle{occupation humaine} (peuplement, activités économiques, types d'habitat, modes de vie).
\end{enumerate}
Ces composantes sont en \textbf{interaction permanente} : le climat commande la végétation et la pédogenèse ; la végétation protège le sol et influence le microclimat ; l'homme adapte ses activités au milieu et le transforme en retour.
\end{definition}

\begin{methode}[Analyser une photographie de paysage géographique]
Pour « lire » un paysage comme un géographe, suivez cette démarche en quatre temps :
\begin{enumerate}
    \item \textbf{Décrire} : repérez les plans (premier plan, plan moyen, arrière-plan) ; notez les formes du relief, la couverture végétale, la présence d'eau, les indices d'activités humaines (champs, cases, routes, lignes électriques).
    \item \textbf{Identifier} : nommez les éléments naturels (type d'arbres, aspect du sol, cours d'eau) et humains (type d'habitat, cultures, élevage, densité de peuplement apparente).
    \item \textbf{Localiser} : à l'aide des indices climatiques (aridité / humidité), végétaux (espèces indicatrices) et humains, proposez une localisation régionale (ex. : « Grand Nord soudano-sahélien », « Plateau de l'Adamaoua », « Littoral équatorial »).
    \item \textbf{Expliquer} : replacez le paysage dans son ensemble biogéographique : quel climat ? quel type de sol ? quelles activités dominantes ? quelles contraintes et quels atouts pour les populations ?
\end{enumerate}
\end{methode}

\begin{exemplebox}[Application : deux photographies contrastées]
\textbf{Photo A (Maroua, mai) :} Premier plan : sol nu, latéritique, quelques touffes d'herbe sèche. Plan moyen : acacias \emph{Acacia seyal}, baobabs (\emph{Adansonia digitata}), cases en banco à toit conique. Arrière-plan : plaine infinie, horizon flou (harmattan). \textbf{Analyse :} Végétation xérophile (adaptée à la sécheresse), sols ferrugineux tropicaux lessivés, climat soudano-sahélien (saison sèche 7--8 mois, pluviométrie \SIrange{700}{1000}{mm}). Activités : élevage extensif, culture pluviale de mil/sorgho. \textbf{Ensemble :} \textbf{Domaine soudanien / sahélien (Grand Nord)}.

\textbf{Photo B (Ebolowa, octobre) :} Premier plan : litière épaisse, champignons, racines tabulaires. Plan moyen : fûts droits d'arbres géants (iroko, ayous), lianes, strate arbustive dense. Arrière-plan : canopée continue, pas d'horizon visible. \textbf{Analyse :} Forêt dense humide \textbf{semi-caducifoliée / de transition}, sols ferrallitiques (lessivés, acides, pauvres en bases), climat équatorial de transition (deux saisons des pluies, pluviométrie \SIrange{1500}{2000}{mm}). Activités : cueillette, chasse, agriculture itinérante sur brûlis, plantations pérennes (cacao). \textbf{Ensemble :} \textbf{Domaine équatorial (Sud-Cameroun)}.
\end{exemplebox}

\courssub{Le Cameroun, carrefour géographique : position de transition}

La diversité exceptionnelle du Cameroun tient d'abord à sa \textbf{position géographique unique} sur le continent africain. C'est un pays \textbf{charnière}, un \textbf{carrefour} où se rencontrent et s'interpénètrent de grands ensembles continentaux.

\begin{popfigure}
\begin{tikzpicture}[scale=1.2]
% Coastline and countries - simplified West/Central Africa
\fill[popblue!15] (-3,-4) rectangle (6,3); % Ocean
\fill[popgreen!20] (-3,-1) .. controls (-1,0) and (0,1) .. (2,1.5) .. controls (3,1) and (4,0) .. (5,-1) -- (5,-4) -- (-3,-4) -- cycle; % Land mass approx
% Cameroon shape (simplified polygon)
\fill[poporange!40] (0.5,-0.5) -- (2.5,0.2) -- (3.5,1.2) -- (3.2,2.5) -- (1.5,2.8) -- (0.2,1.8) -- (-0.2,0.5) -- cycle;
\draw[thick, poporange] (0.5,-0.5) -- (2.5,0.2) -- (3.5,1.2) -- (3.2,2.5) -- (1.5,2.8) -- (0.2,1.8) -- (-0.2,0.5) -- cycle;
% Neighboring countries labels
\node[popdark, font=\small\bfseries] at (-2, 1.5) {Nigéria};
\node[popdark, font=\small\bfseries] at (4.5, 1.5) {Tchad};
\node[popdark, font=\small\bfseries] at (4.5, -1.5) {RCA};
\node[popdark, font=\small\bfseries] at (1.5, -2.5) {Congo / Gabon / Guinée Eq.};
\node[popdark, font=\small\bfseries, rotate=90] at (-2.5, -1) {Océan Atlantique};
\node[popdark, font=\small\bfseries, text width=2.5cm, align=center] at (1.8, 1.2) {\textbf{CAMEROUN}};
% Capital
\fill[popdark] (1.8, 1.2) circle (1.5pt);
\node[above, font=\tiny] at (1.8, 1.2) {Yaoundé};
% Maroua
\fill[popdark] (2.8, 0.0) circle (1pt);
\node[below right, font=\tiny] at (2.8, 0.0) {Maroua};
% Kribi
\fill[popdark] (0.8, 0.5) circle (1pt);
\node[below left, font=\tiny] at (0.8, 0.5) {Kribi};
% North arrow
\draw[->, thick, popdark] (4.5, 2.5) -- (4.5, 2.9);
\node[right, font=\tiny] at (4.5, 2.9) {N};
% Scale bar (approx)
\draw[thick, popdark] (4.5, -3.5) -- (5.5, -3.5);
\node[below, font=\tiny] at (5, -3.5) {~\SI{500}{km}};
\end{tikzpicture}
\legende{Carte de situation du Cameroun en Afrique (contours simplifiés). Le pays fait l'articulation entre l'Afrique de l'Ouest (au nord-ouest) et l'Afrique centrale (au sud-est), entre le domaine sahélien (au nord) et le domaine équatorial / golfe de Guinée (au sud).}
\end{popfigure}

\begin{itemize}
    \item \textbf{Entre Afrique de l'Ouest et Afrique centrale :} Le Cameroun prolonge vers l'est les reliefs de la ligne volcanique du Cameroun (monts Mandara, Alantika, Adamaoua, Manengouba, Cameroun) qui se prolongent dans l'océan (Bioko, Annobón, São Tomé, Principe). Il marque la transition entre les bassins versants du Niger (Bénoué) et du Congo (Sangha, Oubangui, Sanaga).
    \item \textbf{Entre Sahel et Golfe de Guinée :} Le 10\textsuperscript{e} parallèle nord (latitude de Kousseri) reçoit moins de \SI{800}{mm} de pluie/an (influence sahélienne) ; le 2\textsuperscript{e} parallèle nord (Kribi, Campo) en reçoit plus de \SI{3000}{mm} (influence équatoriale maritime). C'est un \cle{gradient climatique nord-sud} de près de \SI{1000}{km} qui condense les grands types climatiques africains.
    \item \textbf{Une « charnière » géologique :} Le socle précambrien (au nord et au sud), les bassins sédimentaires (bassin du lac Tchad, bassin côtier), les chaînes volcaniques récentes (ligne du Cameroun) et les plateaux anciens (Adamaoua) offrent une mosaïque de substrats géologiques.
\end{itemize}

\begin{savaistu}[Le Cameroun, champion de la diversité ethnolinguistique]
Le Cameroun compte \textbf{plus de 250 groupes ethniques} (certains recensements en dénombrent près de 300) et \textbf{autant de langues nationales}, réparties en trois grandes familles linguistiques : \textbf{afro-asiatiques} (tchadiques : peul, mofu, kotoko... au Nord), \textbf{nilo-sahariennes} (sara, touri... à l'Est) et \textbf{nigéro-congolaises} (bantoues : bassa, bamiléké, beti, douala... au Sud/Centre/Ouest ; bantoïdes : tikar, bamoun... à l'Ouest). Cette diversité fait du Cameroun l'un des pays les plus diversifiés d'Afrique, un véritable « musée des peuples » où cohabitent éleveurs nomades, pêcheurs lacustres, agriculteurs de montagne, planteurs de forêt et citadins cosmopolites.
\end{savaistu}

\courssub{Problématique et annonce du plan}

La diversité physique et humaine que nous venons d'effleurer n'est pas un simple inventaire. Elle pose une question fondamentale pour la géographie et pour la construction nationale :

\begin{aretenir}[Problématique de la leçon]
\textbf{Comment la diversité physique et humaine du Cameroun s'organise-t-elle en \cle{ensembles biogéographiques} cohérents, et en quoi cette organisation fonde-t-elle l'unité nationale dans la diversité ?}
\end{aretenir}

Pour y répondre, nous suivrons le plan suivant :

\begin{enumerate}
    \item \textbf{La diversité physique (I) :} nous identifierons les grands ensembles de reliefs, de climats, de végétations et de sols qui découpent le territoire en \textbf{cinq grands domaines biogéographiques} (Domaine sahélien, Domaine soudanien, Domaine de l'Adamaoua, Domaine équatorial, Domaine littoral et volcanique).
    \item \textbf{Les grands groupes humains (II) :} nous étudierons la répartition des populations, les grandes familles linguistiques, les modes de vie traditionnels (élevage, agriculture, pêche, chasse-cueillette) et leur adaptation aux milieux, ainsi que la dynamique urbaine actuelle.
    \item \textbf{L'unité dans la diversité (III) :} nous analyserons comment les complémentarités économiques (échanges Nord/Sud), les institutions nationales (bilinguisme, décentralisation), les symboles partagés et l'histoire commune transforment la diversité en ciment de la nation camerounaise.
\end{enumerate}

\begin{attention}[Piège à éviter : Diversité $\neq$ Division]
Il est fréquent de confondre \textbf{diversité} (constat objectif : variété des milieux, des cultures, des langues) et \textbf{division} (jugement subjectif : opposition, conflit, éclatement). La diversité est un \textbf{fait géographique et anthropologique} ; la division est un \textbf{risque politique}. L'histoire du Cameroun montre que la diversité, bien gérée (fédéralisme, puis État unitaire décentralisé, bilinguisme officiel), est un \textbf{atout} (résilience, complémentarités, rayonnement culturel). Ne confondez pas « pays divers » et « pays divisé ».
\end{attention}

\begin{exoresolu}[Exercice : Caractériser la position de transition du Cameroun]
\textbf{Énoncé :} À l'aide de la carte de situation et de vos connaissances, montrez que le Cameroun occupe une position de transition entre trois grands ensembles africains. Citez pour chacun un indice physique et un indice humain.

\textbf{Solution détaillée :}
\begin{enumerate}
    \item \textbf{Transition Sahel / Soudan / Guinée (gradient Nord-Sud) :}
        \begin{itemize}
            \item \textit{Indice physique :} Gradient pluviométrique : $<$ \SI{800}{mm/an} à Kousseri (climat sahélien) $\rightarrow$ \SIrange{1500}{2000}{mm/an} à Yaoundé (climat équatorial de transition) $\rightarrow$ $>$ \SI{3000}{mm/an} à Campo / Debundscha (climat équatorial montagnard/maritime).
            \item \textit{Indice humain :} Élevage nomade/transhumant (Peuls, Choa) au Nord $\leftrightarrow$ Agriculture sédentaire intensive (Bamiléké, Bassa) au Centre/Sud $\leftrightarrow$ Pêche lagunaire et plantation (Bakweri, Douala) sur la côte.
        \end{itemize}
    \item \textbf{Transition Afrique de l'Ouest / Afrique Centrale (gradient Ouest-Est) :}
        \begin{itemize}
            \item \textit{Indice physique :} Ligne volcanique du Cameroun (chaîne Yadé - Mont Cameroun) qui sépare le bassin du Niger (Bénoué) du bassin du Congo (Sanaga, Nyong, Oubangui).
            \item \textit{Indice humain :} Populations soudanaises (langues tchadiques : Mofu, Mandara) à l'Ouest/Nord $\leftrightarrow$ Populations bantoues (Beti, Fang, Bassa) au Sud/Est ; présence de populations « charnières » (Tikar, Bamoun) sur l'Adamaoua.
        \end{itemize}
    \item \textbf{Transition Continent / Océan (facade atlantique) :}
        \begin{itemize}
            \item \textit{Indice physique :} Étroite plaine côtière (0--\SI{200}{m}) bordant le golfe de Guinée, mangroves à l'embouchure du Wouri et de la Sanaga.
            \item \textit{Indice humain :} Métropoles portuaires (Douala, Kribi, Limbé) tournées vers l'exportation (bois, cacao, pétrole, coton) et l'importation ; populations côtières (Duala, Bakweri, Batanga) historiquement intermédiaires dans le commerce atlantique.
        \end{itemize}
\end{enumerate}
\end{exoresolu}

\courssub{Schémas comparatifs : profils paysagers Nord / Sud}

Pour ancrer visuellement le contraste évoqué en introduction, observons deux profils topographiques et paysagers schématiques, l'un dans le Grand Nord (axe Maroua -- Lac Tchad), l'autre dans le Sud forestier (axe Ebolowa -- Kribi).

\begin{popfigure}
\begin{tikzpicture}[scale=0.9]
% Style definitions
\tikzset{
    terrain/.style={thick, draw=popdark},
    veg/.style={fill opacity=0.6},
    label/.style={font=\small, text=popdark},
    city/.style={circle, fill=popdark, inner sep=1.5pt},
    arrow/.style={->, thick, popblue}
}

% PROFIL 1 : GRAND NORD (Maroua -> Lac Tchad)
\begin{scope}[xshift=0cm, yshift=0cm]
% Relief baseline
\draw[terrain] (0,0.5) -- (0.5,0.6) -- (1.2,0.55) -- (2,0.7) -- (2.5,0.65) -- (3.2,0.8) -- (4,0.75) -- (4.5,0.7) -- (5,0.8) -- (5.5,0.75) -- (6,0.7);
% Vegetation / Land cover zones
\fill[poporange!30, veg] (0,0) rectangle (6,0.5); % Soil base
% Savane arbustive (Maroua area)
\fill[popgreen!40!poporange, veg] (0,0.5) -- (0,0.8) -- (0.5,0.9) -- (1.2,0.85) -- (1.2,0.55) -- cycle;
\foreach \x in {0.2, 0.7, 1.1} {
    \draw[popgreen!60!popdark] (\x,0.55) -- (\x-0.05,0.9) -- (\x+0.05,0.9) -- cycle; % Acacia
    \draw[popgreen!60!popdark] (\x,0.55) -- (\x,0.85);
}
% Baobab
\draw[popdark, thick] (0.3,0.55) -- (0.3,1.3);
\draw[popgreen!40!popdark] (0.3,1.3) circle (0.25);
% Plaine de la Bénoué / Logone
\fill[popblue!20, veg] (2,0.7) -- (2,0.9) -- (3.5,0.95) -- (3.5,0.8) -- cycle;
\draw[popblue, thick] (2.2,0.85) -- (3.3,0.88); % River
\node[label, popblue] at (2.8, 1.05) {Logone / Bénoué};
% Lac Tchad (loin)
\fill[popblue!40, veg] (5,0.7) -- (5,0.9) -- (6,0.85) -- (6,0.7) -- cycle;
\node[label, popblue] at (5.5, 1.0) {Lac Tchad};
% Human elements
\node[city] (maroua) at (0.8, 0.55) {};
\node[label, above] at (maroua) {Maroua};
\draw[popdark] (1.5,0.55) -- (1.6,0.7) -- (1.5,0.85); % Hut
\draw[popdark] (1.55,0.55) -- (1.65,0.7) -- (1.55,0.85);
\node[label] at (1.5, 0.4) {Cases en banco};
% Zebus
\foreach \x in {2.5, 2.8, 3.1} {
    \draw[popdark] (\x,0.65) -- (\x+0.15,0.65) -- (\x+0.1,0.8) -- (\x-0.05,0.8) -- cycle;
    \draw[popdark] (\x+0.05,0.65) -- (\x+0.05,0.5);
}
\node[label] at (2.8, 0.4) {Troupeaux (zébus)};
% Climate arrow
\draw[arrow] (3, -0.5) -- (3, 0);
\node[label, below, popblue] at (3, -0.5) {Saison sèche : 7--8 mois};
\node[label, below] at (3, -0.7) {Pluviométrie : \SIrange{700}{1000}{mm}};
% Title
\node[label, font=\bfseries\large, align=center] at (3, 1.4) {Profil Nord : Maroua $\rightarrow$ Lac Tchad\\(Domaine sahélo-soudanien)};
\end{scope}

% PROFIL 2 : SUD FORESTIER (Ebolowa -> Kribi)
\begin{scope}[xshift=8cm, yshift=0cm]
% Relief baseline (hills)
\draw[terrain] (0,0.5) -- (0.5,0.8) -- (1.2,0.6) -- (2,1.0) -- (2.5,0.7) -- (3.2,1.2) -- (3.8,0.8) -- (4.5,1.1) -- (5,0.9) -- (5.5,0.6) -- (6,0.5);
% Soil
\fill[poporange!40!popdark, veg] (0,0) rectangle (6,0.5); % Ferrallitic soil (reddish)
% Forest canopy - dense
\foreach \x in {0.2, 0.5, 0.9, 1.3, 1.7, 2.1, 2.6, 3.0, 3.4, 3.9, 4.3, 4.7, 5.2, 5.7} {
    \pgfmathsetmacro{\h}{0.8 + rnd*0.6}
    \pgfmathsetmacro{\r}{0.2 + rnd*0.15}
    \draw[popgreen!60!popdark, thick] (\x,0.5) -- (\x,0.5+\h); % Trunk
    \fill[popgreen!50!popdark, veg] (\x,0.5+\h) circle (\r); % Crown
    % Lianas
    \draw[popgreen!40!popdark, thin] (\x-0.1,0.5) .. controls (\x-0.2,0.5+\h*0.5) .. (\x-0.1,0.5+\h);
}
% Understory
\fill[popgreen!30!popdark, veg] (0,0.5) -- (0,0.8) -- (6,0.8) -- (6,0.5) -- cycle;
% River / Coast
\fill[popblue!30, veg] (5.5,0.5) -- (5.5,0.7) -- (6,0.6) -- (6,0.5) -- cycle;
\draw[popblue, thick] (5.7,0.55) -- (5.9,0.55);
\node[label, popblue] at (5.7, 0.85) {Océan Atlantique};
% Human elements
\node[city] (ebolowa) at (1.0, 0.5) {};
\node[label, above] at (ebolowa) {Ebolowa};
\draw[popdark] (2.2,0.5) rectangle (2.5,0.8); % House
\draw[popdark] (2.35,0.8) -- (2.35,1.0); % Roof
\node[label] at (2.35, 0.35) {Cases / Maison en dur};
% Cocoa plantation (near coast)
\foreach \x in {4.0, 4.3, 4.6, 4.9} {
    \draw[poporange!60!popdark] (\x,0.5) -- (\x,1.0);
    \fill[popgreen!40!poporange] (\x,1.0) circle (0.1);
}
\node[label] at (4.5, 0.3) {Plantation cacaoyère};
% Climate arrow
\draw[arrow] (3, -0.5) -- (3, 0);
\node[label, below, popblue] at (3, -0.5) {Saison sèche : 2--3 mois (courte)};
\node[label, below] at (3, -0.7) {Pluviométrie : \SIrange{1500}{2000}{mm}};
% Title
\node[label, font=\bfseries\large, align=center] at (3, 1.6) {Profil Sud : Ebolowa $\rightarrow$ Kribi\\(Domaine équatorial de transition)};
\end{scope}

% Comparison legend in middle
\node[label, font=\small\bfseries, align=center, text width=3.5cm] at (7, 1.8) {
\textbf{Contraste majeur} \\
Relief : Plaine vs Plateau vallonné \\
Végétation : Savane ouverte vs Forêt fermée \\
Sol : Ferrugineux (jaune/brun) vs Ferrallitique (rouge) \\
Climat : 1 saison des pluies vs 2 saisons des pluies \\
Activités : Élevage / Mil vs Forêt / Cacao
};
\end{tikzpicture}
\legende{Schéma comparatif de deux profils paysagers types : (gauche) Grand Nord sahélo-soudanien autour de Maroua -- plaine sableuse, savane arbustive à acacias et baobabs, rivière saisonnière, élevage bovin ; (droite) Sud forestier équatorial Ebolowa -- Kribi -- plateau vallonné, forêt dense humide semi-caducifoliée à canopée continue, sols ferrallitiques rouges, plantations pérennes (cacao) vers la côte. Échelle horizontale non respectée (distance ~ \SI{800}{km}).}
\end{popfigure}

\medskip
\noindent
Ces deux profils résument le \cle{gradient bioclimatique majeur} qui structure le Cameroun. Mais entre ces deux extrêmes, il existe des mondes intermédiaires -- les hautes herbes de l'Adamaoua, les forêts de galerie le long des rivières, les savanes d'altitude des monts Mandara, les mangroves de l'estuaire du Wouri -- que nous allons maintenant explorer systématiquement. C'est l'objet de la première grande partie de cette leçon : \textbf{la diversité physique du Cameroun}.

\coursec{La diversité des reliefs et des climats}

\courssub{Transition : le décor physique d'une « \cle{Afrique en miniature} »}

Après avoir posé le cadre général du Cameroun comme synthèse du continent africain, nous devons maintenant examiner le socle physique qui conditionne tout le reste : le relief et le climat. C'est leur combinaison qui dessine les grands ensembles biogéographiques (végétation, sols, hydrographie) et oriente l'installation des populations et des activités économiques. Comprendre cette diversité, c'est saisir pourquoi le Cameroun concentre, sur \SI{475 442}{km²}, une variété de paysages et de milieux qu'on ne retrouve ailleurs qu'à l'échelle de l'Afrique entière.

\courssub{Les grandes unités de relief : une architecture en gradins}

Le relief camerounais s'organise en quatre grands ensembles morphostructuraux, disposés grossièrement du Sud vers le Nord, mais avec une charnière majeure : la \cle{dorsale camerounaise}.

\begin{definition}[Plaines et bas plateaux du Sud (0--650 m)]
Vaste ensemble sédimentaire et précambrien couvrant tout le Sud-Cameroun (régions du Centre, Sud, Est, Littoral, Sud-Ouest hors montagne). L'altitude y est faible, les interfluves sont arrondis (pénéplaine), entaillés par des vallées larges. Les sols sont souvent profonds, latéritiques, lessivés par les fortes pluies équatoriales.
\end{definition}

\begin{definition}[Hautes terres de l'Ouest : massifs bamiléké, monts Bamboutos, mont Oku]
Zone de volcanisme récent (tertiaire--quaternaire) sur socle ancien. Altitudes comprises entre \SI{1 000}{m} et \SI{3 000}{m} (mont Oku : \SI{3 011}{m}). Paysages de cônes volcaniques, de caldeiras (lac Baleng, lac Oku), de plateaux basaltiques fertiles. Climat tempéré d'altitude, forte densité de population, agriculture intensive (café arabica, maraîchage).
\end{definition}

\begin{definition}[Dorsale camerounaise et mont Cameroun (\SI{4 095}{m})]
Alignement volcanique SW--NE (ligne du Cameroun) reliant l'île d'Annobón (Guinée équatoriale) au lac Tchad. Le mont Cameroun (Mont Fako) en est l'extrémité sud-ouest active, stratovolcan culminant à \textbf{\SI{4 095}{m}}, point culminant de l'Afrique de l'Ouest et centrale. Pentes raides, coulées de lave récentes, végétation étagée.
\end{definition}

\begin{definition}[Plateaux de l'Adamaoua (\SI{1 000}{m}--\SI{1 500}{m}) : « \cle{château d'eau} »]
Vaste plateau gréseux et volcanique, charnière climatique et hydrographique. Source des grands fleuves : Sanaga (vers l'Atlantique), Bénoué (vers le Niger), Logone/Chari (vers le lac Tchad). Climat tropical d'altitude, végétation de savane arborée, pâturages extensifs.
\end{definition}

\begin{definition}[Plaines du Nord : Logone, lac Tchad, Bénoué]
Plaines alluviales et sédimentaires basses (\SI{200}{m}--\SI{400}{m}). Plaine du Logone (inondable, riziculture), cuvette du lac Tchad (régression récente), plaine du Bénoué (vallée large, sols fertiles, coton). Climat soudanien à sahélien.
\end{definition}

\begin{popfigure}
\begin{tikzpicture}[scale=0.55]
% Contour très simplifié du Cameroun (polygone approximatif)
\draw[popdark, thick] (0,0) -- (2,0) -- (3,1) -- (4,0.5) -- (5,1.5) -- (6,1) -- (7,2) -- (7.5,3) -- (7,4) -- (6.5,5) -- (6,6) -- (5.5,7) -- (4.5,8) -- (3.5,8.5) -- (2.5,8) -- (1.5,7.5) -- (0.5,6.5) -- (0,5) -- cycle;
% Zones de relief : couleurs semi-transparentes
% 1. Plaines Sud (bas)
\fill[popblue!20] (0.5,0.5) -- (2,0.5) -- (3,1.5) -- (4,1) -- (4.5,2) -- (4,3) -- (3.5,4) -- (2.5,4.5) -- (1.5,4) -- (0.8,3) -- cycle;
\node[popblue, font=\bfseries\small, align=center] at (2,2.5){Plaines \& bas plateaux\\ du Sud (0--650 m)};
% 2. Hautes terres Ouest
\fill[poppurple!25] (3.5,3) -- (5,2.5) -- (5.5,3.5) -- (5,4.5) -- (4.5,5) -- (3.8,4.5) -- cycle;
\node[poppurple, font=\bfseries\small, align=center] at (4.5,3.8){Hautes terres\\ de l'Ouest};
% Mont Oku
\fill[poppurple!50] (3.2,5.5) circle (0.35);
\node[poppurple, font=\tiny] at (3.2,5.5) {Oku (\SI{3011}{m})};
% 3. Dorsale & Mont Cameroun
\fill[poporange!30] (5.5,1.5) -- (6.5,1.2) -- (7,2) -- (6.5,2.5) -- (5.8,2.2) -- cycle;
% Mont Cameroun sommet
\fill[poporange!80] (6.2,1.6) circle (0.25);
\node[poporange, font=\bfseries\tiny, white] at (6.2,1.6) {Fako \SI{4095}{m}};
\node[poporange, font=\bfseries\small, align=center] at (6.5,0.5){Dorsale\\ camerounaise};
% 4. Adamaoua
\fill[popgreen!25] (3.5,5.5) -- (5.5,5) -- (6,6) -- (5.5,7) -- (4.5,7.5) -- (3.5,7) -- cycle;
\node[popgreen!70!black, font=\bfseries\small, align=center] at (4.8,6.3){Plateaux de\\ l'Adamaoua\\ (\SI{1000}{m}--\SI{1500}{m})\\(Château d'eau)};
% 5. Plaines Nord
\fill[popgold!30] (3.5,7) -- (5.5,6.5) -- (6,7.5) -- (5.5,8.5) -- (4,8.5) -- (3,8) -- cycle;
\node[popgold!70!black, font=\bfseries\small, align=center] at (4.5,7.8){Plaines du Nord\\(Logone, Bénoué,\\Lac Tchad)};
% Villes repères
\foreach \x/\y/\name/\col in {
    2.5/2.5/Yaoundé/popblue,
    4.2/3.8/Bafoussam/poppurple,
    6.2/1.8/Buea/poporange,
    4.8/6.3/Ngaoundéré/popgreen!70!black,
    4.5/8.0/Maroua/popgold!70!black,
    1.5/6.5/Bertoua/popblue,
    5.5/4.5/Foumban/poppurple} {
    \fill[\col] (\x,\y) circle (0.08);
    \node[\col, font=\tiny, anchor=west] at (\x+0.15,\y) {\name};
}
% Flèches hydrographiques (schématiques)
\draw[popblue, thick, ->] (4.8,6.3) -- (3.5,4.5) node[midway, above left, font=\tiny, popblue] {Sanaga};
\draw[popblue, thick, ->] (4.8,6.3) -- (5.5,8.0) node[midway, right, font=\tiny, popblue] {Bénoué};
\draw[popblue, thick, ->] (4.8,6.3) -- (3.0,7.5) node[midway, left, font=\tiny, popblue] {Logone};
% Nord
\draw[->, popdark] (7.5,4) -- (8,4) node[right, font=\small] {N};
\end{tikzpicture}
\legende{Carte schématique des grandes unités de relief du Cameroun. Les couleurs distinguent les cinq ensembles morphostructuraux. Le mont Cameroun (\SI{4095}{m}) et le mont Oku (\SI{3011}{m}) sont les points culminants. L'Adamaoua, « château d'eau », alimente trois bassins versants majeurs (Sanaga, Bénoué, Logone). Échelle indicative.}
\end{popfigure}

\begin{exemplebox}[Repérage sur carte muette]
Sur une carte muette du Cameroun, localisez et nommez : le mont Cameroun, le mont Oku, le plateau de l'Adamaoua, la plaine du Bénoué, la cuvette du lac Tchad. Tracez la dorsale camerounaise et indiquez par des flèches les trois directions d'écoulement des fleuves prenant source sur l'Adamaoua.
\end{exemplebox}

\courssub{Coupe topographique schématique : le profil Sud--Nord}

Pour visualiser l'enchaînement des seuils et des bassins, une coupe topographique le long du méridien \SI{12}{\degree} Est (environ) est parlante.

\begin{popfigure}
\begin{tikzpicture}
\begin{axis}[
    width=\linewidth,
    height=7cm,
    xlabel={Distance approximative Sud--Nord (km)},
    ylabel={Altitude (m)},
    xmin=0, xmax=800,
    ymin=0, ymax=4200,
    domain=0:800,
    samples=200,
    axis lines=middle,
    enlargelimits={abs=10},
    grid=major,
    grid style={popline},
    tick label style={font=\small},
    label style={font=\small},
    clip=false
]
% Profil simulé : Sud (0) -> Yaoundé ~200km -> Adamaoua ~500km -> Nord ~800km
% Correction : côte à 0 m
\addplot[popcurve, thick, popblue] coordinates {
    (0, 0)      % Côte (Douala/Kribi) ~ 0 m
    (50, 200)   % Plaine côtière
    (100, 650)  % Seuil du Sud
    (200, 750)  % Yaoundé
    (250, 600)  % Vallée Sanaga
    (300, 1500) % Piémont Adamaoua / Hautes terres Ouest
    (350, 2000) % Hautes terres Ouest (Bafoussam)
    (400, 1800) % Col vers Adamaoua
    (450, 1100) % Dépression de Ngaoundéré
    (500, 1400) % Plateau Adamaoua central
    (550, 1300) % Vers le Nord
    (600, 600)  % Seuil de Poli / Bénoué
    (700, 350)  % Plaine du Bénoué (Garoua)
    (800, 300)  % Extrême Nord (Maroua / Lac Tchad)
};
% Mont Cameroun (hors profil principal mais annoté)
\addplot[only marks, mark=*, mark size=3, poporange] coordinates {(50, 4095)};
\node[poporange, font=\bfseries\tiny, anchor=south west] at (50, 4095) {Mont Cameroun (\SI{4095}{m})};
% Annotations zones
\node[popblue, font=\small, align=center] at (100, 3500) {Plaines \& bas\\ plateaux du Sud};
\node[poppurple, font=\small, align=center] at (350, 3800) {Hautes terres\\ de l'Ouest};
\node[popgreen!70!black, font=\small, align=center] at (500, 3800) {Plateaux de\\ l'Adamaoua};
\node[popgold!70!black, font=\small, align=center] at (700, 3800) {Plaines du Nord};
% Ligne zéro mer
\draw[popblue, dashed] (0,0) -- (800,0) node[right, font=\tiny, popblue] {Niveau de la mer};
\end{axis}
\end{tikzpicture}
\legende{Profil topographique schématique Sud--Nord (méridien ~\SI{12}{\degree} E). On observe l'enchaînement : plaine côtière (0 m) → seuil du Sud → hautes terres de l'Ouest → plateau de l'Adamaoua (château d'eau) → seuil de Poli → plaine du Bénoué → cuvette du Lac Tchad. Le mont Cameroun (hors profil direct) domine l'ensemble.}
\end{popfigure}

\begin{aretenir}[Architecture du relief]
Le relief camerounais s'organise en \textbf{quatre marches} du Sud au Nord : (1) Plaines et bas plateaux du Sud (0--650 m), (2) Hautes terres de l'Ouest et Dorsale volcanique (point culminant : mont Cameroun, \SI{4 095}{m}), (3) Plateau de l'Adamaoua (\SI{1 000}{m}--\SI{1 500}{m}, château d'eau), (4) Plaines du Nord (Logone, Bénoué, Lac Tchad, < 500 m). Cette structure en gradins commande la circulation atmosphérique et le réseau hydrographique.
\end{aretenir}

\courssub{La diversité des climats : du équatorial au sahélien}

Le climat est le facteur dynamique qui, combiné au relief, différencie les milieux. On distingue trois grands domaines climatiques, avec des transitions et des microclimats locaux.

\courssub{Le domaine équatorial humide (Sud-Cameroun)}

Il concerne les régions du Centre, Sud, Est, Littoral (hors côte montagneuse) et Sud-Ouest (piémonts). Caractéristiques :
\begin{itemize}
    \item \textbf{Températures :} Élevées et stables toute l'année (moyennes annuelles \SIrange{24}{26}{\celsius}), faible amplitude thermique annuelle (< \SI{3}{\celsius}), amplitude diurne modérée.
    \item \textbf{Précipitations :} Abondantes, \textbf{\SIrange{1 500}{4 000}{mm/an}}. Deux maxima pluviométriques (équinoxes) et deux minima (solstices) $\rightarrow$ \textbf{4 saisons} :
    \begin{enumerate}
        \item Grande saison des pluies (mars--juin)
        \item Petite saison sèche (juillet--août)
        \item Petite saison des pluies (septembre--novembre)
        \item Grande saison sèche (décembre--février)
    \end{enumerate}
    \item \textbf{Hygrométrie :} Forte (souvent > 80 \%), nébulosité importante, ensoleillement réduit.
\end{itemize}

\begin{definition}[Climat équatorial de transition (Yaoundé)]
Yaoundé (alt. ~\SI{750}{m}) illustre le climat équatorial « atténué » par l'altitude : 4 saisons marquées, pluviométrie ~\SI{1 600}{mm}, températures douces (moyennes \SIrange{23}{24}{\celsius}). La petite saison sèche (juillet--août) est bien individualisée.
\end{definition}

\courssub{Le domaine tropical humide d'altitude (Adamaoua et Hautes terres de l'Ouest)}

L'altitude (\SIrange{1 000}{2 000}{m}) modifie le régime équatorial : baisse des températures, accentuation de l'amplitude thermique diurne, réduction du nombre de saisons (2 saisons nettes).

\begin{itemize}
    \item \textbf{Températures :} Moyennes annuelles \SIrange{18}{22}{\celsius}. Fraîcheur nocturne marquée (peut descendre < \SI{15}{\celsius} en saison sèche). Gelées possibles au sommet de l'Oku.
    \item \textbf{Précipitations :} \SIrange{1 500}{2 500}{mm/an}. \textbf{2 saisons} : saison des pluies (avril--octobre), saison sèche (novembre--mars).
    \item \textbf{Vent :} \cle{Harmattan} (vent sec NE) en saison sèche, \cle{mousson} (SW) en saison des pluies.
\end{itemize}

\begin{exemplebox}[Ngaoundéré (Adamaoua, ~\SI{1 100}{m})]
Pluviométrie ~\SI{1 500}{mm}. Saison des pluies unique (mai--octobre), maximum en août--septembre. Saison sèche longue (novembre--avril) avec harmattan. Températures moyennes : \SI{22}{\celsius} (max \SI{28}{\celsius}, min \SI{16}{\celsius}). Climat propice aux savanes arborées et à l'élevage.
\end{exemplebox}

\courssub{Le domaine soudanien puis sahélien (Grand Nord)}

Au nord du \SI{7}{\degree}--\SI{8}{\degree} parallèle, la durée de la saison sèche s'allonge, les pluies diminuent drastiquement.

\begin{itemize}
    \item \textbf{Zone soudanaire (Garoua, ~\SI{200}{m}) :} \SIrange{900}{1 200}{mm/an}. Saison des pluies courte (mai--septembre), saison sèche longue (octobre--avril). Températures très élevées mars--mai (> \SI{40}{\celsius} max).
    \item \textbf{Zone sahélienne (Maroua, Extrême-Nord, ~\SI{350}{m}) :} \textbf{\SIrange{400}{900}{mm/an}}. Saison des pluies très courte (juin--septembre), irrégulière, variabilité interannuelle forte. Saison sèche 8--9 mois. Harmattan intense. Fortes chaleurs permanentes.
\end{itemize}

\begin{attention}[Piège fréquent : uniformité du Sud]
\emph{Erreur :} « Tout le Sud du Cameroun a le même climat équatorial. »
\emph{Réalité :} L'altitude crée une rupture nette. Le climat de \textbf{Douala} (côte, \SI{4 000}{mm}, 4 saisons, chaud humide) diffère de celui de \textbf{Yaoundé} (\SI{750}{m}, \SI{1 600}{mm}, 4 saisons, plus frais) et surtout de celui de \textbf{Dschang/Bafoussam} (\SI{1 400}{m}, climat tropical d'altitude, 2 saisons, frais la nuit). Il faut toujours croiser \textbf{latitude, altitude et exposition aux vents}.
\end{attention}

\courssub{Le microclimat exceptionnel de la côte du mont Cameroun (Debundscha)}

Le mont Cameroun, par sa masse imposante (\SI{4 095}{m}) face aux alizés de sud-ouest (mousson), provoque un \cle{forçage orographique} extrême. Les vents humides sont contraints de monter, se refroidissent, se condensent et déversent des pluies diluviennes sur le flanc ouest/sud-ouest.

\begin{exemplebox}[Debundscha : record mondial]
Station située au pied ouest du mont Cameroun (alt. ~\SI{30}{m}). Pluviométrie moyenne annuelle \textbf{~\SI{10 000}{mm}} (record > \SI{14 000}{mm}). C'est l'un des lieux les plus arrosés de la planète (avec Mawsynram en Inde). Pluies quasi quotidiennes, brouillards permanents, forêt hygrophile unique. Contraste saisissant : à \SI{100}{km} au nord-est (Maroua), il tombe < \SI{800}{mm/an}.
\end{exemplebox}

\begin{savaistu}[Mont Cameroun : volcan actif]
Le mont Cameroun (Mont Fako) est un stratovolcan basaltique toujours en activité. Ses dernières éruptions effusives (coulées de lave) datent de \textbf{1999} et \textbf{2000} (coulées ayant atteint l'océan et coupé la route). Il est surveillé par l'Observatoire volcanologique de Buéa. Son activité façonne les sols (andosols très fertiles) et le microclimat local.
\end{savaistu}

\courssub{Lecture de diagrammes ombrothermiques : méthode et application}

Le \cle{diagramme ombrothermique} (ou diagramme de Bagnouls-Gaussen) est l'outil standard pour caractériser un climat. Il superpose la courbe des températures moyennes mensuelles (en \si{\celsius}) et l'histogramme des précipitations moyennes mensuelles (en \si{mm}), avec une échelle adaptée : \textbf{1 \si{\celsius} = 2 \si{mm}} (ou $P = 2T$). Cela permet de visualiser directement la \textbf{saison sèche} (barres de pluie \emph{sous} la courbe de température) et la \textbf{saison humide} (barres \emph{au-dessus}).

\begin{methode}[Lire et interpréter un diagramme ombrothermique]
\begin{enumerate}
    \item \textbf{Identifier les échelles :} Température (axe gauche, \si{\celsius}), Précipitations (axe droit, \si{mm}). Vérifier le rapport $P = 2T$ (ou $P = 3T$ pour climats très secs).
    \item \textbf{Repérer la courbe des températures :} Forme (plate = équatorial ; cloche = tropical/sahélien), valeurs min/max, amplitude annuelle.
    \item \textbf{Analyser l'histogramme des pluies :} Total annuel, nombre de pics (1 ou 2 = nombre de saisons des pluies), mois secs ($P < 2T$ ou $P < 3T$), mois humides ($P > 2T$).
    \item \textbf{Déterminer les saisons :} Compter les mois secs consécutifs. Saison sèche marquée si $P < 2T$ (ou $P < 3T$) pendant au moins 2 mois.
    \item \textbf{Conclure sur le type de climat :} Équatorial (4 saisons, $P > 2T$ quasi toute l'année), Tropical humide (2 saisons, $P > 2T$ 5--7 mois), Soudanien (2 saisons, saison sèche 5--7 mois), Sahélien (saison sèche > 8 mois, $P$ faible).
    \item \textbf{Relier au milieu :} Végétation potentielle, activités agricoles possibles, contraintes (érosion, stress hydrique).
\end{enumerate}
\end{methode}

\begin{popfigure}
\begin{tikzpicture}
\begin{axis}[
    width=\linewidth,
    height=9cm,
    axis y line*=left,
    axis x line=bottom,
    ybar,
    bar width=6pt,
    ymin=0, ymax=450,
    xmin=0.5, xmax=12.5,
    xtick={1,2,3,4,5,6,7,8,9,10,11,12},
    xticklabels={J,F,M,A,M,J,J,A,S,O,N,D},
    ylabel={Précipitations (mm)},
    ylabel style={popblue},
    tick label style={font=\small},
    label style={font=\small},
    legend style={at={(0.5,-0.18)}, anchor=north, legend columns=4, font=\small},
    grid=none,
]
% Précipitations Yaoundé (barres bleues)
\addplot[popblue!60, fill=popblue!40] coordinates {
    (1, 18) (2, 35) (3, 110) (4, 170) (5, 195) (6, 155)
    (7, 70) (8, 100) (9, 230) (10, 280) (11, 130) (12, 30)
};
\addlegendentry{Pluies Yaoundé}
% Précipitations Maroua (barres oranges)
\addplot[poporange!80, fill=poporange!50] coordinates {
    (1, 0) (2, 0) (3, 2) (4, 10) (5, 50) (6, 110)
    (7, 190) (8, 240) (9, 160) (10, 30) (11, 2) (12, 0)
};
\addlegendentry{Pluies Maroua}
% Température Yaoundé (courbe popdark continue)
\addplot[popdark, thick, mark=*, mark size=2, smooth] coordinates {
    (1, 24.5) (2, 25.2) (3, 25.0) (4, 24.5) (5, 24.0) (6, 23.2)
    (7, 22.5) (8, 22.3) (9, 22.8) (10, 23.0) (11, 23.8) (12, 24.2)
};
\addlegendentry{Temp. Yaoundé}
% Température Maroua (courbe poppurple pointillée)
\addplot[poppurple, thick, dashed, mark=square*, mark size=2, smooth] coordinates {
    (1, 26.5) (2, 29.0) (3, 32.0) (4, 33.5) (5, 32.0) (6, 29.5)
    (7, 27.0) (8, 26.0) (9, 26.5) (10, 28.0) (11, 27.5) (12, 26.0)
};
\addlegendentry{Temp. Maroua}
\end{axis}
% Deuxième axe pour l'échelle de température (invisible, juste pour l'échelle droite)
\begin{axis}[
    width=\linewidth,
    height=9cm,
    axis y line*=right,
    axis x line=none,
    ymin=15, ymax=35,
    xmin=0.5, xmax=12.5,
    xtick={1,2,3,4,5,6,7,8,9,10,11,12},
    xticklabels={},
    ylabel={Température (\si{\celsius})},
    ylabel style={popdark},
    tick label style={font=\small},
    label style={font=\small},
    hide x axis,
    y axis line style={popdark},
    y tick label style={popdark},
]
% Pas de données, juste l'échelle
\addplot[draw=none] coordinates {(1,20) (12,20)};
\end{axis}
\end{tikzpicture}
\legende{Diagramme ombrothermique comparatif (échelle $P = 2T$) : \textbf{Yaoundé} (Sud, climat équatorial atténué) et \textbf{Maroua} (Extrême-Nord, climat sahélien). Barres = précipitations (bleu/Orange). Courbes = températures (noir continu/pourpre pointillé). À Yaoundé, les barres dépassent presque toujours la courbe $P=2T$ (4 saisons, pas de vraie saison sèche stricte à l'échelle $2T$ sauf grande sèche). À Maroua, les barres sont sous la courbe $P=2T$ 8 mois/an (saison sèche marquée), nette supériorité en août seulement. Source : données moyennes ASECNA / Météo Cameroun.}
\end{popfigure}

\begin{exoresolu}[Analyse comparative Yaoundé vs Maroua]
\textbf{Énoncé :} À l'aide du diagramme ombrothermique ci-dessus, comparez les climats de Yaoundé et de Maroua en précisant pour chacun : type de climat, pluviométrie annuelle approximative, nombre et durée des saisons, températures, et conséquences sur la végétation naturelle.

\textbf{Solution détaillée :}

\noindent\textbf{1. Yaoundé (Centre, ~\SI{750}{m}) :}
\begin{itemize}
    \item \textbf{Type :} Climat équatorial de transition (ou équatorial atténué d'altitude).
    \item \textbf{Pluviométrie :} $\sum P \approx \SI{1 528}{mm}$ (somme des barres). Abondante, bien répartie.
    \item \textbf{Saisons :} \textbf{4 saisons} nettes. Deux saisons des pluies (mars--juin : $\sim$\SI{630}{mm} ; sept.--nov. : $\sim$\SI{640}{mm}) séparées par une petite saison sèche (juil.--août : $\sim$\SI{170}{mm}, mais $P > 2T$ ? Vérif : Juil $T=\SI{22.5}{\celsius} \rightarrow 2T=\SI{45}{mm}$, $P=\SI{70}{mm} > 45$ ; Août $T=\SI{22.3}{\celsius} \rightarrow 2T=\SI{44.6}{mm}$, $P=\SI{100}{mm} > 44.6$. \textbf{Attention :} À l'échelle $P=2T$, il n'y a \textbf{pas} de mois sec strict ($P < 2T$) en juillet-août. La « petite saison sèche » est une \emph{diminution relative} des pluies, pas une aridité. La grande saison sèche (déc.--fév.) voit $P$ faible mais $T$ élevée $\rightarrow$ $P < 2T$ possible en janv/fév. (Janv: $T=\SI{24.5}{\celsius}, 2T=\SI{49}{mm}, P=\SI{18}{mm} < 49$). Donc 2--3 mois secs ($P<2T$) en grande saison sèche seulement.
    \item \textbf{Températures :} Moyenne ~\SI{23,8}{\celsius}. Amplitude annuelle faible (~\SI{3}{\celsius}). Pas de stress thermique.
    \item \textbf{Végétation :} Forêt dense humide semi-caducifoliée (forêt atlantique biafréenne), exploitée pour le bois d'œuvre.
\end{itemize}

\noindent\textbf{2. Maroua (Extrême-Nord, ~\SI{350}{m}) :}
\begin{itemize}
    \item \textbf{Type :} Climat sahélien (tropical sec).
    \item \textbf{Pluviométrie :} $\sum P \approx \SI{794}{mm}$. Faible, très concentrée.
    \item \textbf{Saisons :} \textbf{2 saisons} très contrastées. Saison des pluies courte (juin--sept, 4 mois, pic août \SI{240}{mm}). Saison sèche très longue (oct.--mai, 8 mois). Mois secs ($P < 2T$) : oct. à mai (8 mois). Mois humides ($P > 2T$) : juil., août, sept. seulement.
    \item \textbf{Températures :} Moyenne ~\SI{28,5}{\celsius}. Amplitude annuelle marquée (~\SI{7,5}{\celsius}). Pic de chaleur avril--mai (> \SI{32}{\celsius}). Amplitude diurne forte.
    \item \textbf{Végétation :} Steppe sahélienne (herbacées annuelles, arbustes épineux \emph{Acacia}, \emph{Balanites}), galerie forestière le long des cours d'eau (Mayo Tsanaga, Logone).
\end{itemize}

\noindent\textbf{Synthèse :} Le gradient climatique Sud--Nord est extrême : facteur 2 sur les pluies (\SI{1 500}{mm} vs \SI{800}{mm}), inversion du régime des saisons (4 vs 2), augmentation de la continentalité (amplitude T). L'altitude (Yaoundé) adoucit l'équatorial ; la latitude (Maroua) assèche le tropical.
\end{exoresolu}

\begin{exercice}[Entraînement : Diagramme de Ngaoundéré]
Les données moyennes mensuelles pour Ngaoundéré (Adamaoua, \SI{1 100}{m}) sont :
\begin{center}
\begin{tabular}{|c|c|c|c|c|c|c|c|c|c|c|c|c|}\hline
Mois & J & F & M & A & M & J & J & A & S & O & N & D \\\hline
$T$ (\si{\celsius}) & 22,5 & 24,0 & 25,0 & 24,5 & 23,5 & 22,0 & 21,0 & 20,5 & 21,0 & 22,0 & 22,0 & 21,5 \\\hline
$P$ (mm) & 2 & 5 & 25 & 80 & 160 & 190 & 230 & 260 & 220 & 80 & 10 & 2 \\\hline
\end{tabular}
\end{center}
\begin{enumerate}
    \item Tracer le diagramme ombrothermique (échelle $P = 2T$) sur papier millimétré.
    \item Calculer la pluviométrie annuelle totale.
    \item Identifier les mois secs ($P < 2T$) et humides ($P > 2T$).
    \item Déduire le nombre de saisons et le type de climat.
    \item Expliquer l'influence de l'altitude sur les températures par rapport à Maroua (même latitude ~\SI{7}{\degree} N pour Ngaoundéré, Maroua à ~\SI{10,5}{\degree} N).
\end{enumerate}
\end{exercice}

\begin{corrige}[Exercice : Diagramme de Ngaoundéré]
\begin{enumerate}
    \item \emph{Tracé :} Courbe T en cloche (max \SI{25}{\celsius} mars/avril, min \SI{20,5}{\celsius} août). Barres P : nulles nov--mars, montée avril, pic août (\SI{260}{mm}), chute oct.
    \item $\sum P = \SI{1 264}{mm}$.
    \item Seuil $2T$ : ~\SIrange{41}{50}{mm}. Mois secs ($P < 2T$) : Nov, Déc, Jan, Fév, Mars (5 mois). Mois humides ($P > 2T$) : Avr--Oct (7 mois).
    \item \textbf{2 saisons nettes} : Saison sèche nov--mars (5 mois), saison des pluies avr--oct (7 mois). \textbf{Climat tropical humide d'altitude} (ou tropical de transition Adamaoua).
    \item Ngaoundéré (\SI{7,3}{\degree} N, \SI{1 100}{m}) vs Maroua (\SI{10,6}{\degree} N, \SI{350}{m}). Malgré une latitude plus basse (plus au sud, donc théoriquement plus chaud), l'altitude (\SI{1 100}{m} vs \SI{350}{m}) abaisse les températures de ~\SIrange{7}{8}{\celsius} (gradient ~\SI{0,6}{\celsius}/\SI{100}{m}), réduit l'amplitude annuelle, supprime les pics de chaleur extrêmes (> \SI{40}{\celsius} à Maroua vs < \SI{26}{\celsius} à Ngaoundéré). L'altitude favorise aussi la condensation orographique $\rightarrow$ pluies plus abondantes (\SI{1 264}{mm} vs \SI{794}{mm}).
\end{enumerate}
\end{corrige}

\courssub{Synthèse : le couple relief--climat structure le territoire}

\begin{propriete}[Relations relief--climat au Cameroun]
\begin{itemize}
    \item \textbf{Barrière orographique :} La dorsale camerounaise et les hautes terres de l'Ouest bloquent les alizés de SW (mousson) $\rightarrow$ surpluviométrie côtière (Debundscha), ombre pluviométrique à l'est (bassin du Mbam, plus sec).
    \item \textbf{Étagement altitudinal :} Sur le mont Cameroun et les hautes terres, la température baisse de ~\SI{0,6}{\celsius}/\SI{100}{m} $\rightarrow$ étagement végétal : forêt de basse altitude $\rightarrow$ forêt de montagne $\rightarrow$ prairie subalpine (sommet Fako).
    \item \textbf{Château d'eau :} L'Adamaoua (\SIrange{1 000}{1 500}{m}) capte les pluies de mousson et alimente 3 bassins versants (Atlantique, Niger, Lac Tchad). Son climat tropical d'altitude (fraîcheur, 2 saisons) permet l'élevage bovin (absence de mouche tsé-tsé).
    \item \textbf{Continentalité croissante vers le Nord :} Éloignement de l'océan + latitude + relief bas $\rightarrow$ air sec, harmattan, amplitude thermique forte, pluviométrie faible et irrégulière.
\end{itemize}
\end{propriete}

\begin{aretenir}[Essentiel à retenir pour le Bac]
\begin{itemize}
    \item \textbf{Relief :} 4 ensembles : Plaines Sud (0--650 m), Hautes terres Ouest + Dorsale (Mont Cameroun \SI{4 095}{m}), Adamaoua (\SIrange{1 000}{1 500}{m}, château d'eau), Plaines Nord (< 500 m).
    \item \textbf{Climats :} 3 domaines + microclimats. (1) Équatorial humide Sud (4 saisons, \SIrange{1 500}{4 000}{mm}). (2) Tropical humide altitude Adamaoua/Hautes terres (2 saisons, fraîcheur). (3) Soudanien/Sahélien Nord (saison sèche longue, \SIrange{400}{1 200}{mm}). (4) Microclimat Debundscha (~\SI{10 000}{mm}).
    \item \textbf{Ombothermique :} Échelle $P=2T$. Saison sèche = $P < 2T$. Yaoundé = 4 saisons (pas de mois sec strict à l'échelle $2T$ sauf grande sèche). Maroua = 8 mois secs. Ngaoundéré = 5 mois secs / 7 mois humides.
    \item \textbf{Couplage :} Relief $\rightarrow$ Climat (altitude, barrière) $\rightarrow$ Végétation/Sols/Hydrographie $\rightarrow$ Activités humaines.
\end{itemize}
\end{aretenir}

\coursec{La diversité des végétations et des sols}

Après avoir analysé la mosaïque des reliefs et la variété des climats qui structurent le territoire camerounais (leçon précédente), nous comprenons maintenant que ces deux facteurs agissent comme les architectes principaux des paysages biologiques. Le relief, en modifiant l'altitude et l'exposition, crée des microclimats ; le climat, par la distribution des pluies et des températures, dicte la nature de la couverture végétale. Cette section explore la réponse du vivant à ces contraintes physiques : comment la végétation s'organise en grands domaines, quels sols se forment sous ces couverts, et quelle faune y trouve refuge.

\courssub{Le domaine forestier du Sud : le royaume de l'humidité}

Le sud du Cameroun, correspondant à la région de climat équatorial (guinéen et camerounais), abrite l'un des massifs forestiers les plus remarquables d'Afrique : la forêt dense humide sempervirente. C'est le prolongement occidental du bassin du Congo.

\begin{definition}[Forêt dense humide équatoriale]
Formation végétale climax des climats chauds et humides sans saison sèche marquée (ou très courte, < 2 mois), caractérisée par une stratification verticale complexe (strates émergente, principale, intermédiaire, herbacée), une biodiversité exceptionnelle et une productivité biologique maximale.
\end{definition}

Cette forêt ne se laisse pas réduire à un « bloc vert » uniforme. On y distingue des faciès précis :
\begin{itemize}
    \item \textbf{La forêt de plaine et de piémont :} Sur les sols profonds du plateau sud-camerounais, elle atteint son développement optimal. C'est le domaine des \cle{essences précieuses} : l'\textbf{ébène} (\emph{Diospyros crassiflora}), bois noir dense pour l'ébénisterie ; l'\textbf{acajou} (\emph{Entandrophragma cylindricum} -- sapelli, \emph{Khaya ivorensis} -- acajou d'Afrique), bois rouge d'œuvre prisé ; l'\textbf{iroko} (\emph{Milicia excelsa}), imputrescible pour la construction navale ; l'\textbf{ayous} (\emph{Triplochiton scleroxylon}), bois blanc léger pour le contreplaqué. Ces espèces « nobles » ont fait la richesse forestière du Cameroun mais subissent aujourd'hui une pression d'exploitation intense.
    \item \textbf{La forêt littorale et les mangroves :} À l'interface terre-mer (embouchures du Wouri, du Nyong, de la Sanaga, Rio del Rey), la forêt s'adapte à la salinité et aux marées. La \cle{mangrove} (palétuviers \emph{Rhizophora} spp., \emph{Avicennia} spp.) forme un rideau protecteur contre l'érosion côtière, une nurserie pour la pêche (crevettes, poissons) et un piège à carbone majeur (« carbone bleu »).
\end{itemize}

\begin{exemplebox}[Observation de terrain : la strate émergente]
En survolant la réserve du Dja ou le parc de Lobéké, on aperçoit des géants (60--70 m) comme le \emph{Entandrophragma} (sapelli) ou le \emph{Piptadeniastrum africanum} (dabéma) qui percent la canopée continue à 30--40 m. Leur stratégie : capter la lumière au-dessus de la concurrence. Leur fruit/légèreté permet une dispersion anémochore (par le vent) sur de grandes distances.
\end{exemplebox}

\courssub{Le domaine des savanes : l'adaptation à la saison sèche}

Au nord de la latitude 6°--7° N, la durée de la saison sèche s'allonge (4 à 7 mois). La forêt recule, cédant la place aux \cle{savanes}, formations où la strate herbacée domine, entrecoupée d'arbres plus ou moins denses.

\begin{definition}[Savane]
Formation végétale herbacée (graminées hautes) parsemée d'arbres et d'arbustes, propre aux climats tropicaux à saison sèche marquée (soudanien, soudan-sahélien). La strate herbacée sèche et brûle en saison sèche ; la strate ligneuse perd ses feuilles (caducité) ou possède des écorces épaisses (résistance au feu).
\end{definition}

On distingue un gradient nord-sud net :
\begin{itemize}
    \item \textbf{Savane arborée et herbeuse de l'Adamaoua et de l'Ouest (soudano-guinéenne) :} Transition forêt-savane. Arbres denses (\emph{Isoberlinia doka}, \emph{Afzelia africana}, \emph{Daniellia oliveri}) formant par endroits des « boisements » fermés. Herbes hautes (\emph{Andropogon gayanus}, \emph{Hyparrhenia}). Altitude adoucissant le climat (Adamaoua : 1000--1500 m).
    \item \textbf{Savane soudanienne du Nord (régions Nord, Nord-Est) :} Arbres plus espacés, plus petits (\emph{Acacia seyal}, \emph{Balanites aegyptiaca}, \emph{Vitellaria paradoxa} -- karité, \emph{Parkia biglobosa} -- néré). Herbes plus courtes. Parcs nationaux de la Bénoué, Bouba Ndjida, Faro.
    \item \textbf{Steppe sahélienne de l'Extrême-Nord (région Extrême-Nord, bassin du Lac Tchad) :} Formation xérique. Arbres rares, épineux (\emph{Acacia nilotica}, \emph{Balanites aegyptiaca}, \emph{Calotropis procera}, \emph{Leptadenia pyrotechnica}). Herbes annuelles à cycle très court. Risque de désertification élevé.
\end{itemize}

\begin{aretenir}[Gradient phytogéographique Nord-Sud]
\begin{center}
\begin{tabularx}{\linewidth}{|l|X|X|X|}\hline
\textbf{Domaine} & \textbf{Climat} & \textbf{Essences caractéristiques} & \textbf{Strate herbacée} \\ \hline
Forêt dense & Équatorial (Guinéen/Camerounais) & Ayous, Iroko, Acajou, Ébène, Sapelli & Faible (ombre) \\ \hline
Savane arborée (Adamaoua/Ouest) & Soudano-guinéen & Isoberlinia, Afzelia, Daniellia, Karité & Haute (Andropogon, Hyparrhenia) \\ \hline
Savane soudanienne (Nord) & Soudanien & Acacia seyal, Karité, Néré, Baobab & Moyenne \\ \hline
Steppe sahélienne (Extrême-Nord) & Sahélien & Acacia nilotica, Balanites, Leptadenia & Courte, annuelle \\ \hline
\end{tabularx}
\end{center}
\end{aretenir}

\courssub{La végétation d'altitude : les « îles » fraîches}

Sur les hautes terres de l'Ouest (Bamenda Highlands, Monts Mandara) et sur le \cle{Mont Cameroun} (4095 m), l'altitude impose un climat de montagne : températures basses, fortes pluies orographiques, vents violents. La végétation s'étage :

\begin{itemize}
    \item \textbf{Étage submontagnard (1000--1800 m) :} Forêt de transition, moins haute, épiphytes abondants (mousses, orchidées, fougères).
    \item \textbf{Étage montagnard (1800--2500 m) :} Forêt nuageuse (\emph{Podocarpus milanjianus}, \emph{Olea capensis}, \emph{Schefflera abyssinica}), troncs tordus, tapis de mousses.
    \item \textbf{Étage subalpin / alpin (au-dessus de 2500 m, sommet Mont Cameroun) :} \cle{Prairies et pelouses} (bambous \emph{Arundinaria alpina} sur le Mont Cameroun, landes à \emph{Erica mannii}, \emph{Helichrysum}, \emph{Lobelia} géants). Absence d'arbres. Endémisme fort (ex: \emph{Impatiens} spp., chiroptères, francolin de Cameroun).
\end{itemize}

\begin{savaistu}[Le Mont Cameroun, laboratoire naturel]
Le Mont Cameroun est l'un des rares volcans actifs au monde où l'on observe une succession primaire complète : laves récentes stériles $\rightarrow$ mousses/lichens $\rightarrow$ fougères $\rightarrow$ bambous $\rightarrow$ forêt montagnarde. C'est un site clé pour étudier la recolonisation végétale et l'évolution des sols volcaniques (andosols).
\end{savaistu}

\courssub{Les sols : le miroir du climat et de la végétation}

Le sol n'est pas un simple support inerte ; c'est le fruit de l'altération de la roche mère sous l'action du climat (température, eau) et des organismes (végétation, faune du sol, micro-organismes). Au Cameroun, la zonalité des sols suit fidèlement celle du climat et de la végétation.

\begin{definition}[Sol ferrallitique (Sud-Cameroun)]
Sol rouge, profond (souvent > 10 m), très ancien, caractéristique des régions intertropicales chaudes et très humides. Il résulte d'un \cle{lessivage} intense : les pluies abondantes lessivent les bases (Ca, Mg, K), la silice et une partie de l'alumine, ne laissant s'accumuler que les oxydes de fer (hématite = rouge) et d'alumine (gibbsite). \textbf{Pauvres en nutriments, acides (pH 4--5), à faible capacité d'échange cationique (CEC).} La fertilité est entièrement portée par le cycle organique (humus de surface).
\end{definition}

\begin{definition}[Sol ferrugineux (Savanes)]
Sol rouge-brun à jaune, moins profond, moins lessivé que le ferrallitique (saison sèche stoppe la migration). Présence d'une \cle{croûte ferrugineuse} (cuirasse / pétroferrique) en surface ou à faible profondeur, témoin d'une alternance humide/sec intense. Meilleure réserve en bases que le ferrallitique, mais sensibilité à l'érosion (croûte = imperméabilité $\rightarrow$ ruissellement).
\end{definition}

\begin{itemize}
    \item \textbf{Sols hydromorphes :} Plaines inondables du Logone, de la Bénoué, rives du Lac Tchad, bas-fonds (marais). Saturés en eau une partie de l'année $\rightarrow$ gleyfication (couleurs bleutées/grises, réduction du fer). Riches en matière organique et nutriments $\rightarrow$ fort potentiel agricole (riziculture, cultures de décrue) mais drainage difficile.
    \item \textbf{Sols volcaniques (Andosols) :} Hautes terres de l'Ouest (Bamenda, Dschang, Bafoussam, flancs Mont Cameroun). Jeunes, issus de cendres et laves basaltiques. \textbf{Très fertiles} : structure grumeleuse, forte porosité, réserve en eau élevée, richesse en minéraux primaires (K, Mg, Ca) et forte fixation du phosphore (allophane). C'est le « grenier » du Cameroun (café arabica, légumes, pommes de terre).
    \item \textbf{Sols halomorphes :} Bas-fonds du Logone, Lac Tchad. Salinisation/sodification par évaporation capillaire. Difficiles pour l'agriculture sans aménagements (lessivage, drainage).
\end{itemize}

\begin{experience}[Observation de profils de sol (TP terrain ou documents)]
\textbf{Matériel :} Photos de coupes de profils (Sud : ferrallitique ; Adamaoua : ferrugineux ; Logone : hydromorphe ; Ouest : volcanique) ou carottes de prélèvement.
\textbf{Protocole :}
\begin{enumerate}
    \item Identifier les horizons (O, A, E, B, C) par leur couleur (référentiel Munsell), texture, structure, présence de nodules/concrétions, racines.
    \item Mesurer le pH (eau et KCl) sur chaque horizon.
    \item Tester la réaction à l'acide chlorhydrique (carbonates) et au peroxyde d'hydrogène (matière organique).
\end{enumerate}
\textbf{Observations attendues :}
\begin{itemize}
    \item \textit{Ferrallitique :} Horizon A fin, sombre ; transition diffuse vers B jaune-rouge uniforme, profond, structure microgrumeleuse, pH 4--5, pas de concrétions dures en surface.
    \item \textit{Ferrugineux :} Horizon A plus clair ; horizon B avec concrétions ferrugineuses (gravelle) ou cuirasse continue (pétroferrique) ; pH 5--6.
    \item \textit{Volcanique :} Horizon A épais, noir, humifère, structure grumeleuse forte ; horizon B brun foncé, pH 6--7.
    \item \textit{Hydromorphe :} Taches de rouille (oxydation) et zones grises/bleues (réduction/gley) dans le profil ; matière organique accumulée en surface.
\end{itemize}
\textbf{Interprétation :} Relier la morphologie du profil aux processus pédogénétiques dominants (lessivage, rubéfaction, gleyfication, altération volcanique) et au climat/végétation.
\end{experience}

\courssub{Faune associée : la biodiversité comme indicateur de milieu}

La répartition de la faune sauvage suit strictement les grands ensembles biogéographiques. Le Cameroun, par sa position charnière, abrite une faune à la fois forestière (bassin du Congo) et soudano-sahélienne.

\begin{center}
\begin{tabularx}{\linewidth}{|l|X|X|}\hline
\textbf{Domaine biogéographique} & \textbf{Espèces emblématiques (Grands mammifères)} & \textbf{Statut / Aires protégées clés} \\ \hline
Forêt dense (Sud, Est) & Gorille des plaines de l'Ouest (\emph{Gorilla gorilla gorilla}), Chimpanzé d'Afrique centrale (\emph{Pan troglodytes troglodytes}), Éléphant de forêt (\emph{Loxodonta cyclotis} -- plus petit, défenses droites), Bongo, Sitatunga, Mandrill, Drill, Pangolin géant. & Parcs : Lobéké, Boumba Bek, Nki, Campo Ma'an, Dja (Réserve de biosphère UNESCO). Braconnage (ivoire, viande de brousse) = menace majeure. \\ \hline
Savane arborée / Soudanienne (Adamaoua, Nord) & Éléphant de savane (\emph{Loxodonta africana}), Lion, Léopard, Buffle, Hippotrague, Cob de Buffon, Girafe (kordofan), Lycaon, Guépard (rare). & Parcs : Bénoué, Bouba Ndjida, Faro, Mbam et Djerem (transition). Conflits Homme-Faune (cultures). \\ \hline
Steppe / Sahélien (Extrême-Nord) & Antilope dama, Oryx algazelle (réintroduit), Gazelle dorcas, Autruche, Damans, Chacal. Oiseaux migrateurs (zone Ramsar Waza-Logone). & Parc : Waza (Ramsar, Réserve de biosphère). Sécheresses, pression pastorale, insécurité (Boko Haram) = menaces. \\ \hline
Montagne (Ouest, Mont Cameroun) & Endémiques : Rat à queue blanche du Mont Cameroun, Musaraignes, Chiroptères, Oiseaux (Francolin de Cameroun, Turaco de Bannerman). & Mont Cameroun (Parc national), Kilum-Ijim, Monts Mandara. Déforestation, agriculture. \\ \hline
\end{tabularx}
\end{center}

\begin{savaistu}[Le Parc National de Waza : un patrimoine menacé]
Créé en 1934 (réserve de faune) puis parc national en 1968, le parc de Waza (170 000 ha) dans l'Extrême-Nord est une zone humide soudano-sahélienne (plaine d'inondation du Logone). Il est classé \textbf{Ramsar} (1987) et \textbf{Réserve de biosphère UNESCO} (1979). Il abrite l'une des dernières populations de girafes de Kordofan, des éléphants, des lions, et une avifaune exceptionnelle (pélicans, marabouts, canards). La construction du barrage de Maga (amont) a modifié le régime hydraulique (assèchement partiel), et l'insécurité régionale a freiné le tourisme et la surveillance.
\end{savaistu}

\courssub{Relations Climat -- Végétation -- Sol -- Activités humaines : une chaîne fonctionnelle}

Ces quatre composantes ne sont pas juxtaposées ; elles interagissent en boucle rétroactive. Comprendre cette chaîne est la clé de la géographie appliquée (aménagement, agriculture, conservation).

\begin{popfigure}
\begin{tikzpicture}[
    node distance=1.2cm and 1.5cm,
    box/.style={draw=popblue, thick, rounded corners, fill=popblue!10, minimum width=2.8cm, minimum height=1.2cm, align=center, font=\small},
    arrow/.style={-Stealth, thick, popdark},
    link/.style={font=\footnotesize, popdark, midway, sloped}
]
\node[box, align=center] (climat){CLIMAT\\ (Température, Pluviométrie,\\ Saisonnalité, Évapotranspiration)};
\node[box, right=of climat, align=center] (veg){VÉGÉTATION\\ (Type, Structure, Biomasse,\\ Cycle phénologique)};
\node[box, right=of veg, align=center] (sol){SOL\\ (Type, Fertilité, Profondeur,\\ Structure, Stock carbone)};
\node[box, below=of veg, align=center] (act){ACTIVITÉS HUMAINES\\ (Agriculture, Élevage, Forêt,\\ Urbanisme, Conservation)};
\draw[arrow] (climat) -- node[link, above] {Détermine le potentiel} (veg);
\draw[arrow] (veg) -- node[link, above, align=center]{Litière, Racines,\\ Protection érosion} (sol);
\draw[arrow] (sol) -- node[link, below] {Support, Réserve eau/nutriments} (veg);
\draw[arrow] (veg) -- node[link, left, align=center]{Évapotranspiration,\\ Rugosité, Albedo} (climat);
\draw[arrow] (act) -- node[link, left, align=center]{Déforestation, Feux,\\ Intrants, Amendements} (veg);
\draw[arrow] (act) -- node[link, right, align=center]{Labour, Irrigation,\\ Érosion, Pollution} (sol);
\draw[arrow, dashed] (act) -- node[link, above right, align=center]{Émissions GES,\\ Changement d'affectation} (climat);
\end{tikzpicture}
\legende{Schéma des interactions systémiques entre Climat, Végétation, Sol et Activités humaines. Les flèches pleines indiquent les interactions directes fortes ; la flèche pointillée l'impact global à long terme.}
\end{popfigure}

\begin{methode}[Analyser un ensemble biogéographique (Démarche d'investigation)]
Pour comprendre un paysage rural ou naturel au Cameroun, appliquez cette grille de lecture en 4 étapes :
\begin{enumerate}
    \item \textbf{Diagnostiquer le climat :} Localiser la station météo la plus proche. Noter $P$ (mm/an), $T$ (°C), nombre de mois secs ($P < 50$ mm), régime de pluies.
    \item \textbf{Identifier la formation végétale :} Observer la strate dominante (forêt, savane, steppe, prairie), la densité, les espèces indicatrices (ex: \emph{Isoberlinia} = savane soudano-guinéenne ; \emph{Acacia seyal} = soudanien ; \emph{Rhizophora} = mangrove).
    \item \textbf{Inférer le type de sol :} Utiliser la règle zonale (Ferrallitique $\leftrightarrow$ Forêt équatoriale ; Ferrugineux $\leftrightarrow$ Savane ; Hydromorphe $\leftrightarrow$ Bas-fond ; Volcanique $\leftrightarrow$ Hautes terres). Vérifier par indices de terrain (couleur, texture, présence de cuirasse, culture pratiquée).
    \item \textbf{Évaluer les potentialités et risques :} Croiser sol + climat + pente + population $\rightarrow$ Aptitude culturale (riz, coton, café, pâturage), risque érosion, besoin en amendements, enjeu conservation.
\end{enumerate}
\end{methode}

\coursec{Les grands groupes humains du Cameroun}

La diversité physique (relief, climat, végétation, sols) structure l'installation et les activités des populations. Le Cameroun compte plus de 250 groupes ethniques, regroupés en grands ensembles linguistiques et culturels dont la répartition géographique recoupe largement les ensembles biogéographiques.

\courssub{Les grands ensembles linguistiques et leur répartition}

\begin{itemize}
    \item \textbf{Les populations bantoues (Zone Sud, Centre, Est, Littoral, Sud-Ouest) :} Regroupent les Beti-Fang (Ewondo, Boulou, Fang), les Bassa, les Douala, les Bakoko, les Makaa-Njem, les Kako, les Konabembe. \textbf{Milieu :} Forêt dense, climat équatorial. \textbf{Activités traditionnelles :} Agriculture itinérante sur brûlis (manioc, macabo, plantain, igname), cueillette, chasse, pêche fluviale. Habitat dispersé (cases en matériaux végétaux) ou groupé le long des axes routiers.
    \item \textbf{Les populations semi-bantoues (Zone Ouest, Nord-Ouest, Adamaoua occidental) :} Bamileke, Bamoun, Tikar, Bafia, Widikum, Mbo. \textbf{Milieu :} Hautes terres volcaniques (climat tempéré, sols fertiles), piémonts forestiers. \textbf{Activités :} Agriculture intensive de haute altitude (café arabica, maïs, haricot, pomme de terre, légumes), élevage bovin (race locale, métis), artisanat, commerce. Habitat groupé (chefferies), forte densité de population.
    \item \textbf{Les populations soudano-sahéliennes (Zone Nord, Adamaoua oriental, Extrême-Nord) :} 
    \begin{itemize}
        \item \textit{Groupe tchadique (Kanouri, Kotoko, Mousgoum, Massa, Arabes Choa) :} Bassin du Lac Tchad, plaine du Logone. Pêche, agriculture de décrue (riz, sorgho), élevage, commerce transfrontalier.
        \item \textit{Groupe adamawa-oubanguien (Mbum, Dii, Gbaya, Mboum, Tupuri, Mundang, Toupouri) :} Adamaoua, Bénoué, Mayo-Kebbi. Agriculture (sorgho, maïs, mil, coton), élevage, chasse.
        \item \textit{Peuls (Fulbe) :} Présents sur tout le Nord, l'Adamaoua et les piémonts de l'Ouest. \textbf{Activité majeure :} Élevage bovin extensif (zébu), transhumance saisonnière (recherche pâture/eau). Rôle économique central (approvisionnement en viande des villes du Sud). Sédentarisation croissante (agro-pastoralisme).
    \end{itemize}
    \item \textbf{Les populations pygmées (Baka à l'Est, Bakola/Bagyeli au Sud, Bedzang au Centre) :} Premiers habitants de la forêt. \textbf{Mode de vie :} Chasse-cueillette nomade/séminomade, connaissance intime de la forêt (pharmacopée, miel, chenilles). \textbf{Enjeu actuel :} Sédentarisation forcée, perte d'accès aux ressources forestières, marginalisation, non-reconnaissance des droits coutumiers.
\end{itemize}

\begin{aretenir}[Répartition spatiale des grands groupes humains]
\begin{center}
\begin{tabularx}{\linewidth}{|l|X|X|}\hline
\textbf{Ensemble} & \textbf{Localisation principale (Régions)} & \textbf{Activité dominante traditionnelle} \\ \hline
Bantous & Sud, Centre, Est, Littoral, Sud-Ouest & Agriculture itinérante (tubercule), cueillette, chasse \\ \hline
Semi-Bantous & Ouest, Nord-Ouest, Adamaoua (Ouest) & Agriculture intensive (café, vivriers), élevage, commerce \\ \hline
Soudano-Sahéliens (Tchadiens, Adamawa) & Nord, Extrême-Nord, Adamaoua (Est) & Agriculture pluviale (céréales, coton), pêche, élevage \\ \hline
Peuls (Fulbe) & Nord, Adamaoua, Nord-Ouest, Ouest & Élevage bovin transhumant / agro-pastoralisme \\ \hline
Pygmées (Baka, Bakola, Bedzang) & Est, Sud, Centre (forêts) & Chasse-cueillette, exploitation produits forestiers non ligneux \\ \hline
\end{tabularx}
\end{center}
\end{aretenir}

\courssub{Dynamiques démographiques et recompositions spatiales}

\begin{itemize}
    \item \textbf{Forte croissance démographique :} Taux d'accroissement $\approx$ 2,6 \%/an. Population jeune (moins de 25 ans : $> 60\%$). Pression foncière accrue sur les zones fertiles (Ouest, piémonts, plaines alluviales).
    \item \textbf{Migrations internes :} 
    \begin{itemize}
        \item \textit{Sud $\rightarrow$ Sud (front pionnier) :} Vers l'Est (Sangmélima, Yokadouma) et le Sud-Ouest (Kumba, Mamfé) pour cultures de rente (cacao, café, palme à huile) et exploitation forestière.
        \item \textit{Nord $\rightarrow$ Sud (main-d'œuvre) :} Vers les plantations (CDC, SOCAPALM, HEVECAM) et les villes (Douala, Yaoundé).
        \item \textit{Montagnes (Ouest) $\rightarrow$ Villes / Littoral / Adamaoua :} Surpopulation rurale, recherche de terres.
    \end{itemize}
    \item \textbf{Urbanisation rapide :} Taux d'urbanisation $\approx$ 58 \% (2023). Douala (capitale économique, $\sim$ 4 M hab.) et Yaoundé (capitale politique, $\sim$ 4 M hab.) concentrent l'essentiel. Problèmes : habitat précaire, chômage, gestion déchets/eau, étalement sur zones agricoles/périlleuses.
\end{itemize}

\coursec{L'unité dans la diversité : atouts et défis}

La « diversité » décrite ci-dessus est dynamique et fragile. Les pressions anthropiques s'accélèrent, modifiant la donne climatique locale et globale.

\courssub{Menaces sur les milieux : l'empreinte humaine}

\begin{attention}[Erreur fréquente : la désertification, une fatalité climatique ?]
\emph{Attention :} On attribue souvent la dégradation des terres au Grand Nord (Extrême-Nord) au seul « climat qui se sèche ». \textbf{Faux.} La désertification est définie par l'ONU (CNULCD) comme la « dégradation des terres en zones arides, semi-arides et sub-humides sèches résultant de \textbf{facteurs divers, notamment des variations climatiques et des activités humaines} ».
Au Cameroun septentrional, les activités humaines \textbf{accélèrent dramatiquement} le processus :
\begin{itemize}
    \item \textbf{Surpâturage :} Densité de bétail $>$ capacité de charge des parcours $\rightarrow$ destruction de la strate herbacée, mise à nu du sol $\rightarrow$ érosion éolienne (vents d'harmattan) et hydrique (orages violents).
    \item \textbf{Bois de chauffe / charbon :} Prélèvement d'énergie bois $>$ productivité du boisement $\rightarrow$ déboisement des corridors, disparition des arbres fourragers et fruitiers (karité, néré, balanites).
    \item \textbf{Agriculture itinérante sur brûlis mal maîtrisée / cultures de rente (coton) :} Labour à nu, perte de matière organique, tassement, érosion.
    \item \textbf{Barrages amont (Maga, Yagoua, Lagdo) :} Modification du régime d'inondation du Logone $\rightarrow$ assèchement des plaines de décrue (pâturages, pêche, recharge nappe) $\rightarrow$ extension des zones salées.
\end{itemize}
La lutte contre la désertification (Grande Muraille Verte, reboisement, agroforesterie, gestion concertée des parcours) est une priorité du Plan National de Développement (PND 2020-2030).
\end{attention}

\begin{aretenir}[Principales menaces par domaine]
\begin{itemize}
    \item \textbf{Forêt dense (Sud/Est) :} Exploitation forestière industrielle (grumes) + artisanale (scieries mobiles) + agriculture itinérante sur brûlis + agro-industrie (palmier à huile, hévéa) + mines (fer, cobalt) + braconnage (ivoire, pangolin). \textbf{Taux de déforestation :} $\approx$ 0,6--1\%/an (FAO), soit $\sim$ 200 000 ha/an.
    \item \textbf{Savanes (Adamaoua, Nord) :} Feux de brousse précoces/tardifs (chasse, pâturage, miel) $\rightarrow$ rajeunissement de la strate ligneuse, érosion, émissions CO$_2$. Extension cultures (coton, maïs) $\rightarrow$ fragmentation habitats faune.
    \item \textbf{Steppe (Extrême-Nord) :} Désertification (voir encadré Attention), insécurité, pression réfugiés/déplacés sur ressources (bois, eau, terre).
    \item \textbf{Montagnes (Ouest) :} Érosion hydrique massive (pentes fortes + pluies intenses + déforestation cultures), tassement sols volcaniques, pollution intrants (pesticides horticulture).
    \item \textbf{Mangroves/Littoral :} Coupe bois de chauffe/fumage poisson, pollution hydrocarbures (Douala), urbanisation galopante, élévation niveau mer.
\end{itemize}
\end{aretenir}

\begin{exemplebox}[Chiffres clés : la forêt camerounaise en 2023-2024]
\begin{itemize}
    \item Superficie forestière totale : \textbf{$\approx$ 20 millions d'hectares} (près de \textbf{40 \%} du territoire national).
    \item Forêts de production permanente (UFA, ventes de coupe) : $\approx$ 7,5 M ha.
    \item Aires protégées (forestières) : $\approx$ 3,5 M ha (Parcs, Réserves, Sanctuaires).
    \item Forêts communautaires / forêts communales : $\approx$ 1,5 M ha (gestion décentralisée).
    \item Contribution au PIB : $\approx$ 6--8 \% (secteur formel), bien plus si secteur informel inclus.
    \item Emplois directs/indirects : $\approx$ 45 000 directs, 200 000+ indirects.
\end{itemize}
\emph{Source : MINFOF, FAO/FRA 2020, ITTO, rapports annuels.}
\end{exemplebox}

\courssub{Exercice résolu : Analyse intégrée d'un terroir}

\begin{exoresolu}[Terroir de la plaine de la Bénoué (région du Nord)]
\textbf{Énoncé :} À partir des données suivantes, caractérisez l'ensemble biogéographique de la plaine de la Bénoué (alt. 200--300 m) et proposez les principales orientations d'aménagement agricole durable.
\begin{itemize}
    \item Climat : Soudanien. $P \approx 1000$ mm/an. Saison sèche 5--6 mois (nov--avr). $T$ moy. 28°C.
    \item Végétation naturelle : Savane arborée à \emph{Isoberlinia doka}, \emph{Anogeissus leiocarpus}, \emph{Vitellaria paradoxa} (karité). Galerie forestière le long Bénoué.
    \item Sol : Ferrugineux lessivé (sol tropical lessivé) sur alluvions anciennes. Croûte ferrugineuse affleurante par endroits. Bas-fonds : sols hydromorphes (plans d'eau temporaires « mares »).
    \item Faune : Parc de la Bénoué (hippopotames, cobs, buffles, lions, éléphants).
    \item Activités actuelles : Coton (SODECOTON), maïs, sorgho, élevage bovin (Peuls), pêche, cueillette karité/néré.
\end{itemize}

\textbf{Solution détaillée :}
\tcblower
\textbf{1. Caractérisation biogéographique :}
C'est un \textbf{ensemble biogéographique de transition soudano-guinéen / soudanien} (plaine alluviale de la Bénoué).
\begin{itemize}
    \item \textit{Climat $\rightarrow$ Végétation :} 1000 mm avec 5-6 mois secs $\rightarrow$ savane arborée (caducifoliée). La galerie forestière marque l'hydromorphie permanente (nappe phréatique haute).
    \item \textit{Végétation $\rightarrow$ Sol :} Feuilles caduques $\rightarrow$ litière annuelle $\rightarrow$ humification modérée $\rightarrow$ sol ferrugineux (lessivage saisonnier, formation croûte). Bas-fonds $\rightarrow$ gleyfication $\rightarrow$ sols hydromorphes.
    \item \textit{Sol $\rightarrow$ Végétation/Agriculture :} Croûte ferrugineuse = obstacle racinaire + ruissellement. Bas-fonds = réserve en eau + fertilité naturelle (limons).
\end{itemize}

\textbf{2. Potentialités / Contraintes :}
\begin{center}
\begin{tabularx}{\linewidth}{|l|X|X|}\hline
\textbf{Milieu} & \textbf{Potentialités} & \textbf{Contraintes / Risques} \\ \hline
Interfluves (sols ferrugineux) & Coton, maïs, sorgho, arachide ; pâturage saisonnier & Croûte ferrugineuse (peu profond), érosion laminaire/gullies, acidité, faible CEC, dépendance engrais. \\ \hline
Bas-fonds / Bas-fonds alluviaux & Riz pluvial/irrigué, légumes (contre-saison), pâturage sec, pêche & Ennoyage, salinisation risque si irrigation mal gérée, conflits agriculteurs/éleveurs. \\ \hline
Galeries forestières / Parc & Biodiversité, tourisme (photo-safari), services écosystémiques (régulation eau, carbone) & Braconnage, feux de brousse, pression riverains (bois, cultures), corridor faune fragmenté. \\ \hline
\end{tabularx}
\end{center}

\textbf{3. Orientations d'aménagement durable (Schéma Directeur d'Aménagement) :}
\begin{enumerate}
    \item \textbf{Zonage strict :} Sanctuariser le Parc de la Bénoué + zones tampons (chasse contrôlée / tourisme). Définir des couloirs de transhumance fixes (balisés) pour éviter cultures/parc.
    \item \textbf{Gestion des sols interfluves :} \textbf{Agriculture de conservation (AC)} : non-labour, couverture permanente (mucuna, crotalaire), rotation coton/céréales/légumineuses. Brisure mécanique de la croûte (scarificateur) si nécessaire. Apport matière organique (compost, fumure) pour remonter CEC/pH.
    \item \textbf{Valorisation bas-fonds :} Aménagement hydro-agricole maîtrisé (maîtrise de l'eau, drainage) pour riziculture intensive + maraîchage. Association riz-pisciculture. Gestion concertée eau (Comités de bassin).
    \item \textbf{Filières locales durables :} Structuration filière karité/néré (beurre, huile) $\rightarrow$ revenus femmes $\rightarrow$ incitation protection arbres parcs (agroforesterie). Apiculture (miel de savane).
    \item \textbf{Gestion du feu :} Feux dirigés précoces (déc--janv) par comités villageois pour protéger pâturages tardifs et limite parc.
\end{enumerate}
\textbf{Conclusion :} La durabilité repose sur l'articulation \textbf{Parc (conservation) -- Bas-fonds (intensification maîtrisée) -- Interfluves (agroécologie/AC) -- Pâturage géré}, avec une gouvernance locale forte (COGER, COVAREF).
\end{exoresolu}

\courssub{Illustration : Superficie des principaux parcs nationaux du Cameroun}

Le graphique ci-dessous donne un ordre de grandeur des surfaces consacrées à la conservation intégrale (parcs nationaux), témoins de la diversité des écosystèmes protégés.

\begin{popfigure}
\begin{tikzpicture}
\begin{axis}[
    ybar,
    bar width=12pt,
    width=\linewidth,
    height=10cm,
    enlargelimits=0.15,
    ylabel={Superficie (km$^2$)},
    ylabel style={font=\small},
    xlabel={Parcs nationaux (Échantillon)},
    xlabel style={font=\small},
    symbolic x coords={Waza, Bouba Ndjida, Bénoué, Faro, Korup, Lobéké, Campo Ma'an, Mbam\&Djerem, Mont Cameroun, Douala-Edéa},
    xtick=data,
    xticklabel style={rotate=45, anchor=east, font=\tiny},
    ymin=0,
    ymax=3000,
    ymajorgrids=true,
    grid style=dashed,
    title={Superficie approximative des principaux Parcs Nationaux du Cameroun},
    title style={font=\small\bfseries, yshift=0.3cm},
    nodes near coords,
    nodes near coords align={vertical},
    nodes near coords style={font=\tiny, rotate=90, anchor=west, yshift=2pt},
    popcurve/.style={draw=popblue, fill=popblue!30},
]
\addplot[popcurve] coordinates {
    (Waza,1700)
    (Bouba Ndjida,2200)
    (Bénoué,1800)
    (Faro,3300)
    (Korup,1260)
    (Lobéké,2178)
    (Campo Ma'an,2640)
    (Mbam\&Djerem,4165)
    (Mont Cameroun,581)
    (Douala-Edéa,2629)
};
\end{axis}
\end{tikzpicture}
\legende{Superficies approximatives (km$^2$) des principaux parcs nationaux. Source : MINFOF, WDPA. Note : Mbam et Djerem est le plus vaste ; le Mont Cameroun est le plus petit de cet échantillon mais à fort endémisme.}
\end{popfigure}

\begin{exercice}[Application : Diagnostic d'un terroir]
Vous êtes consultant pour une ONG d'appui au développement rural dans l'arrondissement de \textbf{Bankim} (région de l'Adamaoua, limite Nord-Ouest).
\begin{enumerate}
    \item Données : Alt. 1100 m. Climat soudano-guinéen ($P \approx 1500$ mm, 5 mois secs). Végétation : savane arborée à \emph{Isoberlinia}, \emph{Afzelia}, \emph{Daniellia}. Sols : ferrugineux lessivés (cuirasse profonde). Population : Mbororo (éleveurs) + cultivateurs (maïs, igname, patate). Parc de Mbam et Djerem au Sud.
    \item \textbf{Questions :}
    \begin{enumerate}
        \item Situez Bankim dans le gradient biogéographique national (climat, végétation, sol).
        \item Identifiez les deux principaux conflits d'usage du sol actuels ou potentiels.
        \item Proposez 3 mesures techniques d'agroforesterie ou de gestion des parcours adaptées au sol ferrugineux et à la cohabitation agriculture/élevage.
        \item Justifiez pourquoi la création de « corridors écologiques » vers le parc de Mbam et Djerem est cruciale pour la viabilité des populations d'éléphants et de grands ongulés.
    \end{enumerate}
\end{enumerate}
\end{exercice}

\begin{corrige}[Exercice : Diagnostic de Bankim]
\begin{enumerate}
    \item \textbf{Situation biogéographique :} Zone de transition \textbf{soudano-guinéenne} (Adamaoua). Climat : pluviométrie élevée (1500 mm) mais saison sèche marquée (5 mois) $\rightarrow$ savane arborée dense (\emph{Isoberlinia} = indicateur fort). Sol : ferrugineux lessivé, cuirasse ferrugineuse (pétroferrique) à profondeur variable, moins lessivé que le ferrallitique du Sud, meilleure réserve en bases mais contrainte physique (croûte) et érosion.
    \item \textbf{Conflits d'usage :}
        \begin{enumerate}
            \item \textbf{Agriculture vs Élevage (transhumance) :} Extension cultures (maïs/igname) sur parcours traditionnels + résidus de récolte non valorisés $\rightarrow$ blocage couloirs, dégâts cultures par bétail (période sèche/soudure).
            \item \textbf{Exploitation ressources vs Conservation (Parc Mbam-Djerem) :} Braconnage (viande/ivoire), feux de brousse précoces (chasse, pâturage) partant des zones villageoises vers le parc, prélèvement bois d'œuvre/énergie aux abords.
        \end{enumerate}
    \item \textbf{Mesures techniques adaptées (Sol ferrugineux + Cohabitation) :}
        \begin{enumerate}
            \item \textbf{Parcs à karité / néré / acacia (Agroforesterie) :} Protection/plantation d'arbres fourragers/légumineux (fixation N, ombre, fruits) dans les champs $\rightarrow$ amendement organique (feuilles), fourrage sec (gousses), lutte érosion (racines brisent croûte), revenu (karité).
            \item \textbf{Cultures fourragères en bande / haies vives :} \emph{Stylosanthes}, \emph{Brachiaria}, \emph{Leucaena} en bordure champs/parcours $\rightarrow$ stock fourrager saison sèche (fenaison/ensilage), coupe-pare-feu, enrichissement sol (N, M.O.), délimitation visuelle couloirs transhumance.
            \item \textbf{Aménagement points d'eau pastoraux (forages/mares) HORS des couloirs/cultures :} Réduit concentration bétail/piétinement sols sensibles, limite divagation vers champs/parc. Gestion par comité mixte éleveurs/cultivateurs.
        \end{enumerate}
    \item \textbf{Corridors écologiques :} Le parc Mbam-Djerem (4165 km$^2$) est une « île » relative. Les éléphants et grands ongulés (hippotragues, cobs, buffles) ont des aires vitales immenses (éléphant : 500--3000 km$^2$) et effectuent des migrations saisonnières (eau, sel, pâture, reproduction). Sans corridors (galeries forestières le long rivières Mbam, Djerem, affluents), les populations sont isolées $\rightarrow$ consanguinité, extinction locale, conflits Homme-Faune accrus (éléphants dans champs). Les corridors assurent la \textbf{connectivité fonctionnelle} du paysage (métapopulation).
\end{enumerate}
\end{corrige}

\coursec{Introduction : le Cameroun, « Afrique en miniature »}

Le Cameroun est souvent qualifié d'\textbf{« Afrique en miniature »} tant il condense, sur \SI{475\,442}{km^2}, la quasi-totalité des grands paysages, climats, végétations, sols et groupes humains du continent africain. Cette diversité exceptionnelle résulte de trois facteurs structurels :
\begin{enumerate}
    \item \textbf{Une position charnière} : à la jonction de l'Afrique de l'Ouest (golfe de Guinée), de l'Afrique centrale (bassin du Congo) et de l'Afrique sahélienne (bassin du lac Tchad), entre le 2\textsuperscript{e} et le 13\textsuperscript{e} parallèle nord.
    \item \textbf{Un relief structurant} : la \textbf{ligne du Cameroun} (alignement volcanique SW-NE : mont Cameroun \SI{4095}{m}, Manengouba, Bamboutos, Adamaoua, Mandara) qui barre le pays, créant des étages altitudinaux et une barrière climatique.
    \item \textbf{Un gradient climatique nord-sud} : de l'équatorial humide (sud) au sahélien sec (extrême-nord), en passant par le tropical de transition (Adamaoua) et le tempéré d'altitude (Ouest).
\end{enumerate}
Cette diversité physique fonde une mosaïque de \textbf{quatre grands ensembles biogéographiques} aux potentialités complémentaires, mais dont l'intégration nationale reste un défi permanent d'aménagement du territoire.

\coursec{Les grands ensembles biogéographiques du Cameroun}

La diversité des végétations et des sols que nous venons d'étudier ne se disperse pas au hasard : elle s'organise en grands ensembles cohérents, où le climat, le relief, la végétation, les sols et les activités humaines forment un système interdépendant. C'est ce que les géographes appellent des \cle{ensembles biogéographiques} : des vastes unités spatiales dont les composantes naturelles et humaines présentent une forte homogénéité interne et des contrastes nets avec les ensembles voisins. Le Cameroun, par sa position charnière entre l'Afrique équatoriale et l'Afrique soudano-sahélienne, et par son relief très contrasté, se divise naturellement en \textbf{quatre grands ensembles biogéographiques}. Cette structuration n'est pas seulement académique : elle fonde l'organisation de l'espace national, les flux d'échanges internes et les politiques d'aménagement du territoire.

\courssub{Découpage synthétique : les quatre grands ensembles}

Le territoire camerounais se partage en quatre grandes entités, chacune caractérisée par un couplage climat--végétation--sols--activités original :

\begin{enumerate}
    \item \textbf{L'ensemble forestier du Sud} (régions du Sud, du Littoral, du Centre, de l'Est, partie du Sud-Ouest) : domaine de la forêt dense humide, du climat équatorial, de l'agriculture de plantation (cacao, café, hévéa, palmier à huile) et vivrière (plantain, manioc, igname), de la chasse et de la pêche fluviale.
    \item \textbf{L'ensemble des hautes terres de l'Ouest} (régions de l'Ouest et du Nord-Ouest) : domaine de montagne, de climat tempéré d'altitude, de sols volcaniques très fertiles, de forte densité de population et d'agriculture intensive sur terrasses.
    \item \textbf{L'ensemble des savanes de l'Adamaoua} (région de l'Adamaoua) : domaine de plateaux herbeux, de transition climatique, d'élevage bovin dominant et de très faible densité humaine.
    \item \textbf{L'ensemble soudano-sahélien du Grand Nord} (régions du Nord et de l'Extrême-Nord) : domaine de savanes arbustives puis de steppes, de climat tropical sec (soudanien au sud, sahélien au nord), d'agriculture pluviale (mil, sorgho, maïs, coton), d'élevage extensif et de pêche dans le Logone et le lac Tchad.
\end{enumerate}

\begin{aretenir}[Les quatre ensembles biogéographiques du Cameroun]
Le Cameroun se structure en quatre grands ensembles biogéographiques aux ressources complémentaires :
\begin{itemize}
    \item \textbf{Ensemble forestier du Sud} : forêt dense, climat équatorial, cultures d'exportation et vivrières, forte biodiversité.
    \item \textbf{Hautes terres de l'Ouest} : relief accidenté, climat frais, sols volcaniques, très forte densité, agriculture intensive.
    \item \textbf{Savanes de l'Adamaoua} : plateaux, transition, élevage bovin, très faible densité.
    \item \textbf{Grand Nord soudano-sahélien} : savanes/steppes, climat sec, mil/sorgho/coton, élevage, pêche, densité contrastée.
\end{itemize}
Ces ensembles ne sont pas des blocs figés : leurs limites sont des \emph{zones de transition progressives} (écotones).
\end{aretenir}

\begin{attention}[Piège fréquent : des ensembles figés]
Ne présentez jamais ces ensembles comme des cases hermétiques aux frontières nettes. Dans la réalité, les limites sont des \textbf{zones de transition (écotones)} larges de plusieurs dizaines de kilomètres : la lisière forêt-savane fluctue selon les incendies, le climat et l'action humaine ; le piémont des hautes terres s'étage progressivement vers l'Adamaoua ; le Nord passe graduellement du soudanien au sahélien. Sur un croquis, représentez ces transitions par des zones hachurées ou des traits discontinus, jamais par un trait unique brutal.
\end{attention}

\courssub{L'ensemble forestier du Sud : le domaine de l'humide et de la forêt}

\courssub{Climat, végétation et sols}

Cet ensemble couvre près de \SI{45}{\%} du territoire national. Il bénéficie d'un \textbf{climat équatorial de type guinéen} (au Sud, au Centre, à l'Est) et \textbf{équatorial de type camerounais} (au Littoral et au Sud-Ouest, avec une courte saison sèche). Les précipitations annuelles dépassent \SI{1500}{mm} partout, atteignant \SI{3000}{mm} à \SI{4000}{mm} sur le piémont du Cameroun (Debundscha : \SI{10000}{mm/an}, record africain). Les températures sont élevées et stables (\SI{24}{\celsius}--\SI{26}{\celsius} de moyenne annuelle), l'hygrométrie saturante.

La végétation climax est la \textbf{forêt dense humide sempervirente} (persistante) au Sud et à l'Est, et \textbf{semi-décidue} au Centre et à l'Est (perte partielle des feuilles en saison sèche). C'est une forêt stratifiée (strates émergente, principale, intermédiaire, herbacée), à haute biodiversité (essences : ayous, sapelli, iroko, moabi, azobé, etc.). Les sols sont des \textbf{ferralsols} (sols ferralitiques) : profonds, très lessivés, acides, pauvres en bases échangeables, mais à bonne réserve en eau et structure stable. Leur fertilité chimique est faible ; la fertilité repose sur le cycle rapide de la matière organique (humus de surface).

\begin{exemplebox}[Le cycle forestier et la fertilité]
En forêt dense, 80 à 90\% des nutriments sont stockés dans la biomasse aérienne (bois, feuilles) et non dans le sol. La décomposition ultra-rapide (chaleur + humidité) recycle la litière en quelques semaines. Toute coupe rase sans remise en culture immédiate ou sans jachère longue entraîne une perte irréversible de fertilité : le sol se « ferralitise » davantage, durcit (formation de cuirasse latéritique) et devient impropre à l'agriculture sans apports massifs.
\end{exemplebox}

\courssub{Activités humaines et peuplement}

L'agriculture domine : \textbf{plantes pérennes d'exportation} (cacao au Centre, Sud, Sud-Ouest ; café arabica à l'Ouest du Sud et au Littoral ; hévéa et palmier à huile au Littoral, Sud-Ouest, Sud) et \textbf{vivriers} (plantain, manioc, igname, taro, maïs en saison sèche). L'agroforesterie traditionnelle (champs sous couvert forestier) imite la structure de la forêt. La \textbf{chasse} (faune de brousse) et la \textbf{pêche artisanale} (rivières Nyong, Sanaga, Ntem, lacs de barrage) complètent l'alimentation protéinée. L'exploitation forestière industrielle (bois d'œuvre) et artisanale génère des revenus mais pose la question de la gestion durable.

La densité de population est \textbf{faible à moyenne} (\SI{10}{\text{à}} \SI{50}{hab/km^2} en moyenne), concentrée le long des axes routiers (Yaoundé--Bertoua, Douala--Yaoundé--Ngaoundéré) et dans les villes (Yaoundé, Douala, Ebolowa, Bertoua, Sangmélima). L'Est et le Sud profond restent des « vides » relatifs.

\courssub{Les hautes terres de l'Ouest : le jardin du Cameroun}

\courssub{Relief, climat et sols : un milieu de montagne tropicale}

Cet ensemble (régions de l'Ouest et du Nord-Ouest) correspond au \textbf{massif volcanique du Cameroun} (ligne du Cameroun) : monts Bamboutos (\SI{2740}{m}), mont Manengouba (\SI{2411}{m}), monts Bamiléké, plateau de Bamenda. Le relief est \textbf{très accidenté} : pentes fortes, vallées encaissées, cônes volcaniques, plaines de lave. L'altitude (souvent \SI{1000}{\text{à}} \SI{2500}{m}) crée un \textbf{climat tempéré d'altitude} : températures moyennes \SI{18}{\celsius}--\SI{22}{\celsius}, amplitudes thermiques journalières marquées, gelées possibles au-dessus de \SI{2000}{m}. Les précipitations sont abondantes (\SI{1800}{\text{à}} \SI{3000}{mm/an}) avec une saison sèche courte (novembre--mars).

Les sols sont des \textbf{andosols} (sols volcaniques jeunes) et des \textbf{nitisols} (sols rouges profonds sur roches basaltiques altérées) : riches en matière organique, bien structurés, à forte capacité d'échange cationique, bien drainés mais retenant l'eau. Ce sont les \textbf{sols les plus fertiles du pays}.

\courssub{Agriculture intensive et densité record}

La combinaison sols fertiles + climat favorable + forte main-d'œuvre a permis le développement d'une \textbf{agriculture intensive sur terrasses} (cultures en courbes de niveau, haies vives, fumure organique, association cultures--élevage). Productions : \textbf{vivriers} (maïs, haricot, pomme de terre, taro, igname, banane plantain, chou, carotte, poireau) et \textbf{perennes} (café arabica, thé à Ndu et Tole, arbres fruitiers : avocatier, agrumes, prunier d'Afrique \emph{Dacryodes edulis}). L'élevage (bovins, porcs, volailles, lapins) est en stabulation ou pâturage tournant, intégré aux cultures (fumure).

\begin{exemplebox}[Densité de population record]
Les hautes terres de l'Ouest atteignent des densités de \textbf{300 à plus de 500 habitants/km²} (départements du Ndé, de la Menoua, du Haut-Nkam, du Mezam), parmi les plus élevées d'Afrique rurale. Cette surpopulation relative pousse à l'exode vers les villes (Bafoussam, Bamenda, Dschang) et vers les fronts pionniers du Sud (Est, Sud) et du Nord (Bénoué). La pression foncière y est extrême : taille moyenne des exploitations < \SI{1}{ha}, morcellement extrême.
\end{exemplebox}

\begin{aretenir}[Contraste démographique majeur]
\begin{itemize}
    \item \textbf{Hautes terres de l'Ouest} : \textbf{> 300 hab/km²} (pointes > 500 hab/km²).
    \item \textbf{Adamaoua} : \textbf{< 15 hab/km²} (parfois < 5 hab/km² dans le Nord de la région).
    \item \textbf{Grand Nord} : contrasté (basse vallée du Logone > 100 hab/km² ; Diamaré, Mayo-Tsanaga 50--80 hab/km² ; zones pastorales < 10 hab/km²).
\end{itemize}
Ce gradient démographique structure les migrations internes (Sud-Nord, Ouest-Est, rurales-urbaines).
\end{aretenir}

\courssub{L'ensemble des savanes de l'Adamaoua : la transition pastorale}

\courssub{Un plateau de transition}

L'Adamaoua (région de l'Adamaoua, chef-lieu Ngaoundéré) est un \textbf{vaste plateau} (\SI{1000}{\text{à}} \SI{1500}{m d'altitude}), large de \SI{200}{km} est-ouest, charnière entre le Sud forestier et le Nord soudano-sahélien. Le climat est \textbf{tropical de transition (soudano-guinéen)} : une saison sèche marquée (novembre--mars, harmattan), une saison des pluies unique (avril--octobre), précipitations \SI{1300}{\text{à}} \SI{1600}{mm/an}, températures douces (\SI{22}{\celsius} moyenne annuelle). La végétation est une \textbf{savane arborée à \emph{Isoberlinia doka}, \emph{Terminalia}, \emph{Anogeissus}} (savane claire), avec des îlots de forêt galerie le long des cours d'eau (Mbéré, Vina, Lom, Pangar). Les sols sont des \textbf{ferralsols lessivés} (sols ferrugineux tropicaux lessivés), moins profonds, plus sableux, à fertilité moyenne, sensibles à l'érosion.

\courssub{L'élevage bovin, activité dominante}

L'Adamaoua est le \textbf{premier bassin d'élevage bovin du Cameroun} (race \emph{Goudali}, \emph{Red Fulani}, \emph{White Fulani}, croisements). L'élevage est \textbf{extensif, transhumant} : migration saisonnière des troupeaux vers le Nord (Bénoué) en saison sèche (pâturages verts, points d'eau) et retour sur le plateau en saison des pluies (éviter les mouches tsé-tsé, valoriser l'herbe nouvelle). L'agriculture vivrière (mil, sorgho, maïs, manioc, igname, arachide) est pratiquée en marge des villages, souvent en système itinérant sur brûlis. La densité de population est \textbf{très faible} (\SI{5}{\text{à}} \SI{15}{hab/km^2}), l'espace est perçu comme « vide » mais il est en réalité géré par les chefferies peules (lamidats) selon un code pastoral coutumier.

\begin{propriete}[La complémentarité des ensembles biogéographiques]
La spécialisation régionale crée une \textbf{interdépendance économique nationale} :
\begin{itemize}
    \item \textbf{Grand Nord + Adamaoua} $\rightarrow$ \textbf{viande bovine} (bœufs sur pied) vers les villes du Sud et de l'Ouest.
    \item \textbf{Ensemble forestier du Sud} $\rightarrow$ \textbf{café, cacao, bois, hévéa, palmier à huile, banane plantain} vers le Nord et l'Ouest.
    \item \textbf{Hautes terres de l'Ouest} $\rightarrow$ \textbf{vivriers frais} (maïs, haricot, pomme de terre, légumes, fruits) vers les villes du Sud (Yaoundé, Douala) et du Nord (Garoua, Maroua, Ngaoundéré).
    \item \textbf{Grand Nord} $\rightarrow$ \textbf{coton, oignon, riz, poisson (Logone, lac Tchad)} vers le Sud.
\end{itemize}
Ces flux empruntent l'axe routier et ferroviaire \textbf{Douala--Yaoundé--Ngaoundéré--Garoua--Maroua}, colonne vertébrale du territoire.
\end{propriete}

\courssub{L'ensemble soudano-sahélien du Grand Nord : le domaine de la saison sèche}

\courssub{Du soudanien au sahélien : un gradient climatique fort}

Le Grand Nord (régions du Nord et de l'Extrême-Nord) s'étend du 7\textsuperscript{e} au 13\textsuperscript{e} parallèle nord. Il présente un \textbf{gradient climatique nord-sud} net :
\begin{itemize}
    \item \textbf{Zone soudanaire (sud du Nord, Bénoué)} : \SI{900}{\text{à}} \SI{1200}{mm/an}, saison des pluies 5--6 mois (mai--octobre), savane arborée (\emph{Isoberlinia}, \emph{Terminalia}, \emph{Combretum}), sols ferrugineux tropicaux lessivés.
    \item \textbf{Zone soudano-sahélienne (nord du Nord, Diamaré, Mayo-Tsanaga)} : \SI{600}{\text{à}} \SI{900}{mm/an}, saison des pluies 3--4 mois (juin--septembre), savane arbustive épineuse (\emph{Acacia}, \emph{Balanites}, \emph{Commiphora}), sols ferrugineux peu lessivés, sols hydromorphes en bas-fonds.
    \item \textbf{Zone sahélienne (Extrême-Nord, Logone, lac Tchad)} : \SI{400}{\text{à}} \SI{600}{mm/an}, saison des pluies 2--3 mois (juillet--septembre), steppe à \emph{Acacia raddiana}, \emph{Leptadenia}, \emph{Panicum}, sols sableux (arénosols), sols salins (solonchaks) près du lac.
\end{itemize}
L'irrégularité interannuelle des pluies (coefficient de variation 20--30\%) est le risque majeur : sécheresses récurrentes (1973, 1984, 2010, 2021).

\courssub{Agriculture pluviale, élevage et pêche : des systèmes adaptés}

L'agriculture est \textbf{pluviale, extensive, à jachère courte ou inexistante} (pression foncière). Cultures dominantes : \textbf{mil (pénicillaire), sorgho (sorgho rouge, sorgho blanc), maïs} (en extension), \textbf{coton} (culture de rente majeure, encadrée par la SODECOTON, zone Nord), \textbf{arachide, niébé, sésame, riz} (bas-fonds, plaine de la Bénoué, Yaérés). L'élevage (bovins, ovins, caprins, ânes, chevaux) est \textbf{extensif, transhumant Nord-Sud} (vers l'Adamaoua et la Bénoué en saison sèche). La \textbf{pêche} est vitale dans le \textbf{Logone}, la \textbf{Bénoué} et le \textbf{lac Tchad} (poissons : capitaine, tilapia, silure, carpe) ; elle assure protéines et revenus (femmes transformatrices : séchage, fumage, salage).

La densité de population est \textbf{contrastée} : forte dans la \textbf{basse vallée du Logone et les Yaérés} (\SI{100}{\text{à}} \SI{200}{hab/km^2}, cultures de décrue, pêche, coton), moyenne dans le \textbf{Diamaré et le Mayo-Tsanaga} (\SI{50}{\text{à}} \SI{80}{hab/km^2}), faible dans les \textbf{zones pastorales pures} du Nord-Nord (\SI{<10}{hab/km^2}). Les villes (Garoua, Maroua, Ngaoundéré, Kousseri) concentrent la croissance urbaine.

\begin{exoresolu}[Analyse de la complémentarité Nord-Sud à partir de données]
\textbf{Énoncé :} À partir des données ci-dessous, montrez que les ensembles biogéographiques du Cameroun sont complémentaires.

\begin{center}
\begin{tabular}{|l|c|c|c|}\hline
\textbf{Produit} & \textbf{Zone de production principale} & \textbf{Zone de consommation principale} & \textbf{Flux} \\ \hline
Bœufs sur pied & Adamaoua, Grand Nord & Yaoundé, Douala, Bafoussam & Nord $\rightarrow$ Sud \\ \hline
Cacao, café, bois & Sud, Est, Centre, Littoral & Nord, Extrême-Nord, Ouest & Sud $\rightarrow$ Nord \\ \hline
Maïs, haricot, pomme de terre, légumes & Ouest, Nord-Ouest & Yaoundé, Douala, Garoua, Maroua & Ouest $\rightarrow$ Sud/Nord \\ \hline
Coton, oignon, riz, poisson & Grand Nord & Usines textile (Garoua, Douala), marchés Sud & Nord $\rightarrow$ Sud \\ \hline
\end{tabular}
\end{center}

\textbf{Solution :}
Le tableau met en évidence une \textbf{spatialisation sectorielle} nette : chaque ensemble se spécialise dans les productions pour lesquelles il a un avantage comparatif naturel (climat, sols, eau). Les flux vont \textbf{des zones de production excédentaire vers les zones de consommation déficitaire} (villes, zones non productrices). L'axe Douala--Ngaoundéré--Garoua--Maroua structure ces échanges. Cette complémentarité est la base de l'\textbf{unité économique nationale} : le Nord nourrit le Sud en protéines animales et en coton ; le Sud fournit au Nord les produits d'exportation (cacao, café, bois) et les produits manufacturés ; l'Ouest approvisionne tout le pays en vivriers frais. Toute perturbation (conflit, route coupée, sécheresse) déséquilibre l'ensemble du système.
\end{exoresolu}

\courssub{Tableau comparatif des quatre ensembles biogéographiques}

\begin{popfigure}
\legende{Tableau comparatif des quatre grands ensembles biogéographiques du Cameroun (source : synthèse MINESEC, INS, IRD).}
\begin{center}
\small
\begin{tabularx}{\linewidth}{|l|X|X|X|X|}\hline
\rowcolor{popblueL}
\textbf{Critère} & \textbf{Ensemble forestier du Sud} & \textbf{Hautes terres de l'Ouest} & \textbf{Savanes de l'Adamaoua} & \textbf{Grand Nord soudano-sahélien} \\ \hline
\textbf{Régions concernées} & Sud, Littoral, Centre, Est, partie Sud-Ouest & Ouest, Nord-Ouest & Adamaoua & Nord, Extrême-Nord \\ \hline
\textbf{Relief} & Plaine côtière, plateau sud-camerounais (\SI{200}{\text{à}} \SI{600}{m}), chaos granitiques & Massif volcanique, monts (\SI{1000}{\text{à}} \SI{2740}{m}), vallées encaissées & Grand plateau (\SI{1000}{\text{à}} \SI{1500}{m}), doux ondulé & Plaine de la Bénoué, piémont mandara, plaine du Logone, bassin du lac Tchad \\ \hline
\textbf{Climat} & Équatorial guinéen (S/C/E) \& camerounais (L/SO) : \SI{1500}{\text{à}} \SI{4000}{mm}, 2--3 mois secs & Tempéré d'altitude : \SI{1800}{\text{à}} \SI{3000}{mm}, frais (\SI{18}{\text{à}}\SI{22}{\celsius}), 3--4 mois secs & Tropical de transition (soudano-guinéen) : \SI{1300}{\text{à}} \SI{1600}{mm}, 4--5 mois secs & Gradient soudano-sahélien : \SI{400}{\text{à}} \SI{1200}{mm}, 2--6 mois secs, harmattan \\ \hline
\textbf{Végétation} & Forêt dense humide sempervirente \& semi-décidue & Forêt de montagne, savane d'altitude, cultures & Savane claire à \emph{Isoberlinia}, forêts galeries & Savane arborée $\rightarrow$ arbustive $\rightarrow$ steppe (N) \\ \hline
\textbf{Sols} & Ferralsols (ferralitiques) : profonds, acides, pauvres en bases, humus de surface & Andosols, Nitisols : volcaniques, riches, structurés, fertiles & Ferralsols lessivés (ferrugineux tropicaux lessivés), sableux, érodables & Ferrugineux peu lessivés, arénosols, vertisols (bas-fonds), solonchaks (lac) \\ \hline
\textbf{Activités dominantes} & Agriculture de plantation (cacao, café, hévéa, palmier), vivrière (plantain, manioc), chasse, pêche, exploitation forestière & Agriculture intensive sur terrasses (maïs, haricot, pomme de terre, légumes, café arabica, thé), élevage intégré & Élevage bovin extensif transhumant (race Goudali), agriculture vivrière itinérante (mil, maïs) & Agriculture pluviale (mil, sorgho, maïs, coton, riz), élevage extensif, pêche (Logone, lac Tchad) \\ \hline
\textbf{Densité de population} & Faible à moyenne (\SI{10}{\text{à}} \SI{50}{hab/km^2}), concentrée sur axes & \textbf{Très forte} (\SI{300}{\text{à}} \SI{>500}{hab/km^2}) & \textbf{Très faible} (\SI{5}{\text{à}} \SI{15}{hab/km^2}) & Contrastée : forte (Logone, Yaérés \SI{>100}{hab/km^2}), faible (zones pastorales \SI{<10}{hab/km^2}) \\ \hline
\textbf{Villes principales} & Yaoundé, Douala, Ebolowa, Bertoua, Kribi, Sangmélima & Bafoussam, Bamenda, Dschang, Bafang, Kumbo & Ngaoundéré, Tibati, Meiganga, Tignère & Garoua, Maroua, Kousseri, Mokolo, Yagoua, Guider \\ \hline
\end{tabularx}
\end{center}
\end{popfigure}

\courssub{Les grands groupes humains du Cameroun}

La répartition des populations épouse largement le découpage biogéographique, témoignant de l'adaptation historique des sociétés à leur milieu. On distingue quatre grands ensembles linguistiques et culturels :

\begin{enumerate}
    \item \textbf{Les populations bantoues (Sud, Centre, Est, Littoral, Sud-Ouest)} : Beti-Fang (Bulu, Fang, Eton, Ewondo), Bassa, Douala, Bakoko, Makaa, Kako, Konabembe. Agricultures de forêt (plantain, manioc, cacao), chefferies traditionnelles, organisation lignagère.
    \item \textbf{Les populations semi-bantoues (Ouest, Nord-Ouest)} : Bamiléké (Ghomala', Fe'fe', Medumba, Yemba), Bamoun (Shüpamum), Tikar, groupes des Grassfields (Bamenda, Kom, Nso, Oku). Agriculture intensive sur terrasses, chefferies complexes (fonctions royales, sociétés secrètes), forte densité, dynamisme entrepreneurial et migratoire.
    \item \textbf{Les populations soudano-sahéliennes (Nord, Extrême-Nord, partie Adamaoua)} : Peuls (Fulfulde) majoritaires (éleveurs, islamisés), Massa, Moundang, Toupouri, Mousgoum, Kotoko, Arabes Choa. Systèmes agropastoraux, chefferies (lamidats), islam prédominant, habitat groupé en villages de concession.
    \item \textbf{Les populations de l'Adamaoua (transition)} : Peuls (Fulbe) dominants (lamidats de Ngaoundéré, Tibati, Rey Bouba), Mboum, Dii, Nyem-Nyem, Gbaya. Zone de brassage culturel et linguistique (fulfulde langue véhiculaire), islam et christianisme cohabitent.
\end{enumerate}

\begin{aretenir}[Répartition des grands groupes humains]
\begin{itemize}
    \item \textbf{Bantous} : Sud forestier (agriculteurs forestiers, chefferies lignagères).
    \item \textbf{Semi-Bantous (Grassfields)} : Hautes terres de l'Ouest (agriculteurs intensifs, chefferies royales, forte densité).
    \item \textbf{Soudano-Sahéliens} : Grand Nord (agropasteurs, lamidats, islam).
    \item \textbf{Adamaoua} : Zone de contact Peuls / groupes voltaïques (Mboum, Dii, Gbaya), transition culturelle.
\end{itemize}
La langue officielle (français, anglais) et les langues véhiculaires (fulfulde au Nord/Adamaoua, pidgin au Sud-Ouest, ewondo/bassa/douala au Centre/Sud/Littoral) assurent la communication interethnique.
\end{aretenir}

\courssub{La complémentarité économique : le ciment de l'unité nationale}

Les quatre ensembles ne vivent pas en autarcie. Leurs spécialisations créent des \textbf{flux d'échanges intérieurs massifs} qui structurent l'espace camerounais :

\begin{itemize}
    \item \textbf{Flux Nord $\rightarrow$ Sud} : bœufs sur pied (Adamaoua, Grand Nord $\rightarrow$ abattoirs de Yaoundé, Douala, Bafoussam), coton (Grand Nord $\rightarrow$ usines CICAM à Garoua, puis export via Douala), oignons, poisson séché/fumé (Logone, lac Tchad $\rightarrow$ marchés du Sud).
    \item \textbf{Flux Sud $\rightarrow$ Nord} : café, cacao, bois, hévéa, huile de palme, banane plantain (Sud $\rightarrow$ marchés du Nord), produits manufacturés importés (Douala $\rightarrow$ intérieur).
    \item \textbf{Flux Ouest $\rightarrow$ Sud/Nord} : vivriers frais (maïs, haricots, pommes de terre, légumes, fruits $\rightarrow$ Yaoundé, Douala, Garoua, Maroua, Ngaoundéré).
    \item \textbf{Flux Est/Ouest} : main-d'œuvre agricole (Ouest $\rightarrow$ plantations Sud/Est), migration pendulaire, transfert de compétences.
\end{itemize}

Ces flux empruntent l'\textbf{axe structurant Douala--Yaoundé--Ngaoundéré--Garoua--Maroua} (route nationale N1, chemin de fer Camrail, oléoduc Tchad--Cameroun, fibre optique). La \textbf{politique d'aménagement du territoire} (décentralisation, routes régionales, pôles de croissance régionaux) vise à renforcer cette complémentarité et à réduire les déséquilibres (enclavement de l'Est, de l'Adamaoua, du Grand Nord).

\begin{methode}[Construire un croquis géographique des ensembles biogéographiques]
Pour représenter les quatre ensembles sur un croquis de synthèse :
\begin{enumerate}
    \item \textbf{Tracer le fond de carte} : contour simplifié du Cameroun (polygone approximatif), frontières internationales, capitale Yaoundé, grands axes (N1, chemin de fer), fleuves majeurs (Sanaga, Bénoué, Logone, Nyong), lac Tchad.
    \item \textbf{Délimiter les 4 ensembles} : traits discontinus ou zones hachurées pour marquer les transitions (écotones). Ne pas tracer de traits nets.
    \item \textbf{Choisir des figurés de surface (coloris)} : 
    \begin{itemize}
        \item Vert foncé = Ensemble forestier du Sud
        \item Vert clair / marron = Hautes terres de l'Ouest
        \item Jaune/ocre clair = Savanes de l'Adamaoua
        \item Orange/rouge clair = Grand Nord soudano-sahélien
    \end{itemize}
    \item \textbf{Ajouter des figurés ponctuels/linéaires} : 
    \begin{itemize}
        \item Villes principales (carrés proportionnels)
        \item Flèches de flux (épaisseur $\propto$ volume) : bœufs N$\rightarrow$S, vivriers O$\rightarrow$S/N, cacao/café S$\rightarrow$N
        \item Ressources minières/pétrolières si pertinent
    \end{itemize}
    \item \textbf{Compléter les éléments obligatoires} :
    \begin{itemize}
        \item \textbf{Titre} en haut, centré, précis : « Le Cameroun : les quatre grands ensembles biogéographiques et leurs complémentarités »
        \item \textbf{Légende} en bas à gauche, organisée par type de figurés (surface, ponctuels, linéaires)
        \item \textbf{Échelle} graphique (ex. 1 cm = 100 km)
        \item \textbf{Orientation} : flèche Nord en haut à droite
        \item \textbf{Source} : données INS, MINEPAT, cartes IGN
    \end{itemize}
\end{enumerate}
\end{methode}

\begin{popfigure}
\legende{Croquis de synthèse : les quatre grands ensembles biogéographiques du Cameroun et leurs flux de complémentarité.}
\begin{center}
\begin{tikzpicture}[scale=0.85]
% Contour très simplifié du Cameroun (polygone approximatif)
\draw[popdark, thick] (0,0) -- (10,0) -- (11,2) -- (11,8) -- (9,10) -- (6,11) -- (3,10) -- (1,8) -- (0,6) -- cycle;

% Frontières (simplifiées)
\draw[popmuted, dashed] (0,0) -- (0,6); % Ouest
\draw[popmuted, dashed] (10,0) -- (11,2); % Sud-Ouest (golfe)
\draw[popmuted, dashed] (11,8) -- (9,10); % Nord-Ouest (Nigeria)
\draw[popmuted, dashed] (6,11) -- (3,10); % Nord (Tchad)
\draw[popmuted, dashed] (1,8) -- (0,6); % Est (RCA, Congo)

% Zones biogéographiques (remplissage)
% 1. Ensemble forestier du Sud (Sud, Littoral, Centre, Est, partie SO)
\fill[popgreen!25] (0,0) -- (10,0) -- (8,3) -- (5,4) -- (3,4) -- (1,3) -- (0,2) -- cycle;
% 2. Hautes terres de l'Ouest (Ouest, NO) - zone centrale ouest
\fill[popgreen!45!poporange!70!popdark!30] (4,2) -- (7,2) -- (8,4) -- (6,5) -- (4,4) -- cycle;
% 3. Adamaoua (plateau central)
\fill[popgold!30] (3,5) -- (8,5) -- (9,7) -- (6,8) -- (4,6) -- cycle;
% 4. Grand Nord (Nord, Extrême-Nord)
\fill[poporange!25] (4,6) -- (9,7) -- (9,9) -- (6,10) -- (4,7) -- cycle;

% Villes principales (points)
\foreach \x/\y/\name/\size in {
    5/1.5/Douala/3, 4.5/3/Yaoundé/4, 7/3.5/Bafoussam/2.5, 5.5/4.5/Bamenda/2.5,
    6/6.5/Ngaoundéré/3, 7.5/8.5/Garoua/3, 8.5/9.5/Maroua/3, 9/9/Kousseri/2,
    2/4/Bertoua/2, 1.5/1.5/Ebolowa/2, 2/1/Kribi/2, 8/4.5/Meiganga/2
} {
    \fill[popdark] (\x,\y) circle (\size pt);
    \node[above right, font=\tiny, popdark] at (\x,\y) {\name};
}

% Fleuves
\draw[popblue, thick] (4.5,3) .. controls (5,4) and (4,5) .. (3,5.5) node[left, font=\tiny, popblue] {Sanaga};
\draw[popblue, thick] (6.5,6.5) .. controls (7,7.5) and (8,8.5) .. (8.5,9) node[right, font=\tiny, popblue] {Bénoué};
\draw[popblue, thick] (8,8.5) .. controls (8.5,9) and (9,9.5) .. (9,9.5) node[right, font=\tiny, popblue] {Logone};
\draw[popblue, thick] (2.5,4.5) .. controls (2,5) and (1.5,5.5) .. (1,6) node[left, font=\tiny, popblue] {Nyong};

% Lac Tchad (coin NE)
\fill[popblue!30] (9,9.5) ellipse (0.8 and 0.4);
\node[font=\tiny, popblue] at (9,9.5) {Lac Tchad};

% Axes routiers/ferroviaires
\draw[popdark, thick, decorate, decoration={snake, amplitude=1pt, segment length=4pt}] (5,1.5) -- (4.5,3) -- (6,6.5) -- (7.5,8.5) -- (8.5,9.5);
\node[font=\tiny, popdark, rotate=30] at (5.5,5) {N1 / Camrail};

% Flèches de flux (complémentarité)
% Bœufs Nord -> Sud
\draw[popred, thick, ->] (7,8) -- (5.5,4) node[midway, above left, font=\tiny, popred] {Bœufs};
% Vivriers Ouest -> Sud/Nord
\draw[popgreen!70!popdark, thick, ->] (5.5,4) -- (4.5,3) node[midway, above, font=\tiny, popgreen!70!popdark] {Vivriers};
\draw[popgreen!70!popdark, thick, ->] (5.5,4) -- (6.5,6) node[midway, right, font=\tiny, popgreen!70!popdark] {Vivriers};
% Cacao/Café Sud -> Nord
\draw[poporange!70!popdark, thick, ->] (4,3) -- (6,6) node[midway, above right, font=\tiny, poporange!70!popdark] {Cacao/Café};
% Coton/Poisson Nord -> Sud
\draw[poporange, thick, ->] (8,9) -- (6,5) node[midway, right, font=\tiny, poporange] {Coton/Poisson};

% Légende (simplifiée, dans le croquis)
\begin{scope}[shift={(0.5,0.5)}]
\node[font=\small\bfseries, popdark, anchor=south west] at (0,0) {Légende :};
\fill[popgreen!25] (0,-0.5) rectangle (0.5,-0.9); \node[right, font=\tiny] at (0.6,-0.7) {Forêt du Sud};
\fill[popgreen!45!poporange!70!popdark!30] (0,-1.1) rectangle (0.5,-1.5); \node[right, font=\tiny] at (0.6,-1.3) {Hautes terres Ouest};
\fill[popgold!30] (0,-1.7) rectangle (0.5,-2.1); \node[right, font=\tiny] at (0.6,-1.9) {Savanes Adamaoua};
\fill[poporange!25] (0,-2.3) rectangle (0.5,-2.7); \node[right, font=\tiny] at (0.6,-2.5) {Grand Nord};
\draw[popblue, thick] (0,-3.1) -- (0.5,-3.1); \node[right, font=\tiny] at (0.6,-3.1) {Fleuves};
\draw[popdark, thick, decorate, decoration={snake, amplitude=1pt, segment length=4pt}] (0,-3.5) -- (0.5,-3.5); \node[right, font=\tiny] at (0.6,-3.5) {Axe N1/Camrail};
\draw[popred, thick, ->] (0,-3.9) -- (0.5,-3.9); \node[right, font=\tiny] at (0.6,-3.9) {Flux bœufs};
\draw[popgreen!70!popdark, thick, ->] (0,-4.3) -- (0.5,-4.3); \node[right, font=\tiny] at (0.6,-4.3) {Flux vivriers};
\draw[poporange!70!popdark, thick, ->] (0,-4.7) -- (0.5,-4.7); \node[right, font=\tiny] at (0.6,-4.7) {Flux cacao/café};
\draw[poporange, thick, ->] (0,-5.1) -- (0.5,-5.1); \node[right, font=\tiny] at (0.6,-5.1) {Flux coton/poisson};
\end{scope}

% Échelle et Nord
\node[font=\tiny, popdark] at (10.5,0.5) {\textbf{Échelle : 1 cm $\approx$ 100 km}};
\draw[popdark, ->, thick] (10.5,1.2) -- (10.5,1.8) node[above, font=\tiny] {N};
\end{tikzpicture}
\end{center}
\end{popfigure}

\courssub{Exercice d'application}

\begin{exercice}[Analyse d'un ensemble biogéographique]
Choisissez l'un des quatre ensembles biogéographiques du Cameroun. À partir du tableau comparatif et du croquis, rédigez une fiche structurée (1 page max) présentant :
\begin{enumerate}
    \item La localisation et les régions concernées.
    \item Les caractéristiques physiques (relief, climat, végétation, sols).
    \item Les activités économiques dominantes et leur adaptation au milieu.
    \item La densité et la répartition de la population.
    \item Les flux d'échanges avec les autres ensembles (ce qu'il « donne » et ce qu'il « reçoit »).
    \item Un défi majeur pour son développement durable (ex. : déforestation, érosion, désertification, pression foncière, enclavement).
\end{enumerate}
\end{exercice}

\begin{corrige}[Exercice : Fiche ensemble forestier du Sud (exemple)]
\begin{enumerate}
    \item \textbf{Localisation} : Sud du 4\textsuperscript{e} parallèle, régions Sud, Littoral, Centre, Est, partie Sud-Ouest. Frontières : Guinée équatoriale, Gabon, Congo, RCA, mer (golfe de Guinée).
    \item \textbf{Physique} : Plateau sud-camerounais (200--600 m), chaos granitiques. Climat équatorial guinéen (S/C/E : 1500--2500 mm, 2--3 mois secs) et camerounais (L/SO : 3000--4000 mm, 1--2 mois secs). Forêt dense humide sempervirente (S/E) et semi-décidue (C). Sols : ferralsols profonds, acides, humifères en surface.
    \item \textbf{Activités} : Plantations pérennes d'exportation (cacao, café robusta, hévéa, palmier à huile) + vivriers (plantain, manioc, igname, taro) en agroforesterie. Chasse, pêche fluviale. Exploitation forestière (bois d'œuvre). Adaptation : cultures d'ombrage, jachères améliorées, cueillette.
    \item \textbf{Population} : Densité faible à moyenne (10--50 hab/km²). Concentrée sur axes routiers/ferroviaires et villes (Yaoundé, Douala, Ebolowa, Bertoua, Kribi). Est et Sud profond = faibles densités (< 5 hab/km²). Ethnies : Bantu (Beti, Bassa, Boulou, Fang, Bakoko, Bassa, Makaa, Kako, etc.).
    \item \textbf{Flux} : \textbf{Sortants} : cacao, café, bois, caoutchouc, huile de palme, banane plantain $\rightarrow$ Nord, Ouest, export (Douala/Kribi). \textbf{Entrants} : bœufs (Adamaoua/Nord), vivriers frais (Ouest), coton/oignon/poisson (Nord), produits manufacturés/pétroliers (Douala).
    \item \textbf{Défi majeur} : \textbf{Déforestation et dégradation forestière} (exploitation illégale, agriculture itinérante sur brûlis, expansion plantations agro-industrielles, mines, infrastructures). Perte de biodiversité, érosion des sols, perturbation du cycle de l'eau, émissions CO₂. Réponses : gestion durable (UFA, certification FSC), aires protégées (Dja, Lobéké, Boumba Bek, Nki), reboisement, agroforesterie, REDD+.
\end{enumerate}
\end{corrige}

\begin{savaistu}[Le concept d'écotone et la lisière forêt-savane]
La limite entre l'ensemble forestier du Sud et les savanes de l'Adamaoua/Grand Nord n'est pas une ligne : c'est un \textbf{écotone} large de 50 à 100 km, zone de contact où s'affrontent et s'hybrident deux biomes. On y observe une \textbf{mosaïque forêt-savane} : îlots forestiers (galeries, îlots relictuels sur sols profonds) dans une matrice de savane, ou inversement. Cette lisière est \textbf{dynamique} : elle recule (progression de la savane) sous l'effet combiné du changement climatique (sécheresses des années 1970-80), des feux de brousse anthropiques (chasse, pâturage, nettoyage champs) et de la démographie. Inversement, l'abandon des cultures (exode rural) ou la protection (parcs) permet la \textbf{recolonisation forestière} (succession secondaire : herbes $\rightarrow$ arbustes $\rightarrow$ forêt). Étudier cette lisière (télédétection, inventaires terrain) est un indicateur clé du changement global au Cameroun.
\end{savaistu}

\courssub{Bilan : une géographie de la complémentarité}

Les quatre ensembles biogéographiques du Cameroun ne sont pas de simples cases sur une carte. Ce sont des \textbf{systèmes géographiques vivants}, façonnés par l'histoire géologique (ligne du Cameroun), climatique (gradient équateur--Sahara) et humaine (migrations, techniques agricoles, politiques). Leur \textbf{complémentarité des ressources} (viande du Nord, bois et pérennes du Sud, vivriers de l'Ouest, coton/poisson du Grand Nord) est le fondement matériel de l'unité nationale. Mais cette complémentarité est \textbf{fragile} : elle repose sur des infrastructures de transport (route, rail) vulnérables, sur une paix sociale nécessaire aux flux, et sur une gestion durable des ressources (forêts, pâturages, eau, sols). L'aménagement du territoire camerounais, de la décentralisation aux grands projets structurants (autoroute Yaoundé--Douala, barrage de Nachtigal, corridor Douala--N'Djamena), doit sans cesse négocier entre \textbf{valorisation des potentialités locales} et \textbf{intégration nationale}. C'est tout l'enjeu de la géographie appliquée au développement.

\coursec{Les grands groupes humains du Cameroun}

La diversité physique et climatique que nous venons d'étudier — des plaines soudano-sahéliennes du Nord aux forêts denses et humides du Sud, en passant par les hautes terres volcaniques de l'Ouest — n'est pas seulement un décor naturel. Elle constitue le cadre structurant dans lequel s'est opéré, au fil des millénaires, le \cle{peuplement} du territoire camerounais. Il existe une corrélation forte, bien que non déterministe, entre les \cle{ensembles biogéographiques} et la répartition des \cle{groupes humains}. Comprendre cette géographie humaine, c'est saisir comment les sociétés se sont adaptées à leur milieu, ont valorisé ses ressources et ont construit des identités culturelles riches de plus de 250 expressions distinctes.

\courssub{Définitions et notions fondamentales}

Avant d'analyser la répartition spatiale des populations, il est indispensable de préciser le vocabulaire scientifique qui permet d'éviter les confusions courantes et les généralisations abusives.

\begin{definition}[Ethnie et groupe ethnique]
Une \textbf{ethnie} (ou groupe ethnique) est un ensemble d'individus partageant une \textbf{langue} commune (ou des dialectes mutuellement intelligibles), une \textbf{culture} matérielle et immatérielle (traditions, rites, organisation sociale, savoir-faire techniques), une \textbf{conscience d'appartenance} collective et souvent le récit d'une \textbf{origine commune} (mythe fondateur, ancêtre éponyme). L'ethnie n'est pas une entité biologique figée : elle se construit historiquement par le brassage, les alliances, les conflits et les migrations.
\end{definition}

\begin{definition}[Groupe linguistique et famille de langues]
Un \textbf{groupe linguistique} rassemble plusieurs ethnies dont les langues présentent des parentés structurelles (vocabulaire, grammaire) prouvant une origine historique commune. Au Cameroun, on distingue principalement trois grandes familles linguistiques : les langues \textbf{afro-asiatiques} (au Nord, groupe tchadique), les langues \textbf{nilo-sahariennes} (au Nord-Est, ex. : kanouri, zaghawa) et surtout les langues \textbf{niger-congo} (dont la branche \textbf{bantoue} au Sud, \textbf{semi-bantoue} à l'Ouest et \textbf{atlantique} pour le peul/fulfulde), qui couvrent la majeure partie du territoire.
\end{definition}

\begin{definition}[Peuplement]
Le \textbf{peuplement} désigne le \textbf{processus dynamique et historique} d'installation, de déplacement, de mélange et d'adaptation des populations sur un territoire donné. Il ne s'agit pas d'une « photo » figée, mais d'un « film » : migrations bantoues venues de l'Est/Nord-Est (il y a 3000-2000 ans), poussées peules (XVIII\textsuperscript{e}-XIX\textsuperscript{e} siècles), déplacements liés aux guerres, à la traite, à la colonisation et à l'urbanisation récente.
\end{definition}

\begin{aretenir}[Distinction cruciale]
\begin{itemize}
    \item \textbf{Ethnie} = identité culturelle et linguistique vécue à l'échelle locale (ex. : Bamiléké, Bassa, Toupouri).
    \item \textbf{Groupe linguistique} = classification scientifique basée sur la parenté des langues (ex. : Bantou, Semi-bantou, Tchadique, Atlantique).
    \item \textbf{Peuplement} = dynamique historique d'occupation de l'espace (migrations, brassages).
\end{itemize}
Ne pas confondre « ethnie » et « tribu » (notion coloniale souvent réductrice) ni « ethnie » et « race » (concept biologiquement sans fondement pour l'espèce humaine).
\end{aretenir}

\begin{attention}[Piège à éviter : l'essentialisme ethnique]
\emph{Erreur fréquente :} Attribuer des traits de caractère fixes, immuables et biologiquement déterminés à chaque ethnie (ex. : « Les Bamiléké sont commerçants », « Les Peuls sont éleveurs », « Les Pygmées sont primitifs »).
\emph{Réalité :} Les identités sont \textbf{dynamiques, relationnelles et contextuelles}. Un même individu peut se dire « Bassa » au village, « Douala » en ville, « Camerounais » à l'étranger. Le brassage est ancien : les Peuls se sont sédentarisés et mélangés aux populations soudano-sahéliennes ; les Bantous ont absorbé des groupes pygmées ; les villes sont des creusets multiethniques. Réduire un groupe à un stéréotype, c'est nier l'histoire et la complexité sociale.
\end{attention}

\courssub{Les grands ensembles humains et leur répartition géographique}

La géographie humaine du Cameroun se lit comme un palimpseste où les grandes vagues migratoires ont laissé leurs traces, en forte correspondance avec les ensembles biogéographiques (climat, végétation, relief) vus précédemment.

\paragraph{1. Les peuples soudaniens du Grand Nord (Régions de l'Extrême-Nord, du Nord, de l'Adamaoua-Nord)}
Cet espace correspond à la zone \cle{soudano-sahélienne} (climat tropical sec, savane herbeuse à arbustive). Les populations y pratiquent traditionnellement une agriculture pluviale (mil, sorgho, maïs, riz, coton) et un élevage extensif.

\begin{itemize}
    \item \textbf{Groupes tchadiques (famille afro-asiatique) :} \textbf{Kotoko} (pêcheurs et agriculteurs des rives du Logone et du Lac Tchad, héritiers de la civilisation Sao), \textbf{Mousgoum} (célèbres pour leurs cases coquilles en terre, agriculteurs des plaines inondables), \textbf{Massa} (pêcheurs-agriculteurs du Bas-Logone), \textbf{Toupouri} (agriculteurs des plaines de Kaélé), \textbf{Moundang} (agriculteurs des piémonts des Mandara).
    \item \textbf{Peuls / Foulbé (famille nigéro-congolaise, branche atlantique, langue fulfulde) :} Arrivés par vagues successives (XVII\textsuperscript{e}-XIX\textsuperscript{e} s.), ils forment une aristocratie pastorale et commerciale islamisée. Distinction entre \textbf{Peuls nomades} (Bororo, Wodaabe : transhumance pure, forte mobilité) et \textbf{Peuls sédentaires} (agro-éleveurs, commerçants, élites urbaines : Garoua, Maroua, Ngaoundéré). L'islam est un facteur majeur d'unité culturelle et politique (lamidats).
\end{itemize}

\paragraph{2. Les peuples des montagnes et plateaux (Région de l'Adamaoua, Monts Mandara, Piémonts)}
Zone de transition (climat soudano-guinéen, relief accidenté), refuge historique et carrefour migratoire.

\begin{itemize}
    \item \textbf{Populations des Monts Mandara (Kirdi) :} Terme générique (péjoratif à l'origine, signifiant « païens » en kanouri) recouvrant une mosaïque de petits groupes (Mafa, Mofu, Guiziga, Podoko, Zulgo, etc.) qui ont résisté à l'islamisation et à la razzia peule en se retranchant dans les montagnes. Agriculture en terrasses, organisation sociale segmentaire.
    \item \textbf{Tikar :} Groupe historique majeur de l'Adamaoua et de l'Ouest, à l'origine de nombreux royaumes (Bamoun, Bafut, Kom, Nso, Bankim). Langue semi-bantoue. Artisans réputés (fer, perles, architecture).
    \item \textbf{Bamoun :} Royaume centralisé ancien (Foumban), écriture propre (Shümom), islamisés tardivement (début XX\textsuperscript{e} s.) sous le roi Njoya. Art raffiné (bronze, perles, tissus).
    \item \textbf{Autres groupes de l'Adamaoua :} Mbaya, Vute, Nyamnyam (Djém), Gbaya (venus de RCA), etc. Zone de fort brassage.
\end{itemize}

\paragraph{3. Les peuples semi-bantous des Grassfields (Régions de l'Ouest et du Nord-Ouest)}
Correspondent à l'ensemble \cle{montagnard volcanique} (altitudes 1000-2500 m, climat tempéré, sols ferralitiques riches, forte pluviométrie). Haute densité de population historique.

\begin{itemize}
    \item \textbf{Bamiléké (Ouest) :} Ensemble de groupes (Bafoussam, Bandjoun, Baham, Bangangté, Dschang, etc.) partageant la langue \textbf{ghomala'} (dialectes variés). Organisation en \textbf{chefferies traditionnelles puissantes} (fonction sacrée du chef, société secrète, économie de redistribution). Agriculture intensive (terrasses, fumure, cultures associées : maïs, haricot, taro, café, thé). Forte diaspora commerciale et entrepreneuriale nationale.
    \item \textbf{Groupes du Nord-Ouest (Grassfields anglophones) :} \textbf{Bamenda} (Bali, Bafut, Kom, Nso, Oku, Widikum, etc.). Langues semi-bantoues (groupe Ring, Bantuide). Chefferies (Fons) structurées, agriculture intensive similaire, forte émigration vers les plantations du Sud-Ouest et les villes.
\end{itemize}

\begin{exemplebox}[Densité et intensification agraire chez les Bamiléké]
La région de l'Ouest (territoire Bamiléké) affiche des densités dépassant \textbf{200 à 300 hab./km²} (localement > 500), records nationaux pour le milieu rural. Face à la pression foncière, les paysans ont développé une \textbf{agriculture intensive sur terrasses} (ex. : paysage de Dschang, Bafou), utilisant le fumier animal et les résidus de culture pour maintenir la fertilité. C'est un modèle d'adaptation anthropique au milieu montagnard.
\end{exemplebox}

\paragraph{4. Les peuples bantous du Sud forestier (Régions du Centre, du Sud, de l'Est, du Littoral, du Sud-Ouest)}
Correspondent à l'ensemble \cle{forestier guinéo-congolais} (climat équatorial, forêt dense, sols ferrallitiques lessivés, fragiles). Langues de la famille \textbf{bantoue} (zone A : Beti-Fang ; zone B : Bassa-Douala ; zone C : langues des forêts du Sud-Est).

\begin{itemize}
    \item \textbf{Groupe Beti-Fang (Centre, Sud, Est) :} \textbf{Ewondo} (Yaoundé), \textbf{Bulu} (Sangmélima, Ebolowa), \textbf{Fang} (Sud, frontière Gabon/Congo), \textbf{Eton}, \textbf{Bene}, \textbf{Manguissa}. Organisation en lignages patrilinéaires, chefferies de village. Agriculture itinérante sur brûlis (cacao, café, vivriers : manioc, plantain, macabo), chasse, cueillette.
    \item \textbf{Groupe Bassa-Douala (Littoral, Sud-Ouest) :} \textbf{Bassa} (Edéa, Nkongsamba), \textbf{Douala} (Wouri, côte), \textbf{Bakoko}, \textbf{Mbo}. Peuples côtiers et fluviaux : pêche lagunaire et maritime, commerce précoce avec Européens (XV\textsuperscript{e} s.), agriculture de plantation (palmier à huile, cacao).
    \item \textbf{Autres groupes du Sud-Est :} \textbf{Bakweri} (pentes du Mont Cameroun, plantationnistes), \textbf{Bakola}, \textbf{Bagyeli} (voir Pygmées), \textbf{Kwasio}, \textbf{Mvae}.
\end{itemize}

\paragraph{5. Les peuples premiers de la forêt : les Pygmées (Régions de l'Est et du Sud)}
Ils occupent l'ensemble \cle{forestier dense humide} du Sud-Est (bassin du Dja, Boumba-Ngoko, Sangha). Ce sont des populations \textbf{chasseuses-cueilleuses} (parfois agriculteurs occasionnels aujourd'hui), de petite taille morphologique, maîtrisant une connaissance encyclopédique de la forêt (pharmacopée, pistage, miel, chenilles, gibier).

\begin{itemize}
    \item \textbf{Baka} (Est, Sud-Est, vers Yokadouma, Moloundou, Lomié) : plus nombreux (~ 40 000).
    \item \textbf{Bagyeli / Bakola} (Sud, Océan, Kribi, Campo, Bipindi) : plus petits groupes (~ 5 000-10 000).
\end{itemize}
Leurs modes de vie sont \textbf{menacés} : sédentarisation forcée, perte d'accès aux ressources forestières (parcs nationaux, concessions forestières, agriculture villageoise), marginalisation sociale, non-reconnaissance foncière coutumière.

\begin{savaistu}[Les Pygmées : gardiens de la forêt et patrimoine UNESCO]
Les Pygmées \textbf{Baka} sont parmi les \textbf{plus anciens habitants} de la forêt camerounaise (présence attestée depuis des millénaires). La \textbf{polyphonie vocale} des Pygmées \textbf{Aka} (proches parents des Baka, zone RCA/Congo/Cameroun), avec ses chants yeyi, polyrythmies complexes et technique du yodel, est inscrite depuis 2008 au \textbf{Patrimoine culturel immatériel de l'humanité} par l'UNESCO. La connaissance de la biodiversité forestière par les Baka (plus de 300 espèces végétales utilisées) intéresse la recherche pharmaceutique mondiale. Pourtant, ils restent l'un des groupes les plus vulnérables du pays (accès difficile à l'éducation, à la santé, à l'état civil).
\end{savaistu}

\courssub{Langues et dynamique démographique}

\begin{exemplebox}[Chiffres clés du peuplement camerounais (estimations récentes)]
\begin{itemize}
    \item \textbf{Population totale :} environ \textbf{28 millions d'habitants} (2023-2024).
    \item \textbf{Nombre de groupes ethniques / langues nationales :} \textbf{plus de 250} (selon les recensements linguistiques, ex. : \textit{Atlas linguistique du Cameroun} - ALCA).
    \item \textbf{Langues officielles :} \textbf{Français} (8 régions, ~ 80\% pop.) et \textbf{Anglais} (2 régions, ~ 20\% pop.), héritage colonial (mandat France / Royaume-Uni).
    \item \textbf{Langues véhiculaires :} \textbf{Fulfulde} (Nord, Adamaoua), \textbf{Ewondo} (Centre, Sud, Est), \textbf{Douala} (Littoral), \textbf{Bamiléké} (Ouest), \textbf{Pidgin English} (Nord-Ouest, Sud-Ouest, villes).
    \item \textbf{Taux d'urbanisation :} ~ 58\% (croissance rapide de Yaoundé ~ 4,5 M, Douala ~ 4 M).
\end{itemize}
\end{exemplebox}

\begin{popfigure}
\begin{tikzpicture}[scale=0.85]
% Définition des points clés du contour simplifié (sens horaire depuis Extrême-Nord)
\coordinate (EN) at (0,0);        % Extrême-Nord (Kousseri)
\coordinate (N) at (1,0.2);       % Nord (Garoua approx)
\coordinate (AD_E) at (2.5,0.5);  % Adamaoua Est
\coordinate (AD_C) at (3.5,1.5);  % Adamaoua Centre (Ngaoundéré)
\coordinate (AD_W) at (4,3);      % Adamaoua Ouest / Monts Adamawa
\coordinate (NO) at (4.2,5);      % Nord-Ouest (Bamenda)
\coordinate (O) at (3.5,6.5);     % Ouest (Bafoussam)
\coordinate (C) at (2,7.5);       % Centre (Yaoundé)
\coordinate (ES) at (0.5,7);      % Est (Bertoua)
\coordinate (S) at (-0.5,6);      % Sud (Sangmélima)
\coordinate (SW) at (-1,4);       % Sud-Ouest (Buéa)
\coordinate (LT) at (-1,2);       % Littoral (Douala)
\coordinate (EN_W) at (-0.5,0.5); % Extrême-Nord Ouest (Maroua)

% Contour principal
\draw[popdark, thick] (EN) -- (N) -- (AD_E) -- (AD_C) -- (AD_W) -- (NO) -- (O) -- (C) -- (ES) -- (S) -- (SW) -- (LT) -- (EN_W) -- cycle;

% Zones biogéographiques / humaines (remplissages non chevauchants)
% 1. Nord Soudano-Sahélien (Orange)
\fill[poporange!20] (EN) -- (N) -- (AD_E) -- (AD_C) -- (3,2.5) -- (1.5,2) -- (EN_W) -- cycle;

% 2. Adamaoua / Monts Mandara Transition (Violet)
\fill[poppurple!15] (3,2.5) -- (AD_C) -- (AD_W) -- (3.5,4.5) -- (2,4.5) -- (1.5,2) -- cycle;

% 3. Grassfields / Hautes Terres Volcaniques (Vert)
\fill[popgreen!20] (2,4.5) -- (3.5,4.5) -- (AD_W) -- (3.5,2.5) -- (2.5,3) -- cycle;

% 4. Sud Forestier / Bantous (Bleu) - Le reste du polygone principal
% On utilise le "even odd rule" ou on trace le polygone sud explicitement
\fill[popblue!15] 
    (EN_W) -- (LT) -- (SW) -- (S) -- (ES) -- (C) -- (O) -- (2,4.5) -- (1.5,2) -- (EN_W) -- cycle;

% Villes repères
\node[circle, fill=popdark, inner sep=1.5pt] (Ya) at (1.5,4.5) {};
\node[right, font=\tiny] at (Ya) {Yaoundé};
\node[circle, fill=popdark, inner sep=1.5pt] (Do) at (3.2,3.8) {};
\node[above right, font=\tiny] at (Do) {Douala};
\node[circle, fill=popdark, inner sep=1.5pt] (Ga) at (3.5,2.5) {};
\node[above, font=\tiny] at (Ga) {Garoua};
\node[circle, fill=popdark, inner sep=1.5pt] (Ma) at (2.5,1.2) {};
\node[below, font=\tiny] at (Ma) {Maroua};
\node[circle, fill=popdark, inner sep=1.5pt] (Ng) at (3.2,4.2) {};
\node[above right, font=\tiny] at (Ng) {Ngaoundéré};
\node[circle, fill=popdark, inner sep=1.5pt] (Bf) at (2.2,3.5) {};
\node[left, font=\tiny] at (Bf) {Bafoussam};
\node[circle, fill=popdark, inner sep=1.5pt] (Be) at (0.5,5.8) {};
\node[left, font=\tiny] at (Be) {Bertoua};

% Étiquettes groupes humains
\node[font=\small\bfseries, poporange!80!popdark, align=center] at (1.5,1) {SOUDANIENS\\(Kotoko, Massa,\\Toupouri, Moundang)};
\node[font=\small\bfseries, poppurple!80!popdark, align=center] at (3.5,3.5) {MONTAGNES /\\(Kirdi, Tikar,\\Bamoun, Peuls)};
\node[font=\small\bfseries, popgreen!70!popdark, align=center] at (2.8,3.2) {SEMI-BANTOUS\\(Bamiléké,\\Bamenda, Widikum)};
\node[font=\small\bfseries, popblue!70!popdark, align=center] at (0,4.5) {BANTOUS\\(Beti, Bassa,\\Douala, Bakweri)};
\node[font=\small\bfseries, poppink!80!popdark, align=center] at (-0.5,3) {PYGMÉES\\(Baka, Bagyeli,\\Bakola)};

% Flèche Nord (orientation)
\node[font=\tiny] at (4.5, 6.5) {N};
\draw[->, thick, popdark] (4.5, 6.2) -- (4.5, 5.7);

% Échelle indicative
\draw[popdark, thick] (4.5, 0.5) -- (5.5, 0.5);
\node[font=\tiny, below] at (5, 0.5) {~ 300 km};

% Flèches Peuls (transhumance)
\draw[->, thick, poporange!70!popdark, decorate, decoration={snake, amplitude=1pt, segment length=4pt}] (3.5,2.5) -- (2.5,3.5);
\node[font=\tiny, poporange!70!popdark, rotate=30] at (3.2,3.1) {Transhumance Peule};

\end{tikzpicture}
\legende{Carte schématique des grands groupes humains du Cameroun. La répartition recoupe largement les ensembles biogéographiques : zone soudano-sahélienne (Nord), zone de transition montagnarde (Adamaoua/Mandara), zone des hautes terres volcaniques (Grassfields Ouest/Nord-Ouest), zone forestière équatoriale (Sud). Les flèches indiquent les axes de transhumance peule.}
\end{popfigure}

\begin{popfigure}
\begin{tikzpicture}
\begin{axis}[
    width=\linewidth,
    height=8cm,
    ybar,
    bar width=12pt,
    enlarge x limits=0.2,
    ylabel={Part de la population (\% approximatif)},
    ylabel style={font=\small, popdark},
    symbolic x coords={Bantous, Semi-Bantous, Soudaniens, Peuls, Autres},
    xtick=data,
    xticklabel style={font=\tiny, rotate=30, anchor=east, popdark},
    ytick={0,10,20,30,40},
    ymin=0, ymax=40,
    ymajorgrids=true,
    grid style={popline},
    legend style={at={(0.5,-0.2)}, anchor=north, legend columns=3, font=\tiny, /tikz/every even column/.append style={column sep=5pt}},
    nodes near coords={\pgfmathprintnumber\pgfplotspointmeta\%},
    nodes near coords style={font=\tiny, popdark, anchor=south},
]
\addplot[fill=popblue, draw=popblue!50!popdark] coordinates {(Bantous,32)};
\addplot[fill=popgreen, draw=popgreen!50!popdark] coordinates {(Semi-Bantous,28)};
\addplot[fill=poporange, draw=poporange!50!popdark] coordinates {(Soudaniens,18)};
\addplot[fill=poppurple, draw=poppurple!50!popdark] coordinates {(Peuls,12)};
\addplot[fill=poppink, draw=poppink!50!popdark] coordinates {(Autres,10)};
\legend{Bantous, Semi-Bantous, Soudaniens, Peuls, Autres}
\end{axis}
\end{tikzpicture}
\legende{Diagramme en barres de la répartition approximative des grands ensembles linguistiques au Cameroun (estimations basées sur données linguistiques et démographiques). Les groupes bantous et semi-bantous représentent ensemble environ 60\% de la population.}
\end{popfigure}

\courssub{L'unité dans la diversité : brassage, mobilité et enjeux}

La carte et le diagramme montrent une structuration claire, mais la réalité vécue est celle du \textbf{brassage permanent}.

\begin{methode}[Analyser la corrélation Milieu / Groupe humain (ex. sujet Bac)]
\begin{enumerate}
    \item \textbf{Identifier l'ensemble biogéographique :} Climat, végétation, relief, ressources (ex. : forêt dense, sol ferrallitique, pluviométrie > 1500 mm).
    \item \textbf{Identifier le(s) groupe(s) humain(s) dominant(s) :} Langue, activités économiques traditionnelles, habitat, organisation sociale (ex. : Beti, agriculture itinérante sur brûlis, cacao, habitat dispersé le long des pistes).
    \item \textbf{Expliquer l'adaptation :} Comment le groupe valorise les potentialités et contourne les contraintes du milieu (ex. : culture du cacao sous couvert forestier, exploitation du bois, chasse).
    \item \textbf{Nuancer :} Mentionner les dynamiques actuelles qui brisent le déterminisme (ex. : urbanisation, plantations agro-industrielles, aires protégées, migrations pendulaires, scolarisation).
    \item \textbf{Conclure :} Le milieu \textbf{oriente} sans \textbf{déterminer} ; l'histoire, la technique et la politique sont des facteurs tout aussi puissants.
\end{enumerate}
\end{methode}

\begin{exoresolu}[Analyse d'une corrélation Milieu / Population]
\textbf{Énoncé :} À l'aide d'exemples précis pris au Cameroun, montrez que la répartition des groupes humains recoupe largement les ensembles biogéographiques, tout en soulignant les limites de ce déterminisme.

\textbf{Solution détaillée :}
\begin{enumerate}
    \item \textbf{Recoupement fort (arguments pour) :}
    \begin{itemize}
        \item \textit{Grand Nord (Soudano-sahélien) :} Climat sec (6-8 mois), savane. $\rightarrow$ Peuples tchadiques (Massa, Mousgoum, Toupouri) : agriculture pluviale (mil/sorgho), pêche, élevage bovin. Peuls : pastoralisme extensif transhumant adapté à la saisonnalité des pâturages.
        \item \textit{Grassfields (Montagnard volcanique) :} Altitude, climat tempéré, sols riches, forte pop. $\rightarrow$ Bamiléké / Bamenda : agriculture intensive sur terrasses, chefferies centralisées gérant la rareté foncière, cultures de rente (café arabica, thé).
        \item \textit{Sud Forestier (Équatorial) :} Forêt dense, sols pauvres (lessivage), pluviométrie forte. $\rightarrow$ Bantous (Beti, Bassa) : agriculture itinérante sur brûlis (manioc, plantain), cultures de rente (cacao, café robusta, palmier) adaptées à l'ombrage, habitat linéaire le long des axes. Pygmées (Baka) : chasse-cueillette nomade adaptée à la forêt primaire.
    \end{itemize}
    \item \textbf{Limites du déterminisme / Facteurs de brassage (arguments contre) :}
    \begin{itemize}
        \item \textbf{Migrations historiques :} Expansion bantoue (venant de l'Est) a traversé la savane pour aller en forêt ; Peuls (venant de l'Ouest) se sont installés en zone soudano-sahélienne puis montagnarde.
        \item \textbf{Urbanisation :} Yaoundé (zone Beti) = ville multiethnique (Bamiléké, Nordistes, Étrangers). Douala (zone Bassa) = capitale économique cosmopolite.
        \item \textbf{Plantations agro-industrielles :} CDC (Sud-Ouest, zone Bakweri) attire main-d'œuvre Bamenda, Bamiléké, Nordistes.
        \item \textbf{Politique / Éducation :} Langues officielles (FR/EN), administration, fonction publique créent une identité nationale par-dessus les identités ethniques.
    \end{itemize}
    \item \textbf{Synthèse :} Le milieu biophysique offre un \textbf{cadre de potentialités} (ex. : la forêt \emph{permet} la culture du cacao, elle ne l'\emph{impose} pas). Les sociétés font des \textbf{choix techniques, économiques et politiques} qui transforment le milieu et leur propre répartition.
\end{enumerate}
\end{exoresolu}

\begin{aretenir}[Synthèse : Diversité humaine et unité nationale]
\begin{itemize}
    \item Le Cameroun compte \textbf{> 250 groupes ethniques} regroupés en \textbf{4 grands ensembles linguistiques} (Afro-asiatique, Nilo-saharien, Niger-Congo Bantou, Niger-Congo Semi-bantou/Atlantique).
    \item La répartition spatiale \textbf{recoupe les ensembles biogéographiques} (Nord soudano-sahélien / Peuls-Tchadiens ; Hautes terres / Semi-bantous ; Forêt / Bantous-Pygmées).
    \item \textbf{Deux langues officielles} (Français, Anglais) et des \textbf{langues véhiculaires} (Fulfulde, Ewondo, Pidgin, Camfranglais) assurent la communication interethnique.
    \item \textbf{Défis :} Gestion de la diversité (fédéralisme/unitarisme, quota ethno-régionaux), reconnaissance des droits des peuples autochtones (Pygmées), lutte contre les stéréotypes et les discours de haine, valorisation du patrimoine culturel immatériel.
\end{itemize}
\end{aretenir}

\coursec{L'unité dans la diversité : atouts et défis}

Après avoir identifié les grands groupes humains qui peuplent le territoire camerounais, il convient désormais de comprendre comment cette mosaïque physique et humaine, loin d'être une source de fragmentation, constitue le socle même de la nation. Le Cameroun ne s'est pas construit \emph{malgré} sa diversité, mais \emph{grâce à elle}. Cette section analyse les mécanismes de cette construction, les richesses qu'elle engendre, les tensions qu'elle traverse et les réponses institutionnelles et citoyennes apportées pour faire de l'unité une réalité vivante.

\courssub{La diversité comme richesse : complémentarité des productions}

La diversité biogéographique du Cameroun n'est pas un décor statique ; elle est le moteur d'une économie d'échanges interne puissante. Chaque ensemble biophysique offre des potentialités spécifiques qui, mises en relation, assurent une relative autonomie alimentaire et économique au pays.

\begin{exemplebox}[La complémentarité Nord-Sud et Ouest-Est]
\emph{Observation :} Au marché central de Yaoundé ou de Douala, on trouve simultanément : du riz de la plaine de la Bénoué (Nord), des ignames et du maïs de l'Ouest et de l'Adamaoua, du plantain et du cacao du Centre-Sud, du poisson du Littoral et du Lac Tchad, du bétail (bœuf, mouton) de l'Adamaoua et du Nord.
\par
\emph{Analyse :} Cette présence simultanée résulte de la complémentarité écologique :
\begin{itemize}
    \item \textbf{Zone soudano-sahélienne (Extrême-Nord, Nord, partie Adamaoua) :} Céréales (mil, sorgho, maïs, riz), coton, élevage bovin extensif, oignons (Logone-et-Chari).
    \item \textbf{Zone de transition guinéenne (Adamaoua, Nord-Ouest, Ouest) :} Maïs, pomme de terre, haricot, café arabica, élevage bovin amélioré, cultures maraîchères (altitude).
    \item \textbf{Zone équatoriale (Centre, Sud, Est, Littoral, Sud-Ouest) :} Cultures d'exportation (cacao, café robusta, palmier à huile, banane, hévéa), cultures vivrières de forêt (plantain, manioc, macabo, igname), pêche maritime et lagunaire, exploitation forestière.
\end{itemize}
\emph{Conclusion :} Les flux commerciaux internes (camions, train, pirogues) structurent le territoire national bien plus efficacement que les discours politiques. La sécurité alimentaire des villes du Sud dépend de la production du Nord ; le revenu d'exportation du Sud finance l'équipement du Nord.
\end{exemplebox}

\begin{aretenir}[Richesse de la complémentarité]
La diversité physique (relief, climat, végétation, sols) génère une \textbf{spécialisation régionale} naturelle. L'unité territoriale permet la \textbf{libre circulation des biens}, transformant cette spécialisation en \textbf{complémentarité économique}. C'est le fondement matériel de la solidarité nationale.
\end{aretenir}

\courssub{Les brassages humains : le creuset national}

Si la complémentarité économique crée une interdépendance matérielle, les migrations intérieures tissent le lien social et culturel. Le Cameroun est un pays de brassage ancien et contemporain.

\begin{definition}[Brassage]
Mélange progressif de populations et de cultures résultant des migrations (déplacements durables), des échanges commerciaux, des mariages mixtes, de la scolarisation commune et de l'administration partagée. Il produit des identités hybrides, urbaines et nationales, sans effacer les identités d'origine.
\end{definition}

\begin{popfigure}
\begin{tikzpicture}[scale=0.8, every node/.style={font=\small}]
    % Contour très simplifié du Cameroun (polygone approximatif)
    \draw[popdark, thick] (0,0) -- (1,2.5) -- (3,3.5) -- (6,3) -- (7.5,1.5) -- (7,0) -- (5,-0.5) -- (3,-0.5) -- (1,0.5) -- cycle;
    
    % Villes principales (points)
    \node[circle, fill=popblue, inner sep=2pt] (douala) at (5.5, 1.2) {};
    \node[circle, fill=popblue, inner sep=2pt] (yaounde) at (4.5, 2.2) {};
    \node[circle, fill=popblue, inner sep=2pt] (ngaoundere) at (3.5, 3.0) {};
    \node[circle, fill=popblue, inner sep=2pt] (maroua) at (2.0, 3.2) {};
    \node[circle, fill=popblue, inner sep=2pt] (bafoussam) at (4.0, 1.8) {};
    \node[circle, fill=popblue, inner sep=2pt] (bamenda) at (3.0, 2.0) {};
    \node[circle, fill=popblue, inner sep=2pt] (bertoua) at (5.5, 2.8) {};
    \node[circle, fill=popblue, inner sep=2pt] (garoua) at (1.5, 2.8) {};

    % Noms des villes
    \node[below right] at (douala) {Douala};
    \node[above right] at (yaounde) {Yaoundé};
    \node[above] at (ngaoundere) {Ngaoundéré};
    \node[above left] at (maroua) {Maroua};
    \node[left] at (bafoussam) {Bafoussam};
    \node[above left] at (bamenda) {Bamenda};
    \node[right] at (bertoua) {Bertoua};
    \node[below left] at (garoua) {Garoua};

    % Flèches migrations Nord -> Sud (Grands axes)
    \draw[poporange, very thick, -{Stealth[length=4mm]}, bend left=15] (maroua) to node[midway, above, sloped, fill=white, inner sep=1pt] {\tiny Fonctionnaires, Commerçants, Étudiants} (yaounde);
    \draw[poporange, very thick, -{Stealth[length=4mm]}, bend right=10] (garoua) to node[midway, below, sloped, fill=white, inner sep=1pt] {\tiny Transhumance, Commerce bétail} (douala);
    \draw[poporange, thick, -{Stealth[length=3mm]}] (ngaoundere) -- (yaounde);
    \draw[poporange, thick, -{Stealth[length=3mm]}] (ngaoundere) -- (douala);

    % Flèches migrations Ouest -> Littoral/Centre (Bamiléké, Bamoun, Anglophones)
    \draw[popgreen, very thick, -{Stealth[length=4mm]}, bend right=15] (bafoussam) to node[midway, right, sloped, fill=white, inner sep=1pt] {\tiny Entrepreneurs, Commerçants, Élites} (douala);
    \draw[popgreen, thick, -{Stealth[length=3mm]}, bend left=10] (bafoussam) to node[midway, left, sloped, fill=white, inner sep=1pt] {\tiny Fonctionnaires} (yaounde);
    \draw[popgreen, thick, -{Stealth[length=3mm]}] (bamenda) -- (douala);
    \draw[popgreen, thick, -{Stealth[length=3mm]}] (bamenda) -- (yaounde);

    % Flèches Est -> Centres urbains
    \draw[poppurple, thick, -{Stealth[length=3mm]}] (bertoua) -- (yaounde);
    \draw[poppurple, thick, -{Stealth[length=3mm]}] (bertoua) -- (douala);

    % Légende
    \begin{scope}[shift={(8.2, 3.5)}]
        \draw[poporange, very thick, -{Stealth[length=3mm]}] (0,0.3) -- (1.5,0.3) node[right] {Nord $\to$ Sud (Administration, Commerce, Études)};
        \draw[popgreen, very thick, -{Stealth[length=3mm]}] (0,0) -- (1.5,0) node[right] {Ouest $\to$ Littoral/Centre (Entrepreneuriat, Fonction publique)};
        \draw[poppurple, thick, -{Stealth[length=3mm]}] (0,-0.3) -- (1.5,-0.3) node[right] {Est $\to$ Centres urbains};
    \end{scope}
    
    % Nord symbolique
    \node[popdark, rotate=90] at (-0.5, 1.5) {N};
\end{tikzpicture}
\legende{Schéma des flux migratoires intérieurs majeurs. Épaisseur des flèches $\approx$ intensité relative. Flèches orange : axe Nord-Sud (fonctionnaires, commerçants bétail, étudiants). Flèches vertes : axe Ouest-Centre/Littoral (diaspora Bamiléké/Bamoun, entrepreneurs). Flèches violettes : axe Est-Centres urbains.}
\end{popfigure}

\begin{exemplebox}[Douala et Yaoundé : miroirs du brassage]
Selon les recensements (RGPH 2005, projections INS), Douala et Yaoundé concentrent à elles deux \textbf{près de 30 \% de la population nationale} (plus de 7 millions d'habitants sur ~28 millions en 2023). 
\begin{itemize}
    \item \textbf{Yaoundé :} Capitale administrative $\to$ attraction massive des fonctionnaires de toutes les régions (Nord, Ouest, Est, Sud). Quartiers à forte identité régionale (Mvog-Ada $\leftarrow$ Bassa ; Nkolbisson $\leftarrow$ Bamiléké ; Melen $\leftarrow$ Beti/Fang ; quartier Hausa $\leftarrow$ Nord).
    \item \textbf{Douala :} Capitale économique $\to$ attraction des entrepreneurs (Bamiléké, Bamoun, Haoussas, Grecs, Libanais, Français), dockers, pêcheurs (Nord, Côte), étudiants.
\end{itemize}
Ces villes sont des \emph{creusets} où le \cle{bilinguisme} (français/anglais) et le \cle{camfranglais} (langue vernaculaire urbaine) s'inventent au quotidien.
\end{exemplebox}

\begin{methode}[Analyser un flux migratoire intérieur (démarche géographique)]
Pour rendre compte d'un brassage dans une copie ou un exposé :
\begin{enumerate}
    \item \textbf{Identifier l'origine et la destination} (régions, villes, milieux ruraux/urbains).
    \item \textbf{Qualifier les migrants} (âge, sexe, qualification : étudiants, fonctionnaires, commerçants, saisonniers, réfugiés climatiques/conflits).
    \item \textbf{Expliquer les causes (facteurs d'expulsion/attraction)} : pression foncière (Ouest), insécurité/emploi (Nord), attractivité administrative (Yaoundé), attractivité marchande (Douala).
    \item \textbf{Décrire les modalités d'insertion} : quartiers d'installation, activités économiques, associations de développement (ex. \emph{Groupements Bamiléké}, \emph{Association des Élèves et Étudiants du Nord}), mariages mixtes.
    \item \textbf{Évaluer les effets} : sur l'origine (revenus par transferts, désertification des campagnes) et sur la destination (dynamisme économique, tensions foncières, enrichissement culturel).
\end{enumerate}
\end{methode}

\courssub{Les instruments de l'unité nationale : du symbole à la pratique}

L'unité ne va pas de soi ; elle est \emph{construite} par des dispositifs juridiques, symboliques et institutionnels.

\begin{propriete}[Les piliers institutionnels de l'unité camerounaise]
\begin{itemize}
    \item \textbf{La Devise :} \og Paix -- Travail -- Patrie \fg. Elle ordonne le contrat social : la \cle{Paix} comme condition (absence de violence, résolution pacifique), le \cle{Travail} comme devoir et mérite (lutte contre la rente, valorisation de l'effort), la \cle{Patrie} comme communauté de destin supérieure aux allégeances primaires.
    \item \textbf{Le Bilinguisme officiel :} Français et Anglais (loi n°96-06 du 18 janvier 1996, Constitution). Il traduit l'héritage historique (réunification 1961) et impose l'apprentissage de la langue de l'autre comme devoir citoyen.
    \item \textbf{Les Symboles nationaux :} Drapeau (vert-rouge-jaune + étoile de l'unité), Hymne (\og Ô Cameroun, berceau de nos ancêtres \fg), Sceau, Armoiries. Ils sont enseignés à l'école, honorés lors des cérémonies (11 février Fête de la Jeunesse, 20 mai Fête Nationale).
    \item \textbf{Le Brassage scolaire et administratif :} Affectation des fonctionnaires (enseignants, médecins, administrateurs) hors de leur région d'origine (principe de la \cle{mobilité nationale}). Internats, lycées techniques et universités (Ngaoundéré, Maroua, Bertoua, Buéa, Dschang, Douala, Yaoundé I/II) mélangent les jeunes de tout le pays.
\end{itemize}
\end{propriete}

\begin{savaistu}[La devise sur les billets de banque]
La devise \og \textbf{Paix-Travail-Patrie} \fg figure sur \textbf{tous les billets de banque} émis par la BEAC pour le Cameroun (série 2002 et série 2022) ainsi que sur les documents officiels (passeports, cartes d'identité, actes notariés). C'est un rappel visuel, quotidien et obligatoire : chaque transaction monétaire, chaque acte administratif engage le citoyen dans le pacte national. C'est une pédagogie d'État par l'objet.
\end{savaistu}

\courssub{Les défis : quand la diversité devient fracture}

La diversité mal gérée engendre des tensions. Le géographe doit les nommer pour les comprendre et proposer des régulations.

\begin{attention}[Piège conceptuel : opposer Unité et Diversité]
\emph{Erreur fréquente :} Dire \og L'unité menace la diversité \fg ou \og La diversité menace l'unité \fg.
\emph{Réalité camerounaise :} L'unité camerounaise \textbf{se construit \emph{à travers} la diversité} (fédéralisme historique, décentralisation actuelle, bilinguisme, quota régionaux dans la fonction publique). Nier la diversité (assimilationnisme) ou nier l'unité (séparatisme, tribalisme) sont les deux symétriques de l'échec. La bonne question : \og Comment l'État organise-t-il la diversité pour produire de l'unité ? \fg
\end{attention}

\begin{center}
\begin{tabularx}{\linewidth}{|l|X|X|}
\hline
\rowcolor{popblueL}
\textbf{Défi} & \textbf{Manifestations géographiques \& sociales} & \textbf{Réponses / Régulations (en cours ou nécessaires)} \\
\hline
\textbf{Tribalisme / Régionalisme} & Vote ethnique, népotisme dans l'embauche, discours de haine sur réseaux sociaux, revendications sécessionnistes (NOSO). & Éducation à la citoyenneté, loi sur le discours de haine (2019), Commission Nationale du Bilinguisme et du Multiculturalisme, décentralisation (régions, conseils régionaux). \\
\hline
\textbf{Inégalités de développement} & Indicateurs (taux de pauvreté, scolarisation, santé, routes) : Extrême-Nord, Nord, Adamaoua, Est $>$ Littoral, Centre, Ouest, Sud-Ouest. & Plan National de Développement (PND 2020-2030), Programme d'Urgence Triennal (PLANUT), discrimination positive (concours spéciaux, seuils d'admission différenciés), investissements infrastructures (autoroute Yaoundé-Ngaoundéré, barrage de Nachtigal, port de Kribi). \\
\hline
\textbf{Conflits agriculteurs-éleveurs} & Nord, Adamaoua, Nord-Ouest, Est. Destruction cultures par bétail, blocage couloirs de transhumance, violence armée. & Commissions de conciliation (préfets, lamibé, chefs coutumiers), aménagement couloirs/aires de pâture, sédentarisation partielle (ranchs), code pastoral. \\
\hline
\textbf{Marginalisation des minorités (Pygmées Baka, Bakola, Bagyeli)} & Non-reconnaissance droits fonciers (forêt), exclusion scolaire/sanitaire, exploitation travail, perte culture. & Plans d'action nationaux (PNDP, ONG), écoles mobiles, reconnaissance statut \og populations autochtones \fg (Déclaration ONU 2007), consultation libre, préalable et éclairée (CLPE) pour projets forestiers/miniers. \\
\hline
\end{tabularx}
\end{center}

\courssub{Application à la vie courante : le \og Vivre-ensemble \fg au quotidien}

Les savoir-faire du programme (\og Appliquer les notions à des situations concrètes \fg, \og Rédiger et justifier une démarche \fg) trouvent ici leur terrain d'exercice immédiat.

\begin{experience}[Enquête de terrain : \og Mon quartier, mon marché, ma classe \fg]
\textbf{Objectif :} Observer et analyser le brassage réel dans son environnement immédiat.
\par
\textbf{Protocole :}
\begin{enumerate}
    \item \textbf{Observation (1h) :} Au marché ou dans la cour de l'école, noter sur 30 minutes : langues entendues, tenues vestimentaires, produits vendus (origine régionale), interactions (marchand/client, élèves entre eux).
    \item \textbf{Entretien court (5 min) :} Interroger 3 personnes (commerçants, élèves) : \og D'où venez-vous ? Depuis quand ici ? Quelle langue parlez-vous à la maison ? Avez-vous des amis d'autres régions ? \fg
    \item \textbf{Cartographie mentale :} Dessiner un croquis du lieu : zones d'activités, flux, groupes repérés.
\end{enumerate}
\textbf{Interprétation attendue :} Constater que le \cle{vivre-ensemble} n'est pas un slogan mais une pratique quotidienne : le boulanger haoussa vend au client bamiléké en camfranglais ; l'élève du Nord aide son camarade du Sud en maths ; le chef de quartier médie un conflit foncier entre autochtones et allogènes.
\textbf{Rédaction :} Produire un rapport structuré (Introduction, Méthode, Résultats, Analyse, Conclusion) en mobilisant le vocabulaire : \emph{brassage, complémentarité, tolérance, identité multiple, citoyenneté}.
\end{experience}

\begin{methode}[Rédiger et justifier une conclusion en géographie (exigible au Bac)]
Pour clore un devoir ou une dissertation sur ce thème :
\begin{enumerate}
    \item \textbf{Rappeler le problème posé} (ex. : \og La diversité camerounaise est-elle un atout ou un handicap pour l'unité nationale ? \fg).
    \item \textbf{Synthétiser les idées clés} (2-3 phrases max) : \og La diversité physique fonde une complémentarité économique Nord-Sud/Ouest-Est. La diversité humaine alimente un brassage urbain majeur. L'État a mis en place des instruments (bilinguisme, décentralisation, devise) pour transformer cette diversité en unité. \fg
    \item \textbf{Ouvrir sur une perspective} (sans idée nouvelle) : \og Toutefois, la persistance des inégalités territoriales et des replis identitaires montre que l'unité reste un \emph{chantier permanent} nécessitant une gouvernance plus équitable et une éducation citoyenne renforcée. \fg
    \item \textbf{Interdiction :} Introduire un nouvel exemple, une nouvelle statistique ou un nouvel argument.
\end{enumerate}
\end{methode}

\courssub{La valorisation de la diversité : leviers économiques et culturels}

La diversité n'est pas seulement gérée (éviter les conflits), elle est \emph{valorisée} (créer de la valeur).

\begin{exemplebox}[Tourisme, terroir et patrimoine : la diversité comme produit]
\begin{itemize}
    \item \textbf{Tourisme écologique :} Parc national de Waza (sahélien), Parc de Korup (forêt primaire), Mont Cameroun (volcan actif), Lac Nyos (crater lake), Chutes de la Lobé (fleuve-mer). $\to$ Revenus de devises, emplois locaux, conservation.
    \item \textbf{Tourisme culturel / patrimonial :} 
    \begin{itemize}
        \item \textbf{Palais des Sultanats (Nord) :} Foumban (Bamoun), Rey Bouba, Mandara. Architecture soudanaise, musées, rites (Ngoun).
        \item \textbf{Cases obus (Nord-Ouest) :} Villages de la plaine de Ndop (Bamessing, Bamali). Architecture circulaire en terre, toits coniques, organisation sociale.
        \item \textbf{Chefferies Bamiléké (Ouest) :} Bafoussam, Bandjoun, Baham, Bangangté. Architecture à pilier central, cases à verandah, fresques, sociétés secrètes, arts de la cour (perles, ndop, masques).
        \item \textbf{Patrimoine immatériel :} Ngondo (Sawa), Medumba (Bamiléké), Nyem-Nyem (Adamaoua), danses, langues, gastronomies (ndolé, koki, fufu, kilichi, miel blanc de l'Adamaoua).
    \end{itemize}
    \item \textbf{Produits du terroir (IGP / Marque collective potentielles) :} Café Arabica (Boyo, Ouest), Cacao (Sud, Centre), Miel (Adamaoua), Piment (Ouest), Riz (Logone), Sel de Bamenjou, Huile de palme (Sanaga-Maritime).
\end{itemize}
\emph{Enjeu :} Passer de l'informel à la labellisation (Indication Géographique Protégée) pour capter la valeur ajoutée sur le territoire.
\end{exemplebox}

\begin{popfigure}
\begin{tikzpicture}[scale=0.9, every node/.style={font=\small, align=center}]
    % Nœuds principaux
    \node[draw=popblue, thick, fill=popblueL, rounded corners, minimum width=3.5cm, minimum height=1.2cm, align=center] (phys) at (0, 3){DIVERSITÉ PHYSIQUE\\Reliefs, Climats, Végétations, Sols};
    \node[draw=poporange, thick, fill=poporangeL, rounded corners, minimum width=3.5cm, minimum height=1.2cm, align=center] (hum) at (0, 0){DIVERSITÉ HUMAINE\\Groupes ethniques, Langues, Cultures, Savoirs-faire};
    
    % Flèches convergence
    \draw[popdark, thick, -{Stealth[length=3mm]}] (phys.south) -- node[right, xshift=2mm, align=center]{Potentialités\\complémentaires} (0, 1.8);
    \draw[popdark, thick, -{Stealth[length=3mm]}] (hum.north) -- node[right, xshift=2mm, align=center]{Brassage,\\Créativité} (0, 1.2);
    
    % Nœud central : Gestion / Valorisation
    \node[draw=popgreen, very thick, fill=popgreenL, rounded corners, minimum width=4.5cm, minimum height=1.5cm, align=center] (gest) at (0, -1.5){
        \textbf{GESTION POLITIQUE \& VALORISATION ÉCONOMIQUE}\\
        \textit{Bilinguisme -- Décentralisation -- Quotas -- Éducation civique}\\
        \textit{Tourisme -- Labellisation produits -- Industries culturelles}
    };
    
    % Flèches vers le nœud central
    \draw[popblue, thick, -{Stealth[length=3mm]}] (0, 1.8) -- (gest.north);
    \draw[poporange, thick, -{Stealth[length=3mm]}] (0, 1.2) -- (gest.north);
    
    % Résultat : Unité Nationale
    \node[draw=popgold, very thick, fill=popgoldL, rounded corners, minimum width=4cm, minimum height=1.2cm, align=center] (unite) at (0, -3.8){\textbf{UNITÉ NATIONALE CAMEROUNAISE}\\\og Paix -- Travail -- Patrie \fg};
    \draw[popgreen, very thick, -{Stealth[length=4mm]}] (gest.south) -- (unite.north);
    
    % Boucle de rétroaction (Unité protège Diversité)
    \draw[poppurple, dashed, thick, -{Stealth[length=3mm]}, bend left=40] (unite.west) to node[left, xshift=-3mm, align=center]{Protection\\Droits minorités\\Équité territoriale} (phys.west);
    \draw[poppurple, dashed, thick, -{Stealth[length=3mm]}, bend right=40] (unite.east) to node[right, xshift=3mm, align=center]{Cadre légal\\Porteur de projets} (hum.east);
\end{tikzpicture}
\legende{Schéma-bilan de synthèse : la diversité physique et humaine, gérée par des instruments politiques (bilinguisme, décentralisation) et valorisée économiquement (tourisme, terroir), produit l'unité nationale. L'unité (flèches violettes pointillées) rétroagit pour protéger la diversité (droits, équité).}
\end{popfigure}

\courssub{Conclusion : la diversité, socle de l'unité camerounaise}

Au terme de cette leçon, il apparaît que la diversité biogéographique et humaine du Cameroun n'est pas un accident de l'histoire ni un obstacle au développement. Elle est la \textbf{condition même de l'existence du pays} en tant qu'entité viable.

\begin{aretenir}[Synthèse finale : La diversité, socle de l'unité]
\begin{enumerate}
    \item \textbf{Fondement biophysique :} Le gradient climatique (sahélien $\to$ équatorial) et le relief (plaines, plateaux, volcans, montagne) créent une \textbf{complémentarité productive} (céréales/bétail au Nord ; plantation/forêt/pêche au Sud ; cultures de hauteur à l'Ouest).
    \item \textbf{Moteur démographique :} Les \textbf{migrations intérieures} (Nord$\to$Sud, Ouest$\to$Littoral/Centre) transforment la complémentarité potentielle en flux réels (biens, capitaux, compétences, gènes, cultures). Douala et Yaoundé en sont les creusets majeurs ($>20\%$ pop.).
    \item \textbf{Construction politique :} L'unité est un \textbf{projet conscient} : Devise (Paix-Travail-Patrie), Bilinguisme, Décentralisation, Brassage administratif/scolaire. Ces instruments transforment la \emph{coexistence} en \emph{citoyenneté partagée}.
    \item \textbf{Vigilance permanente :} Les défis (tribalisme, inégalités Nord/Sud, conflits ressources, marginalisation Pygmées) ne sont pas des fatalités mais des \textbf{indicateurs de dysfonctionnement de la gestion de la diversité}. Ils appellent plus d'équité, plus de droit, plus de dialogue.
    \item \textbf{Valorisation économique :} La diversité devient \textbf{atout comparatif} par le tourisme (écologique/culturel), la labellisation des produits du terroir et les industries créatives.
\end{enumerate}
\emph{En géographie comme en citoyenneté : l'unité camerounaise ne se décrète pas, elle se \textbf{construit chaque jour} dans l'équité territoriale et le respect mutuel.}
\end{aretenir}

\begin{exoresolu}[Sujet type Bac : \og La diversité au Cameroun est-elle un atout ou un frein pour l'unité nationale ? \fg]
\textbf{Analyse du sujet :} Question dialectique (atout / frein) sur le lien Diversité $\leftrightarrow$ Unité. Échelle : nationale. Concepts clés : complémentarité, brassage, inégalités, instruments politiques.
\par
\textbf{Plan détaillé :}
\textbf{Introduction :} Accroche (Cameroun = \og Afrique en miniature \fg). Définition de la diversité (physique/humaine). Problématique : La diversité structure le territoire et la société ; elle est un atout majeur \emph{si} elle est gérée par des institutions équitables, sinon elle devient source de fracture. Annonce du plan.
\par
\textbf{I. La diversité, fondement matériel et humain de l'unité (L'ATOUT).}
\begin{itemize}
    \item A. Complémentarité biophysique et échanges Nord/Sud, Ouest/Est (céréales vs plantation, bétail vs poisson). Chiffres : flux commerciaux, sécurité alimentaire.
    \item B. Brassage humain : migrations internes, métissage urbain (Douala/Yaoundé $>20\%$ pop.), camfranglais, mariages mixtes. Création d'une \og culture nationale \fg par le bas.
    \item C. Valorisation économique : tourisme (Waza, Korup, Sultanats, Chefferies), produits terroir (café, miel, piment), potentiel énergétique (barrages), main-d'œuvre qualifiée diverse.
\end{itemize}
\par
\textbf{II. Les tensions nées d'une gestion inégalitaire de la diversité (LE FREIN POTENTIEL).}
\begin{itemize}
    \item A. Inégalités territoriales structurelles (Indicateurs HDI, pauvreté : Septentrion/Est vs Littoral/Centre/Ouest). Sentiment d'abandon $\to$ repli identitaire.
    \item B. Conflits d'usage de l'espace : Agriculteurs/Éleveurs (couloirs transhumance), Autochtones/Allogènes (foncier urbain/périurbain), Minorités forestières (Pygmées) vs Exploitation forestière.
    \item C. Instrumentalisation politique : Tribalisme, népotisme, discours de haine, crise anglophone (NOSO) $\to$ remise en cause du pacte national.
\end{itemize}
\par
\textbf{III. L'unité comme construction politique permanente : transformer le frein en atout.}
\begin{itemize}
    \item A. Cadre institutionnel : Constitution 1996 (Décentralisation, Bilinguisme, Commission Bilingue/Multiculturalisme), Devise, Symboles.
    \item B. Politiques d'équité : Discrimination positive (concours), Investissements infrastructures (autoroute Ydé-Ngaoundéré, Port Kribi, Barrages), PNDP.
    \item C. Éducation et citoyenneté : Brassage scolaire, Enseignement morale/civique, Médias publics bilingues, Société civile (ONG, Églises, Chefferies modernisées).
\end{itemize}
\par
\textbf{Conclusion :} La diversité est \textbf{structurellement un atout} (ressources, créativité, résilience). Elle devient un frein \textbf{conjoncturellement} par défaut de gouvernance (inégalités, injustice). L'unité nationale n'est pas l'absence de diversité, mais sa \textbf{régulation démocratique et équitable}. Perspective : Le Cameroun de 2035 (Vision 2035) ne sera émergent que s'il fait de sa diversité son \textbf{premier capital}.
\tcblower
\textbf{Éléments de correction / Grille d'évaluation :}
\begin{itemize}
    \item \textbf{Connaissances (6 pts) :} Exemples précis des 3 zones biogéo, flux migratoires nommés, instruments juridiques (Constitution, lois), exemples conflits (NOSO, transhumance), exemples valorisation (sites, produits).
    \item \textbf{Analyse / Argumentation (8 pts) :} Structure dialectique respectée (Atout / Frein / Dépassement). Liens logiques (diversité $\to$ complémentarité $\to$ échange $\to$ interdépendance). Nuance : \og frein \fg = mauvaise gestion, pas diversité en soi. Utilisation vocabulaire : brassage, complémentarité, disparités, décentralisation, vivre-ensemble.
    \item \textbf{Expression / Présentation (6 pts) :} Introduction/Conclusion soignées. Paragraphes équilibrés. Cartes/croquis intégrés (carte flux, carte inégalités). Propreté, syntaxe.
\end{itemize}
\end{exoresolu}

\begin{exercice}[Entraînement : Croquis commenté]
À partir du schéma-bilan de synthèse (figure ci-dessus) et de vos connaissances, réalisez un \textbf{croquis commenté} (format A4 paysage) intitulé : \og \textbf{La diversité camerounaise : fondements, tensions et régulations de l'unité nationale} \fg.
\par
\textbf{Consignes :}
\begin{enumerate}
    \item Tracer un contour simplifié du Cameroun (repérer : Yaoundé, Douala, Ngaoundéré, Maroua, Bafoussam, Bertoua, Foumban, Waza, Korup, Mont Cameroun).
    \item Représenter par des \textbf{flèches de largeurs proportionnelles} : Flux agricoles Nord$\to$Sud (céréales, bétail) ; Flux plantation Sud$\to$Nord (cacao, bois, plantain) ; Migrations Ouest$\to$Littoral/Centre ; Migrations Nord$\to$Villes du Sud.
    \item Localiser par \textbf{symboles ponctuels} : 2 zones de conflits agriculteurs-éleveurs ; 2 zones de marginalisation Pygmées ; 3 sites touristiques majeurs (écologique + culturel) ; 2 grands chantiers d'équipement structurant (ex. autoroute, barrage, port).
    \item Légende organisée en 3 items : \og 1. Complémentarités et flux \fg, \og 2. Tensions et fragilités \fg, \og 3. Instruments d'unité et valorisation \fg.
    \item Rédiger un \textbf{commentaire de 15 lignes} structuré en 3 paragraphes (Atouts, Défis, Réponses) mobilisant la méthode de rédaction de conclusion.
\end{enumerate}
\end{exercice}

\begin{corrige}[Exercice n°1 : Croquis commenté]
\textbf{Attendus du croquis :}
\begin{itemize}
    \item \textbf{Base cartographique :} Contour reconnaissable (golfe de Guinée au Sud, pointe extrême-Nord au Lac Tchad). Villes bien positionnées (Yaoundé centre-sud, Douala littoral, Ngaoundéré centre, Maroua extrême-Nord, Bafoussam Ouest, Bertoua Est).
    \item \textbf{Flux (Flèches) :} 
    \begin{itemize}
        \item Flèche épaisse \textcolor{poporange}{orange} Nord (Maroua/Garoua) $\to$ Sud (Douala/Yaoundé) : \og Céréales, Bétail, Oignons \fg.
        \item Flèche épaisse \textcolor{popgreen}{verte} Sud (Douala/Yaoundé/Bertoua) $\to$ Nord : \og Produits manufacturés, Carburants, Plantain, Cacao (export) \fg.
        \item Flèches \textcolor{popblue}{bleues} Ouest (Bafoussam/Bamenda) $\leftrightarrow$ Douala/Yaoundé : \og Migrations travail/études, Transferts fonds \fg.
    \end{itemize}
    \item \textbf{Symboles ponctuels :} 
    \begin{itemize}
        \item $\otimes$ Rouge (Adamaoua/Nord-Ouest) : Conflits agro-pastoraux.
        \item $\triangle$ Violet (Sud-Est/Océan) : Zones Pygmées (Baka/Bakola).
        \item $\star$ Or (Waza, Korup, Mont Cameroun, Foumban, Bafoussam) : Tourisme.
        \item $\blacksquare$ Noir (Ngaoundéré, Kribi, Nachtigal) : Infrastructures structurantes.
    \end{itemize}
    \item \textbf{Légende :} Propre, hiérarchisée (Surface/Ligne/Point), couleurs cohérentes.
\end{itemize}
\textbf{Attendus du commentaire (exemple) :}
\og Le croquis met en évidence la structuration du territoire camerounais par la complémentarité biophysique. Les flèches orange et vertes illustrent l'interdépendance Nord-Sud : le Nord pourvoit le Sud en protéines animales et céréales, le Sud approvisionne le Nord en biens manufacturés et assure l'exportation des produits de rente. Cette complémentarité est activée par les migrations (flèches bleues), véritables ciment humain. Cependant, des tensions persistent : les conflits agro-pastoraux (symboles $\otimes$) révèlent la pression foncière et l'absence de régulation foncière efficace ; la marginalisation des Pygmées ($\triangle$) signale l'injustice sociale. Face à cela, l'État déploie des instruments d'unité : infrastructures transversales (autoroute, barrage, port) reliant les pôles, et valorisation des patrimoines ($\star$) comme leviers de développement local. L'unité nationale se construit ainsi par l'équité territoriale et la reconnaissance de la diversité. \fg
\end{corrige}

\newpage

\coursec{Activité d'intégration}

\coursec{Leçon 01 : La diversité des ensembles biogéographiques du Cameroun}

\courssub{Objectifs de l'activité}
\begin{itemize}
    \item Construire une carte-croquis de synthèse localisant les quatre grands ensembles biogéographiques du Cameroun et la répartition des grands groupes humains.
    \item Réaliser un diagramme ombrothermique comparatif pour deux stations climatiques types (Yaoundé au Sud, Maroua au Nord).
    \item Interpréter les données physiques et humaines pour expliquer la relation entre diversité des milieux et unité nationale.
    \item Mobiliser les savoir-faire cartographiques : légende organisée, figurés adaptés (surface, linéaire, ponctuel), lecture de tableaux statistiques.
\end{itemize}

\courssub{Situation déclenchante}
Le Cameroun est souvent qualifié d'\og Afrique en miniature \fg{} en raison de la concentration, sur son territoire (\SI{475\,442}{km^2}), de la quasi-totalité des climats, végétations, reliefs et groupes ethniques que l'on rencontre sur le continent africain. Cette diversité n'est pas une juxtaposition hasardeuse : elle structure l'espace national en grands ensembles biogéographiques cohérents, chacun associant un type de relief, un climat, une végétation, des sols et des populations aux activités adaptées.

En tant que géographes en formation, vous êtes missionnés par le Ministère de l'Éducation de Base (MINEDUB) pour produire un \emph{atlas simplifié à l'usage des classes de CM2} illustrant cette maxime. Votre livrable principal est une \textbf{carte-croquis de synthèse au 1/10 000 000} (1 cm = 100 km) accompagnée d'un diagramme ombrothermique comparatif et d'une conclusion argumentée de 10 lignes maximum.

\medskip
\textbf{Dossier documentaire fourni (extraits) :}

\begin{center}
\begin{tabularx}{\linewidth}{|l|X|X|X|X|}\hline
\rowcolor{popblueL}
\textbf{Région (Chef-lieu)} & \textbf{Ensemble biogéographique} & \textbf{Climat type} & \textbf{Pluviométrie annuelle (mm)} & \textbf{Groupes humains majoritaires} \\ \hline
Extrême-Nord (Maroua) & 1. Domaine soudanais & Soudanais (Aw) & 700--900 & Peuls, Mafa, Kotoko, Arabes Choa \\ \hline
Nord (Garoua) & 1. Domaine soudanais & Soudanais (Aw) & 900--1\,200 & Peuls, Mboum, Gbaya \\ \hline
Adamaoua (Ngaoundéré) & 2. Domaine des hautes terres & Montagnard (Cwb) & 1\,500--2\,000 & Peuls, Mboum, Tikar, Gbaya \\ \hline
Est (Bertoua) & 3. Domaine forestier équatorial & Équatorial guinéen (Am) & 1\,500--2\,000 & Baya, Bakas (Pygmées), Kaka \\ \hline
Sud (Ebolowa) & 3. Domaine forestier équatorial & Équatorial guinéen (Am) & 1\,500--2\,000 & Boulou, Fang, Bassa \\ \hline
Centre (Yaoundé) & 3. Domaine forestier équatorial & Équatorial guinéen (Am) & 1\,600 & Ewondo, Bassa, Bamiléké (migrants) \\ \hline
Littoral (Douala) & 4. Domaine côtier & Équatorial camerounais (Am) & 3\,000--4\,000 & Duala, Bassa, Bamiléké, Haoussas \\ \hline
Sud-Ouest (Buea) & 4. Domaine côtier & Équatorial de montagne & 3\,000--10\,000 & Bakweri, Bamiléké, Anglophones \\ \hline
Nord-Ouest (Bamenda) & 2. Domaine des hautes terres & Montagnard (Cwb) & 2\,000--3\,000 & Grassfields (Bamiléké, Bamoun, Tikar) \\ \hline
Ouest (Bafoussam) & 2. Domaine des hautes terres & Montagnard (Cwb) & 1\,800--2\,500 & Bamiléké, Bamoun \\ \hline
\end{tabularx}
\end{center}

\medskip
\textbf{Données climatiques mensuelles pour le diagramme ombrothermique (normales 1981-2010) :}

\begin{center}
\begin{tabular}{|c|cccccccccccc|}\hline
\rowcolor{popblueL}
\textbf{Mois} & \textbf{J} & \textbf{F} & \textbf{M} & \textbf{A} & \textbf{M} & \textbf{J} & \textbf{J} & \textbf{A} & \textbf{S} & \textbf{O} & \textbf{N} & \textbf{D} \\ \hline
\textbf{Yaoundé - T° moy (°C)} & 24.5 & 25.0 & 25.0 & 24.5 & 24.0 & 23.5 & 23.0 & 23.0 & 23.5 & 23.5 & 24.0 & 24.0 \\ \hline
\textbf{Yaoundé - Pluie (mm)} & 18 & 38 & 115 & 165 & 195 & 135 & 70 & 110 & 225 & 280 & 140 & 30 \\ \hline
\textbf{Maroua - T° moy (°C)} & 26.5 & 29.0 & 32.0 & 33.5 & 32.0 & 29.5 & 27.5 & 26.5 & 27.0 & 28.5 & 27.5 & 26.5 \\ \hline
\textbf{Maroua - Pluie (mm)} & 0 & 0 & 5 & 30 & 90 & 130 & 190 & 230 & 160 & 40 & 2 & 0 \\ \hline
\end{tabular}
\end{center}

\medskip
\textbf{Rappel des limites des ensembles biogéographiques (simplifiées pour le 1/10M) :}
\begin{itemize}
    \item \textbf{Domaine soudanais} : Au nord du 7\textsuperscript{e} parallèle (environ), limite sud marquée par l'isohyète 1\,000 mm et la frontière Adamaoua.
    \item \textbf{Domaine des hautes terres} : Massif de l'Adamaoua (au nord), Plateau de l'Ouest et Monts Cameroun/Bamenda (au sud-ouest), altitudes > 1\,000 m.
    \item \textbf{Domaine forestier équatorial} : Au sud du 6\textsuperscript{e} parallèle (environ), hors zone côtière et hautes terres. Forêt dense humide.
    \item \textbf{Domaine côtier} : Bande littorale atlantique (0--50 km), plaines de la Sanaga, du Moungo, du Wouri, piémont du Mont Cameroun. Mangroves, forêt littorale, sols hydromorphes.
\end{itemize}

\courssub{Consignes du travail}
\begin{enumerate}
    \item \textbf{Fond de carte} : Sur une feuille A3 (ou format fourni), tracer le contour du Cameroun au 1/10 000 000 (dimensions utiles : ~8 cm Est-Ouest $\times$ 12 cm Nord-Sud). Positionner le Méridien de Greenwich (traversant Kribi) et l'Équateur (au sud de Campo). Repérer et nommer : Yaoundé, Douala, Garoua, Maroua, Ngaoundéré, Bafoussam, Bamenda, Bertoua, Ebolowa, Buea. Tracer les principaux axes routiers (N1, N2, N3, N6, N10) et le fleuve Bénoué.
    \item \textbf{Ensembles biogéographiques} : Délimiter les 4 ensembles à l'aide de traits de limites distincts (pointillés, tirets, traits mixtes). Les colorier avec 4 teintes distinctes (ex: jaune/ocre pour soudanais, vert foncé pour forestier, vert clair/brun pour hautes terres, bleu clair/vert jaune pour côtier). Construire une \textbf{légende organisée} en 4 items (Relief/Climat/Végétation/Sols/Activités principales).
    \item \textbf{Groupes humains} : Superposer aux ensembles biogéographiques la répartition des \textbf{grands groupes linguistiques/ethniques} (liste fournie dans le tableau) à l'aide de \textbf{figurés de surface (hachures orientées différentes)} ou de \textbf{figurés ponctuels (symboles ethniques)}. Ajouter ces figurés dans la légende.
    \item \textbf{Diagramme ombrothermique} : Sur un graphique unique (abscisses : mois J à D ; ordonnée gauche : T° en °C, échelle 1 cm = 5°C ; ordonnée droite : Pluie en mm, échelle 1 cm = 50 mm pour Yaoundé / 100 mm pour Maroua \emph{ou double échelle}), tracer :
        \begin{itemize}
            \item La courbe des températures (trait \texttt{poppink} continu pour Yaoundé, trait \texttt{poporange} pointillé pour Maroua).
            \item Les histogrammes des précipitations (barres \texttt{popblue} pleines Yaoundé, barres \texttt{poporange} hachurées Maroua, décalées latéralement).
        \end{itemize}
        Identifier visuellement la saison sèche et la saison des pluies pour chaque station. Calculer l'amplitude thermique annuelle et le total pluviométrique annuel pour les deux stations.
    \item \textbf{Rédaction} : En \textbf{10 lignes maximum}, répondre à la question : \og En quoi la diversité des milieux et des hommes fonde-t-elle l'unité du Cameroun ? \fg{} S'appuyer sur le croquis et le diagramme (complémentarité économique, brassage, unité politique).
\end{enumerate}

\courssub{Ressources et méthodes à mobiliser}

\begin{methode}[Réaliser un croquis de synthèse géographique (échelle 1/10M)]
\begin{enumerate}
    \item \textbf{Choix de l'échelle et du format} : 1 cm = 100 km. Format A3 (42 x 29,7 cm) ou A4 paysage. Calculer les dimensions du contour : Larg. $\approx$ 8 cm, Haut. $\approx$ 12 cm.
    \item \textbf{Repérage absolu et relatif} : Tracer le Méridien de Greenwich (Kribi) et l'Équateur (sud Campo). Positionner les 10 chefs-lieux de région par coordonnées approximatives (voir corrigé).
    \item \textbf{Structuration de l'espace (linéaires)} : Frontières internationales, axes routiers structurants (N1, N2, N3, N6, N10), hydrographie majeure (Bénoué, Sanaga, Wouri, Logone).
    \item \textbf{Surfaces (choroplèthe / hachures)} : Délimiter les 4 ensembles biogéographiques (couleurs de fond). Superposer les aires linguistiques (hachures orientées).
    \item \textbf{Points (ponctuels)} : Villes, ressources minières, barrages, parcs nationaux.
    \item \textbf{Légende organisée et hiérarchisée} : Titre, auteur, date, échelle. Ordre : Linéaires $\rightarrow$ Surfaces (Milieux) $\rightarrow$ Surfaces (Humains) $\rightarrow$ Ponctuels. Figurés conformes à la carte.
\end{enumerate}
\end{methode}

\begin{methode}[Construire et lire un diagramme ombrothermique (Gaussen)]
\begin{enumerate}
    \item \textbf{Axes} : Abscisses = 12 mois (J à D). Ordonnée gauche = Température (°C), échelle 1 cm = 5°C (0 à 35°C). Ordonnée droite = Précipitations (mm), échelle 1 cm = 50 mm (0 à 300/400 mm).
    \item \textbf{Convention Gaussen} : Relation $P = 2T$ (ou $P=3T$). Si barre de pluie \textbf{sous} la courbe T° $\rightarrow$ saison sèche. Si barre \textbf{au-dessus} $\rightarrow$ saison humide.
    \item \textbf{Calculs obligatoires} : Amplitude thermique annuelle = T° max -- T° min. Total pluviométrique annuel = $\sum$ pluies mensuelles.
    \item \textbf{Interprétation} : Type de climat (Aw, Am, Cwb...), nombre de saisons, durée saison sèche, adéquation pluviométrie/végétation.
\end{enumerate}
\end{methode}

\begin{attention}[Pièges fréquents au Bac Technique]
\begin{itemize}
    \item \textbf{Échelle} : Confondre 1/10M (1 cm = 100 km) et 1/5M (1 cm = 50 km). Vérifier les conversions (Yaoundé-Douala = 2.5cm, pas 5 cm).
    \item \textbf{Légende} : Figurés non conformes à la carte (ex: hachures sur la carte, couleur unie en légende). Oublier l'échelle ou l'orientation (flèche Nord).
    \item \textbf{Diagramme} : Superposer exactement les barres des deux stations (illisible) $\rightarrow$ utiliser \texttt{bar shift}. Inverser les échelles T°/Pluie. Oublier les calculs (amplitude, total).
    \item \textbf{Vocabulaire} : Dire \og climat équatorial \fg{} pour le Nord (c'est soudanais Aw). Confondre \og population \fg{} et \og groupe ethnique/linguistique \fg{}.
\end{itemize}
\end{attention}

\courssub{Démarche et corrigé commenté}

\begin{exoresolu}{TP : Croquis des ensembles biogéographiques et groupes humains du Cameroun}
\textbf{Étape 1 : Fond de carte et repérage}
\begin{itemize}
    \item Le contour du Cameroun est un triangle pointe au Nord (Lac Tchad) et base au Sud (Atlantique). Largeur max Est-Ouest $\approx$ 800 km (8 cm), Hauteur Nord-Sud $\approx$ 1\,200 km (12 cm).
    \item \textbf{Villes repères} (coordonnées approximatives en cm depuis coin Sud-Ouest) :
    \begin{itemize}
        \item Maroua (Extrême-Nord) : (6,5 ; 10,5)
        \item Garoua (Nord) : (5,0 ; 9,0)
        \item Ngaoundéré (Adamaoua) : (4,5 ; 6,5)
        \item Yaoundé (Centre) : (3,5 ; 3,5)
        \item Douala (Littoral) : (1,5 ; 3,0)
        \item Bafoussam (Ouest) : (2,5 ; 4,5)
        \item Bamenda (Nord-Ouest) : (2,0 ; 6,0)
        \item Bertoua (Est) : (6,5 ; 3,5)
        \item Ebolowa (Sud) : (4,0 ; 1,5)
        \item Buea (Sud-Ouest) : (1,0 ; 2,5)
    \end{itemize}
    \item Axes routiers : N1 (Yaoundé-Douala), N2 (Yaoundé-Garoua-Maroua), N3 (Yaoundé-Bertoua), N6 (Bafoussam-Bamenda), N10 (Bafoussam-Foumban-Ngaoundéré). Fleuve Bénoué (Garoua $\rightarrow$ Nigeria). Fleuve Sanaga (centre).
\end{itemize}

\textbf{Étape 2 : Délimitation et coloriage des 4 ensembles biogéographiques}
\begin{center}
\begin{tabularx}{\linewidth}{|l|X|X|X|}\hline
\rowcolor{popblueL}
\textbf{Ensemble} & \textbf{Coloris suggéré} & \textbf{Figuré de limite} & \textbf{Contenu légende (synthèse)} \\ \hline
1. Domaine soudanais & Jaune ocre / Beige & Trait continu fin & \textbf{Relief :} Plaine du Bénoué, piémonts Adamaoua, bassin du Lac Tchad (alt. < 500 m). \textbf{Climat :} Soudanais (Aw), 1 saison sèche (nov-mars), 1 saison des pluies (avr-oct). \textbf{Végétation :} Savane arborée, herbeuse, galerie forestière. \textbf{Sols :} Ferrugineux tropicaux lessivés, hydromorphes (bas-fonds). \textbf{Activités :} Coton, mil/sorgho, maïs, élevage bovin (Peuls), pêche (Logone). \\ \hline
2. Domaine des hautes terres & Vert clair / Brun clair & Trait pointillé épais & \textbf{Relief :} Massif de l'Adamaoua (1\,100--2\,500 m), Plateau de l'Ouest (1\,000--2\,000 m), Mont Cameroun (4\,095 m), Monts Bamenda. \textbf{Climat :} Montagnard tempéré (Cwb), frais, pluies abondantes. \textbf{Végétation :} Savane d'altitude, forêt de galerie, prairies (Monts). \textbf{Sols :} Ferralitiques lessivés, volcans (andosols) fertiles. \textbf{Activités :} Café arabica, thé, pomme de terre, maraîchage, élevage laitier, tourisme. \\ \hline
3. Domaine forestier équatorial & Vert foncé & Trait mixte (-- --) & \textbf{Relief :} Plateau Sud-Camerounais (500--800 m), cuvettes. \textbf{Climat :} Équatorial guinéen (Am), 4 saisons, transitions courtes. \textbf{Végétation :} Forêt dense humide sempervirente (bois d'œuvre : sapelli, ayous, iroko). \textbf{Sols :} Ferralitiques jaunes/rouges, lessivés, pauvres en bases. \textbf{Activités :} Exploitation forestière, cacao, café robusta, huile de palme, chasse/cueillette (Bakas). \\ \hline
4. Domaine côtier & Bleu clair / Vert jaune & Trait de points rapprochés & \textbf{Relief :} Plaines alluviales (Sanaga, Wouri, Moungo), piémont Mont Cameroun, mangroves. \textbf{Climat :} Équatorial camerounais (Am), hyperhumide (Debundscha > 10\,000 mm). \textbf{Végétation :} Forêt littorale, mangroves, palmiers à huile. \textbf{Sols :} Hydromorphes, alluvionnaires, volcaniques (pentes Mont Cameroun). \textbf{Activités :} Port (Douala/Kribi), agro-industrie (CDC, Pamol : banane, palme, caoutchouc), pêche, pétrole offshore, industrie. \\ \hline
\end{tabularx}
\end{center}

\textbf{Étape 3 : Superposition des grands groupes humains}
Utiliser des \textbf{hachures orientées} différentes pour chaque grand groupe linguistique (superposées aux couleurs biogéographiques) :
\begin{itemize}
    \item \textbf{Soudano-sahéliens (Peuls, Arabes Choa, Kotoko, Mafa)} : Hachures verticales \texttt{poppink} (Nord, Extrême-Nord, Adamaoua). \emph{Note : Peuls présents aussi en transhumance dans les hautes terres et au Sud.}
    \item \textbf{Adamawa-Oubanguiens (Mboum, Gbaya, Baya, Tikar)} : Hachures diagonales \og / \fg{} \texttt{poporange} (Adamaoua, Nord, Est).
    \item \textbf{Bantous des forêts (Beti-Fang : Ewondo, Boulou, Fang ; Bassa ; Douala ; Kaka)} : Hachures diagonales \og \textbackslash \fg{} \texttt{popgreen} (Centre, Sud, Est, Littoral).
    \item \textbf{Semi-Bantous / Grassfields (Bamiléké, Bamoun, Tikar, Bamenda)} : Hachures horizontales \texttt{popblue} (Ouest, Nord-Ouest, + diaspora forte au Centre, Littoral, Sud-Ouest).
    \item \textbf{Peuples côtiers / SW (Bakweri, Duala, Isubu, Anglophones)} : Points noirs denses (Sud-Ouest, Littoral).
    \item \textbf{Pygmées (Bakas/Bagyeli)} : Petits triangles noirs épars (Est, Sud).
\end{itemize}
\emph{Astuce légende} : Grouper sous \og Groupes linguistiques majeurs \fg{} avec 6 figurés de surface (carrés 3mm hachurés) + 1 ponctuel (triangle).

\textbf{Étape 4 : Diagramme ombrothermique comparatif (Yaoundé vs Maroua)}
\begin{center}
\begin{tikzpicture}[scale=0.75]
\begin{axis}[
    width=\linewidth, height=10cm,
    axis y line*=left, axis x line=bottom,
    xmin=0.5, xmax=12.5,
    xtick={1,2,3,4,5,6,7,8,9,10,11,12},
    xticklabels={J,F,M,A,M,J,J,A,S,O,N,D},
    ymin=0, ymax=35,
    ylabel={Température (°C)}, ylabel style={poppink},
    ytick={0,5,10,15,20,25,30,35},
    legend pos=north west,
    grid=major, grid style={popline},
]
% Yaoundé Temp
\addplot[popcurve, poppink, thick, mark=*] coordinates {
(1,24.5)(2,25.0)(3,25.0)(4,24.5)(5,24.0)(6,23.5)(7,23.0)(8,23.0)(9,23.5)(10,23.5)(11,24.0)(12,24.0)
};
\addlegendentry{Yaoundé T°}
% Maroua Temp
\addplot[popcurve, poporange, dashed, thick, mark=square*] coordinates {
(1,26.5)(2,29.0)(3,32.0)(4,33.5)(5,32.0)(6,29.5)(7,27.5)(8,26.5)(9,27.0)(10,28.5)(11,27.5)(12,26.5)
};
\addlegendentry{Maroua T°}
\end{axis}

\begin{axis}[
    width=\linewidth, height=10cm,
    axis y line*=right, axis x line=none,
    xmin=0.5, xmax=12.5,
    ymin=0, ymax=400,
    ylabel={Précipitations (mm)}, ylabel style={popblue},
    ytick={0,100,200,300,400},
    ybar, bar width=0.25cm,
    legend pos=south east,
]
% Yaoundé Pluie (bar shift gauche)
\addplot[popblue, fill=popblue!60, bar shift=-0.15cm] coordinates {
(1,18)(2,38)(3,115)(4,165)(5,195)(6,135)(7,70)(8,110)(9,225)(10,280)(11,140)(12,30)
};
\addlegendentry{Yaoundé Pluie}
% Maroua Pluie (bar shift droite, hachuré)
\addplot[poporange, fill=poporange!60, pattern=north east lines, pattern color=poporange!80!black, bar shift=0.15cm] coordinates {
(1,0)(2,0)(3,5)(4,30)(5,90)(6,130)(7,190)(8,230)(9,160)(10,40)(11,2)(12,0)
};
\addlegendentry{Maroua Pluie}
\end{axis}
\end{tikzpicture}
\legende{Diagramme ombrothermique comparatif Yaoundé (Sud, climat Am) et Maroua (Nord, climat Aw). Courbes T° : \texttt{poppink} continue (Yaoundé), \texttt{poporange} pointillée (Maroua). Histogrammes Pluie : barres \texttt{popblue} pleines (Yaoundé), barres \texttt{poporange} hachurées (Maroua), décalées pour lisibilité. Échelle T° : 1 cm = 5°C. Échelle Pluie : 1 cm = 50 mm (lecture Yaoundé) / 100 mm (lecture Maroua approximative).}
\end{center}

\textbf{Analyse du diagramme (Calculs) :}
\begin{itemize}
    \item \textbf{Yaoundé} : Amplitude thermique = $25,0 - 23,0 = \mathbf{2,0\,^\circ C}$ (faible, typique équatorial). Total pluviométrique $\approx \mathbf{1\,521\,mm}$. Saison sèche courte (déc-janv, juil-août \og petite saison sèche \fg{}). $P > 2T$ presque toute l'année $\rightarrow$ forêt dense.
    \item \textbf{Maroua} : Amplitude thermique = $33,5 - 26,5 = \mathbf{7,0\,^\circ C}$ (marquée, continentalité). Total pluviométrique $\approx \mathbf{877\,mm}$. Saison sèche longue (nov-mai, $P < 2T$). Saison des pluies concentrée (juin-sept). $P < 1000$ mm $\rightarrow$ savane soudanais.
\end{itemize}

\textbf{Étape 5 : Rédaction de la conclusion (Modèle corrigé)}
\og La diversité des milieux (4 ensembles biogéographiques aux climats, sols, ressources contrastés) structure une \textbf{complémentarité économique} forte : le Nord soudanais (coton, bétail, céréales) alimente le Sud forestier et côtier (bois, cacao, café, pétrole, industrie, port de Douala/Kribi). Les hautes terres (café arabica, maraîchage, main-d'œuvre) font le lien. Cette complémentarité impose des \textbf{flux Nord-Sud} (bétail, céréales, migrants Bamiléké/Peuls) et \textbf{Sud-Nord} (produits manufacturés, carburants, revenus), soudant le territoire. Le \textbf{brassage ethnique} (diasporas Bamiléké, Peuls, Haoussas, fonctionnaires) et l'\textbf{unité administrative} (10 régions, bilinguisme officiel, capitale politique Yaoundé / économique Douala) transforment cette diversité en ciment national. L'\og Afrique en miniature \fg{} n'est pas une fragilité mais un \textbf{atout structurel} si l'aménagement du territoire (routes, rail, énergie) réduit l'enclavement du Nord et de l'Est. \fg{} (9 lignes).
\end{exoresolu}

\courssub{Bilan de l'activité}
Cette activité d'intégration a permis de mettre en évidence :
\begin{enumerate}
    \item \textbf{La structuration de l'espace camerounais en 4 ensembles biogéographiques cohérents}, chacun associant un triplet Relief-Climat-Végétation/Sols déterminant des potentialités agricoles et minières spécifiques.
    \item \textbf{L'interdépendance fonctionnelle} entre ces ensembles : le gradient climatique Nord-Sud (aridité $\rightarrow$ humidité) et le gradient altimétrique (plaines $\rightarrow$ hautes terres $\rightarrow$ volcan) créent une complémentarité de productions (coton/bétail vs cacao/bois/pétrole vs café arabica/maraîchage) qui fonde l'unité économique.
    \item \textbf{La superposition non figée des groupes humains} sur les milieux : les grands ensembles linguistiques (Soudano-sahéliens, Adamawa-Oubanguiens, Bantous, Semi-Bantous) correspondent globalement aux grands domaines, mais les migrations séculaires (Peuls, Bamiléké) et récentes (fonctionnaires, commerçants) créent un \textbf{mosaïque ethnique} gérée par l'État unitaire et le bilinguisme.
    \item \textbf{La maîtrise technique} : construction d'un croquis géographique de synthèse (choix d'échelle, figurés, légende hiérarchisée) et d'un diagramme ombrothermique comparatif (double échelle, lecture Gaussen, calculs d'amplitude et totaux), outils indispensables à l'analyse géographique au Bac Technique.
\end{enumerate}

\begin{aretenir}[Synthèse pour le Bac / Probatoire]
Le Cameroun est une \textbf{unité dans la diversité}. Ses 4 ensembles biogéographiques (Soudanais, Hautes terres, Forestier équatorial, Côtier) offrent des ressources complémentaires. Les 250+ groupes ethniques, regroupés en 5 grandes familles linguistiques, s'y répartissent en interaction avec les milieux. L'unité nationale repose sur la \textbf{complémentarité économique Nord-Sud}, le \textbf{brassage démographique} et l'\textbf{organisation politique centralisée (État unitaire, 10 régions, bilinguisme)}. L'enjeu actuel est l'\textbf{aménagement du territoire} (désenclavement Est/Nord, corridor Douala-Ndjamena, barrage de Nachtigal) pour transformer la diversité potentielle en développement partagé.
\end{aretenir}

\newpage

\coursec{Méthodes et exercices résolus}

\coursec{La diversité des ensembles biogéographiques}

\courssub{Méthodes et exercices résolus}

\begin{methode}[Identifier et localiser les grands ensembles biogéographiques sur un croquis]
\textbf{Objectif :} Réaliser un croquis de synthèse localisant les 4 grands ensembles biogéographiques du Cameroun.
\begin{enumerate}
    \item \textbf{Repérer les limites physiques :} Tracer l'isohyète 1500 mm (frontière forêt/savane) et la ligne de partage des eaux (massif de l'Adamaoua, ~1000-1500 m).
    \item \textbf{Délimiter les 4 ensembles :}
    \begin{itemize}
        \item \textbf{Domaine équatorial (Sud) :} Au sud du 5\textsuperscript{e} parallèle environ, pluviométrie > 1500 mm, forêt dense.
        \item \textbf{Domaine tropical de transition (Centre/Adamaoua) :} Zone de contact, savanes boisées, gallery forests, altitude 1000-1500 m.
        \item \textbf{Domaine soudano-sahélien (Nord - Bénoué) :} Plaine de la Bénoué, pluviométrie 1000-1500 mm, savane herbeuse/arborée.
        \item \textbf{Domaine sahélien (Extrême-Nord) :} Au nord du 11\textsuperscript{e} parallèle, pluviométrie < 1000 mm, steppe, éboulis, lac Tchad.
    \end{itemize}
    \item \textbf{Figurer les éléments clés :} Villes repères (Yaoundé, Garoua, Maroua), cours d'eau (Sanaga, Bénoué, Logone), axes routiers structurants (N1, N10).
    \item \textbf{Légende organisée :} Utiliser des couleurs conventionnelles (vert foncé = forêt, vert clair/jaune = savane, ocre = sahélien, marron = montagne). Titre, échelle, orientation (Nord).
\end{enumerate}
\cle{Repères} : Parallèle 5° N (limite forêt/savane), Massif de l'Adamaoua (charnière climatique), Lac Tchad (exutoire nord).
\end{methode}

\begin{exoresolu}[Croquis de synthèse : Les ensembles biogéographiques du Cameroun]
\textbf{Énoncé :} À l'aide d'un croquis commenté (format A4), montrez la diversité des ensembles biogéographiques du Cameroun. La légende doit faire apparaître les quatre grands domaines, les grands axes de communication et deux villes par domaine.
\tcblower
\textbf{Solution détaillée :}

\textbf{1. Construction du croquis (TikZ simplifié)}

\begin{popfigure}
\begin{tikzpicture}[scale=0.8]
% Contour très simplifié du Cameroun (polygone approximatif)
\draw[popdark, thick] (0,0) -- (10,0) -- (11,2) -- (11,10) -- (9,12) -- (4,12) -- (2,10) -- (0,8) -- cycle;
\clip (0,0) -- (10,0) -- (11,2) -- (11,10) -- (9,12) -- (4,12) -- (2,10) -- (0,8) -- cycle;
% Domaine Équatorial (Sud) : y < 1.5
\fill[popgreen!60] (0,0) rectangle (11,1.5);
% Domaine de Transition (Adamaoua/Centre) : 1.5 < y < 4.5
\fill[popgreen!30] (0,1.5) rectangle (11,4.5);
% Massif Adamaoua (hachures)
\draw[pattern=north east lines, pattern color=popgreen!60] (3,2.5) rectangle (8,4.5);
% Domaine Soudano-Sahélien (Nord/Bénoué) : 4.5 < y < 8.5
\fill[popgold!40] (0,4.5) rectangle (11,8.5);
% Domaine Sahélien (Extrême-Nord) : y > 8.5
\fill[poporange!30] (0,8.5) rectangle (11,12);
% Villes (points + labels)
\node[circle, fill=popdark, inner sep=1.5pt] (Ya) at (5,0.8) {};
\node[below right] at (Ya) {\small Yaoundé};
\node[circle, fill=popdark, inner sep=1.5pt] (Dou) at (7.5,0.5) {};
\node[below right] at (Dou) {\small Douala};
\node[circle, fill=popdark, inner sep=1.5pt] (Nga) at (5,3) {};
\node[above right] at (Nga) {\small Ngaoundéré};
\node[circle, fill=popdark, inner sep=1.5pt] (Ber) at (7,3.5) {};
\node[above left] at (Ber) {\small Bertoua};
\node[circle, fill=popdark, inner sep=1.5pt] (Gar) at (5,6) {};
\node[above right] at (Gar) {\small Garoua};
\node[circle, fill=popdark, inner sep=1.5pt] (Mar) at (4,9) {};
\node[above left] at (Mar) {\small Maroua};
\node[circle, fill=popdark, inner sep=1.5pt] (Kou) at (8.5,9.5) {};
\node[above right] at (Kou) {\small Kousséri};
% Axes routiers (N1 Nord-Sud, N10 Est-Ouest)
\draw[popblue, very thick, decoration={snake, amplitude=1pt, segment length=5pt}, decorate] (5,0.8) -- (5,9); % N1 approximatif
\draw[popblue, thick, dashed] (7,3.5) -- (5,3) -- (5,6); % Liaisons Est
% Fleche Nord
\node[draw, circle, fill=popblue!20, inner sep=3pt] at (11.5, 6) {\textbf{N}};
\draw[->, thick] (11.5, 5.5) -- (11.5, 6.5);
% Échelle indicative
\draw[|-|, thick] (0.5, -0.5) -- (3.5, -0.5) node[midway, below] {~300 km};
\end{tikzpicture}
\legende{Croquis simplifié des 4 ensembles biogéographiques du Cameroun. Vert foncé : Domaine équatorial (Forêt dense). Vert clair hachuré : Domaine de transition (Adamaoua, savanes boisées). Jaune : Domaine soudano-sahélien (Plaine de la Bénoué). Orange : Domaine sahélien (Extrême-Nord, steppe). Trait bleu continu : Route nationale N1. Trait bleu pointillé : Liaisons Est-Ouest.}
\end{popfigure}

\textbf{2. Légende structurée (à reporter sur la copie)}

\begin{center}
\begin{tabularx}{\linewidth}{|l|X|}\hline
\rowcolor{popblueL} \textbf{Signe} & \textbf{Signification} \\ \hline
\cellcolor{popgreen!60}\textcolor{white}{\rule{1cm}{4pt}} & Domaine équatorial (Forêt dense, >1500 mm) \\ \hline
\cellcolor{popgreen!30}\rule{1cm}{4pt} & Domaine tropical de transition (Adamaoua, savanes boisées, 1000-1500 m) \\ \hline
\cellcolor{popgold!40}\rule{1cm}{4pt} & Domaine soudano-sahélien (Plaine de la Bénoué, savane arborée, 1000-1500 mm) \\ \hline
\cellcolor{poporange!30}\rule{1cm}{4pt} & Domaine sahélien (Extrême-Nord, steppe, <1000 mm) \\ \hline
\textcolor{popblue}{\rule[2pt]{1cm}{3pt}} & Axe routier principal (N1 Yaoundé-Maroua) \\ \hline
\textcolor{popblue}{\rule[2pt]{1cm}{3pt}\rule[2pt]{1cm}{3pt}} & Axes secondaires (Liaisons Est-Ouest) \\ \hline
$\bullet$ & Villes principales (Chefs-lieux de région) \\ \hline
$\uparrow$ \textbf{N} & Orientation Nord \\ \hline
\end{tabularx}
\end{center}

\textbf{3. Commentaire organisé (Rédaction type Bac)}
\textbf{Titre :} \emph{Le Cameroun : une mosaïque biogéographique structurée en latitude et en altitude.}

\textbf{Introduction :} Le Cameroun, souvent qualifié d'« Afrique en miniature », présente une organisation biogéographique en lames latitudinales superposées, perturbée par le relief (massif de l'Adamaoua).

\textbf{Développement :}
\begin{itemize}
    \item \textbf{Au Sud (Domaine équatorial) :} Climat équatorial de transition (4 saisons) ou guinéen (2 saisons). Forêt dense sempervirente (sud-ouest) ou décidue (sud-est). Sols ferrallitiques lessivés, fragiles. Villes : Yaoundé, Douala, Ebolowa, Sangmélima. Activités : agriculture de rente (cacao, café, palmier à huile), exploitation forestière.
    \item \textbf{Au Centre (Domaine de transition / Adamaoua) :} Climat tropical d'altitude (tempéré). Végétation de savane arborée (herbes hautes, gallery forests). Sols ferrugineux tropicaux lessivés. Ville clé : Ngaoundéré (carrefour). Activités : élevage bovin (ranching), cultures vivrières (maïs, igname, pomme de terre).
    \item \textbf{Au Nord-Centre (Domaine soudano-sahélien) :} Climat tropical soudanien (saison sèche marquée 5-6 mois). Plaine de la Bénoué (colluvions fertiles). Savane arborée (karité, néré). Villes : Garoua, Ngaoundéré (sud). Activités : coton (rente), maïs, sorgho, élevage transhumant.
    \item \textbf{Au Nord (Domaine sahélien) :} Climat sahélien (saison sèche 7-9 mois, pluviométrie < 800 mm). Steppe à \emph{Acacia}, éboulis, bas-fonds (yaérés) du Logone et lac Tchad. Villes : Maroua, Kousséri. Activités : riziculture (irrigation SEMRY), pêche, élevage extensif, cultures de décrue.
\end{itemize}

\textbf{Conclusion :} Cette zonation en « étages » (latitude) et « marches » (altitude) conditionne les potentialités agricoles, la répartition de la population (forte au Sud et dans la Bénoué, faible au Nord extrême) et les défis d'aménagement (désertification nord, déforestation sud, enclavement est).

\textbf{Conclusion finale :} Le croquis et son commentaire démontrent que la diversité biogéographique camerounaise résulte du couplage \textbf{latitude/altitude/pluviométrie}, structurant 4 ensembles distincts aux potentialités et contraintes contrastées.
\end{exoresolu}

\begin{methode}[Analyser les relations Climat - Végétation - Sol (Système pédoclimatique)]
\textbf{Objectif :} Expliquer la nature des sols et de la végétation par les paramètres climatiques (pluviométrie, température, saisonnalité) et le relief.
\begin{enumerate}
    \item \textbf{Identifier le type climatique} à partir de données (omthermogramme, tableau pluviométrique) : Équatorial, Tropical, Sahélien.
    \item \textbf{Déduire la formation végétale climax} : Forêt dense $\leftrightarrow$ Pluviométrie $> 1500$ mm + saison sèche courte ($< 2$ mois) ; Savane $\leftrightarrow$ Saison sèche longue ($> 5$ mois) + incendies ; Steppe $\leftrightarrow$ Pluviométrie $< 800$ mm.
    \item \textbf{Relier au processus pédogénétique dominant} :
    \begin{itemize}
        \item \textbf{Ferrallitisation (Sud) :} Lessivage intense (pluies abondantes) $\rightarrow$ Appauvrissement en bases, accumulation Fe/Al $\rightarrow$ Sols rouges/jaunes, acides, peu fertiles (besoin d'apports).
        \item \textbf{Fersiallitisation (Centre/Nord) :} Lessivage modéré, saison sèche marquée $\rightarrow$ Migration argile/fer $\rightarrow$ Sols ferrugineux tropicaux (rouges/bruns), plus riches, sensibles à l'érosion.
        \item \textbf{Peu d'évolution / Salinisation (Extrême-Nord) :} Évaporation $>$ Précipitations $\rightarrow$ Remontée capillaire des sels $\rightarrow$ Sols bruns tropicaux, sols halomorphes (yaérés), sols peu épais (éboulis).
    \end{itemize}
    \item \textbf{Intégrer le facteur temps/relief :} Pentes raides $\rightarrow$ érosion, sols squelettiques ; Bas-fonds $\rightarrow$ hydromorphie, alluvions récentes fertiles.
\end{enumerate}
\cle{Formule clé} : $\text{Climat} + \text{Relief} + \text{Roche mère} + \text{Temps} + \text{Organismes} \rightarrow \text{Sol} + \text{Végétation}$.
\end{methode}

\begin{exoresolu}[Analyse pédoclimatique : Du Sud forestier à l'Extrême-Nord sahélien]
\textbf{Énoncé :} Le tableau ci-dessous donne des données climatiques et édaphiques pour trois stations camerounaises. Pour chaque station :
\begin{enumerate}
    \item Identifiez le type de climat et le domaine biogéographique.
    \item Nommez le type de sol dominant et expliquez sa genèse (processus principal).
    \item Citez la formation végétale naturelle correspondante.
\end{enumerate}
\begin{center}
\begin{tabular}{|c|c|c|c|c|c|}\hline
\textbf{Station} & \textbf{Lat.} & \textbf{Alt. (m)} & \textbf{P ann. (mm)} & \textbf{T moy. (°C)} & \textbf{Saison sèche (mois)} \\ \hline
Ebolowa & 2°54' N & 600 & 1850 & 24,5 & 2 \\ \hline
Ngaoundéré & 7°22' N & 1100 & 1550 & 22,0 & 5 \\ \hline
Maroua & 10°36' N & 400 & 750 & 28,5 & 7 \\ \hline
\end{tabular}
\end{center}
\tcblower
\textbf{Solution détaillée :}

\textbf{Station 1 : Ebolowa (Région du Sud)}
\begin{enumerate}
    \item \textbf{Climat :} Équatorial de transition (Guinéen) : Pluviométrie élevée ($1850 > 1500$ mm), saison sèche très courte (2 mois), températures stables. \textbf{Domaine :} Équatorial (Forêt dense).
    \item \textbf{Sol dominant :} \textbf{Sol ferrallitique} (ferralitique lessivé, rouge/jaune).
    \item \textbf{Genèse (Ferrallitisation) :} Lessivage intense et continu par les fortes pluies ($P > ETP$). Élimination quasi-totale des bases (Ca, Mg, K) et de la silice. Résidu insoluble : Oxydes de Fer (hématite $\rightarrow$ rouge) et d'Aluminium (gibbsite). Faible CEC (Capacité d'Échange Cationique), pH acide ($< 5,5$), réserve minérale faible.
    \item \textbf{Végétation :} \textbf{Forêt dense humide sempervirente} (essences : Ayous, Sapelli, Iroko, Moabi). Strates multiples, lianes, épiphytes.
\end{enumerate}

\textbf{Station 2 : Ngaoundéré (Région de l'Adamaoua)}
\begin{enumerate}
    \item \textbf{Climat :} Tropical d'altitude (ou soudanien d'altitude) : Pluviométrie forte ($1550$ mm) mais saison sèche nette (5 mois), températures adoucies par l'altitude ($1100$ m, $T_{moy}=22°C$). \textbf{Domaine :} Tropical de transition (Adamaoua).
    \item \textbf{Sol dominant :} \textbf{Sol ferrugineux tropical lessivé} (ou fersiallitique appauvri).
    \item \textbf{Genèse (Fersiallitisation / Ferruginisation) :} Alternance saison humide/sèche favorise la migration verticale de l'argile et du fer. Formation d'un horizon B textural (accumulation argile) et/ou nodulaire (concrétions ferrugineuses = « cuirasse » latéritique en surface par inversion de relief). Moins lessivé qu'au Sud, CEC moyenne, pH neutre à légèrement acide. Fertilité moyenne à bonne.
    \item \textbf{Végétation :} \textbf{Savane arborée à gallery forests} (Herbes hautes : \emph{Andropogon}, \emph{Hyparrhenia} ; Arbres : \emph{Daniellia}, \emph{Isoberlinia}, Karité). Feux de brousse annuels maintiennent la savane.
\end{enumerate}

\textbf{Station 3 : Maroua (Région de l'Extrême-Nord)}
\begin{enumerate}
    \item \textbf{Climat :} Sahélien (Tropical sec) : Pluviométrie faible ($750 < 800$ mm), saison sèche très longue (7 mois), températures élevées ($28,5°C$), forte évapotranspiration ($ETP \approx 2000-2500$ mm). \textbf{Domaine :} Sahélien.
    \item \textbf{Sol dominant :} \textbf{Sol brun tropical} (peu évolué) sur éboulis / \textbf{Sol halomorphe} (yaérés) / \textbf{Sol vertique} (argiles gonflantes) en bas-fond.
    \item \textbf{Genèse :} Pédogenèse limitée par l'aridité. Évaporation $>$ Précipitations $\rightarrow$ remontée capillaire $\rightarrow$ accumulation de sels (carbonates, gypse) en surface (croûtes salines) ou dans le profil (horizon calcique). Sur les pentes (Mandara) : érosion $\rightarrow$ sols squelettiques, lithosols. Dans les yaérés (bas-fonds Logone) : hydromorphie + évaporation $\rightarrow$ sols vertiques (gonflement/retrait), salins.
    \item \textbf{Végétation :} \textbf{Steppe sahélienne} (Herbes annuelles : \emph{Cenchrus biflorus}, \emph{Aristida} ; Arbres épineux : \emph{Acacia raddiana}, \emph{Balanites aegyptiaca}, \emph{Leptadenia}). Dégénérescence vers le nord (désertification).
\end{enumerate}

\textbf{Synthèse comparative (Tableau de révision) :}
\begin{center}
\begin{tabularx}{\linewidth}{|l|X|X|X|}\hline
\rowcolor{popblueL} \textbf{Critère} & \textbf{Ebolowa (Sud)} & \textbf{Ngaoundéré (Centre)} & \textbf{Maroua (Nord)} \\ \hline
\textbf{Climat} & Équatorial (Humide) & Tropical d'altitude (Sub-humide) & Sahélien (Aride) \\ \hline
\textbf{Processus sol} & Ferrallitisation (Lessivage total) & Fersiallitisation/Ferruginisation (Migration) & Pédogenèse faible / Salinisation (Remontée) \\ \hline
\textbf{Type de sol} & Ferrallitique (Rouge/Jaune, acide) & Ferrugineux tropical (Rouge/Brun, concrétions) & Brun tropical / Halomorphe / Vertique \\ \hline
\textbf{Végétation} & Forêt dense sempervirente & Savane arborée + Galeries & Steppe épineuse / Yaérés \\ \hline
\textbf{Fertilité} & Faible (chimique), forte (organique) & Moyenne à bonne & Faible (sauf bas-fonds alluvionnaires) \\ \hline
\end{tabularx}
\end{center}
\end{exoresolu}

\begin{methode}[Caractériser et comparer les grands groupes humains (Répartition, Densité, Activités)]
\textbf{Objectif :} Relier la répartition des populations et leurs activités aux contraintes/atouts des ensembles biogéographiques.
\begin{enumerate}
    \item \textbf{Calculer la densité} : $D = \frac{\text{Population}}{\text{Superficie}}$ (hab/km²). Comparer à la moyenne nationale ($\approx 60$ hab/km² en 2023).
    \item \textbf{Analyser la répartition} : Facteurs d'attraction (sols fertiles, villes, axes, ressources) vs Facteurs de répulsion (maladies, aridité, relief accidenté, enclavement).
    \item \textbf{Identifier les groupes ethno-linguistiques majeurs} par zone (Bantous, Semi-Bantous, Soudanais, Sahéliens) sans essentialiser (mobilité, mélange).
    \item \textbf{Lier activités dominantes / milieu} :
    \begin{itemize}
        \item Forêt : Agriculture itinérante sur brûlis (vivrier) + Perennes (rente), Chasse/Collecte, Exploitation forestière.
        \item Adamaoua : Élevage bovin (transhumance/ranching), Céréales (maïs), Pomme de terre.
        \item Bénoué : Coton (rente), Céréales (sorgho/maïs), Élevage, Pêche (Bénoué).
        \item Extrême-Nord : Riziculture irriguée (SEMRY), Mil/Sorgho (pluvial), Élevage extensif, Pêche (Lac Tchad/Logone), Commerce transfrontalier.
    \end{itemize}
    \item \textbf{Noter les dynamiques récentes :} Exode rural $\rightarrow$ villes (Douala, Yaoundé), Front pionnier (Est, Adamaoua), Pression foncière (Nord), Crises sécuritaires (Extrême-Nord, NOSO).
\end{enumerate}
\cle{Indicateur clé} : Taux d'urbanisation ($\approx 58\%$ en 2023), Accroissement naturel ($\approx 2,6\%$).
\end{methode}

\begin{exoresolu}[Dynamiques démographiques et activités par ensemble biogéographique]
\textbf{Énoncé :} À partir des données ci-après (estimations 2020), réalisez les tâches suivantes :
\begin{enumerate}
    \item Calculez la densité de population pour chaque région et classez-les.
    \item Expliquez les contrastes de densité observés entre le Sud (Région du Sud) et l'Extrême-Nord en les reliant aux milieux biogéographiques.
    \item Pour la région de l'Adamaoua, expliquez la faible densité malgré des potentialités agricoles réelles.
\end{enumerate}
\begin{center}
\begin{tabular}{|l|c|c|c|}\hline
\textbf{Région} & \textbf{Superficie (km²)} & \textbf{Population (hab)} & \textbf{Domaine biogéo principal} \\ \hline
Sud & 47 110 & 780 000 & Équatorial (Forêt) \\ \hline
Adamaoua & 63 701 & 1 200 000 & Transition (Savane altitude) \\ \hline
Nord & 66 090 & 2 500 000 & Soudano-sahélien (Bénoué) \\ \hline
Extrême-Nord & 34 263 & 4 300 000 & Sahélien (Plaine/Yaérés) \\ \hline
\end{tabular}
\end{center}
\tcblower
\textbf{Solution détaillée :}

\textbf{1. Calcul des densités ($D = Pop / Sup$) :}
\begin{itemize}
    \item \textbf{Sud :} $D = \frac{780\,000}{47\,110} \approx \mathbf{16,6\ hab/km^2}$.
    \item \textbf{Adamaoua :} $D = \frac{1\,200\,000}{63\,701} \approx \mathbf{18,8\ hab/km^2}$.
    \item \textbf{Nord :} $D = \frac{2\,500\,000}{66\,090} \approx \mathbf{37,8\ hab/km^2}$.
    \item \textbf{Extrême-Nord :} $D = \frac{4\,300\,000}{34\,263} \approx \mathbf{125,5\ hab/km^2}$.
\end{itemize}
\textbf{Classement croissant :} Sud ($16,6$) < Adamaoua ($18,8$) < Nord ($37,8$) < Extrême-Nord ($125,5$).
\emph{Moyenne nationale $\approx 60$ hab/km² : Seules l'Extrême-Nord et le Nord (partiellement) dépassent la moyenne.}

\textbf{2. Contraste Sud vs Extrême-Nord :}
\begin{itemize}
    \item \textbf{Extrême-Nord (Forte densité : 125 hab/km²) :} Milieu sahélien \textbf{a priori} contraignant (aridité, sols pauvres). \textbf{Explication :} Concentration historique dans les \textbf{yaérés} (bas-fonds alluvionnaires du Logone) et le long du fleuve (sols vertiques fertiles, eau disponible pour riziculture SEMRY, pêche). Agriculture intensive (riz, mil, sorgho) + Élevage dense + Commerce frontalier (Nigéria, Tchad). Forte croissance naturelle. \textbf{Paradoxe :} Forte densité sur zone à risque climatique (sécheresses, inondations).
    \item \textbf{Sud (Faible densité : 17 hab/km²) :} Milieu équatorial \textbf{a priori} favorable (eau, forêt). \textbf{Explication :} Contraintes fortes : \textbf{Sols ferrallitiques acides et lessivés} (faible fertilité chimique, toxicité Al), \textbf{Maladies endémiques} (paludisme, onchocercose, trypanosomose), \textbf{Enclavement} (réseau routier faible, forêt dense), \textbf{Activités extensives} (agriculture itinérante sur brûlis, chasse, cueillette, exploitation forestière peu consommatrice de main-d'œuvre). Exode rural massif vers Douala/Yaoundé.
\end{itemize}

\textbf{3. Cas de l'Adamaoua (Densité faible : 19 hab/km²) :}
\begin{itemize}
    \item \textbf{Atouts :} Climat tempéré d'altitude (santé), Sols ferrugineux moins lessivés (fertiles), Herbage abondant (élevage), Ressources minières (bauxite, or), Ville carrefour (Ngaoundéré, chemin de fer).
    \item \textbf{Contraintes/Freins :} \textbf{Enclavement historique} (loin des ports, route/rail unique), \textbf{Paludisme} (vallées), \textbf{Trypanosomose} (mouches tsé-tsé dans les gallery forests $\rightarrow$ frein à l'élevage bovin intensif hors zones traitées), \textbf{Activité extensive} (élevage nomade/transhumant = faible densité humaine/km²), \textbf{Front pionnier récent} (peuplement tardif par rapport au Nord et au Sud).
\end{itemize}

\textbf{Conclusion :} La densité de population ne dépend pas seulement du climat/pluviométrie, mais de la \textbf{valeur d'usage des sols} (fertilité effective, eau maîtrisable), de la \textbf{santé}, des \textbf{infrastructures} et de l'\textbf{histoire du peuplement}. L'Extrême-Nord sahélien est plus dense que le Sud forestier grâce à l'agriculture hydraulique (riz) et aux alluvions.
\end{exoresolu}

\begin{methode}[Rédiger une dissertation géographique structurée (Sujet type Bac)]
\textbf{Objectif :} Traiter un sujet de composition : « La diversité biogéographique du Cameroun : atouts et défis pour l'aménagement du territoire. »
\begin{enumerate}
    \item \textbf{Analyser le sujet :} Mots-clés : « Diversité biogéographique » (relief, climat, végé, sol, hommes), « Atouts » (potentialités), « Défis » (contraintes, risques), « Aménagement du territoire » (politiques publiques, infrastructures, équilibre).
    \item \textbf{Problématiser :} En quoi la mosaïque biogéographique camerounaise constitue-t-elle une ressource majeure mais source de déséquilibres spatiaux que l'aménagement doit corriger ?
    \item \textbf{Construire le plan détaillé (3 parties, 2 sous-parties) :}
    \begin{itemize}
        \item \textbf{I. Une diversité biogéographique source de potentialités complémentaires (Atouts).}
        \begin{itemize}
            \item A. Des ressources naturelles variées et sectorielles (Forêt/Bois, Mines/Adamaoua, Agro-pastoral/Nord, Pêche/Lac).
            \item B. Une diversité des milieux propice au tourisme et à l'énergie (Hydroélectricité/Sud, Solaire/Nord, Tourisme montagnard/ouest, Safari/Nord).
        \end{itemize}
        \item \textbf{II. Des contraintes physiques et humaines majeures freinant la mise en valeur (Défis).}
        \begin{itemize}
            \item A. Fragilité des écosystèmes et risques climatiques (Déforestation/Sud, Désertification/Nord, Érosion/Adamaoua, Inondations/Yaérés).
            \item B. Déséquilibres démographiques et enclavement (Vide forestier/Est-Sud, Surpeuplement/Extrême-Nord, Enclavement/Adamaoua-Est).
        \end{itemize}
        \item \textbf{III. Les réponses de l'aménagement du territoire : vers une valorisation intégrée et durable.}
        \begin{itemize}
            \item A. Infrastructures structurantes pour relier les ensembles (Route N1, Rail, Port Kribi, Barrages Lom Pangar/Nachtigal, Route Ring-Road).
            \item B. Politiques sectorielles adaptées par zone (PDR, PNDP, Plan d'urgence 3 ans, Grande muraille verte, Zonage forestier, SEMRY II).
        \end{itemize}
    \end{itemize}
    \item \textbf{Rédiger :} Introduction (Accroche, Définitions, Problématique, Annonce du plan). Développement (Arguments $\rightarrow$ Exemples Chiffrés $\rightarrow$ Cartes mentales). Conclusion (Bilan, Ouverture sur émergence 2035).
\end{enumerate}
\cle{Consigne Bac} : « Montrez que... », « Analysez... », « Discutez... » $\rightarrow$ Exige une démonstration argumentée, pas une simple énumération.
\end{methode}

\begin{exoresolu}[Dissertation type : Diversité biogéographique et aménagement]
\textbf{Sujet :} \emph{« Le Cameroun, Afrique en miniature : atouts et contraintes de sa diversité biogéographique pour le développement durable. »}
\tcblower
\textbf{Plan détaillé et éléments de rédaction (Introduction + Plan + Conclusion)}

\textbf{INTRODUCTION}
\begin{itemize}
    \item \textbf{Accroche :} L'expression « Afrique en miniature », popularisée par le géographe \textbf{Jean Gottmann}, résume l'extrême variété des milieux naturels du Cameroun concentrée sur 475 442 km².
    \item \textbf{Définitions :} Diversité biogéographique = combinaison unique de reliefs (plaines, plateaux, montagnes, volcans), climats (équatorial, tropical, sahélien), végétations (forêt, savane, steppe), sols et groupes humains.
    \item \textbf{Constat :} Cette diversité offre un panel de ressources exceptionnel (agricoles, minières, hydrauliques, touristiques) mais génère des disparités spatiales fortes.
    \item \textbf{Problématique :} \textbf{Comment l'aménagement du territoire peut-il transformer la diversité biogéographique camerounaise, source de potentialités contrastées, en levier d'un développement durable et équilibré à l'horizon 2035 ?}
    \item \textbf{Annonce du plan :} Nous analyserons d'abord les atouts d'une mosaïque de ressources complémentaires (I), puis les défis structurels qu'elle impose (II), pour enfin examiner les stratégies d'aménagement visant l'intégration nationale (III).
\end{itemize}

\textbf{DÉVELOPPEMENT}

\textbf{I. Une mosaïque de potentialités complémentaires : le capital naturel du Cameroun}
\textbf{A. Des ressources de rente et vivrières sectorielles, moteurs de l'économie.}
\begin{itemize}
    \item \textbf{Domaine équatorial (Sud/Est) :} 1\textsuperscript{er} producteur africain de bois tropical (Ayous, Sapelli), 5\textsuperscript{ème} mondial cacao, palmier à huile, hévéa. Bassin minier (Fer/Mbalam, Cobalt/Nkamouna, Or/Batouri). Hydroélectricité (Barrages : Lom Pangar, Memve'ele, Nachtigal $\rightarrow$ 70\% potentiel national).
    \item \textbf{Domaine de transition (Adamaoua/Quest/Nord-Ouest) :} « Grenier » céréalier (Maïs, Pomme de terre, Haricot), 1\textsuperscript{er} bassin d'élevage bovin (3-4 millions de têtes), Bauxite (Minim-Martap), Tourisme montagnard (Ring-Road, Lac Nyos/Oku).
    \item \textbf{Domaine soudano-sahélien (Nord/Bénoué) :} Coton (1\textsuperscript{ère} recette d'exportation agricole, SODECOTON), Sorgho/Maïs, Oignon (Galerie), Karité. Parc nationaux (Bénoué, Bouba Ndjida, Faro) $\rightarrow$ Tourisme cynégétique/photo.
    \item \textbf{Domaine sahélien (Extrême-Nord) :} Riziculture irriguée (SEMRY $\rightarrow$ 100 000+ tonnes/an), Pêche (Lac Tchad, Logone, Yaérés), Élevage (cheptel ovin/caprin/bovin important), Commerce transfrontalier (Nigéria, Tchad).
\end{itemize}
\textbf{B. Une diversité énergétique et touristique stratégique.}
Hydroélectricité (Sud) + Solaire/Éolien (Nord, fort ensoleillement) + Biomasse (partout). Tourisme : Balnéaire (Kribi, Limbé), Écologique (Parcs, Réserves Dja, Lobéké, Waza), Culturel (Royaumes Bamoun, Bamiléké, Lamidats Nord), Montagnard (Mont Cameroun, Manengouba).

\textbf{II. Des contraintes structurelles : fragilités écologiques et déséquilibres spatiaux}
\textbf{A. Fragilité des écosystèmes et risques climatiques exacerbés.}
\begin{itemize}
    \item \textbf{Sud :} Déforestation (agriculture itinérante, exploitation forestière illégale, agro-industrie palmier/hévéa) $\rightarrow$ Perte biodiversité, érosion sols ferrallitiques nus, perturbation cycle de l'eau.
    \item \textbf{Nord/Extrême-Nord :} Désertification (avancée du désert, dégradation sols, perte couverture végétale, surexploitation bois-énergie). Sécheresses récurrentes (1973, 1984, 2010+). Inondations paradoxales (Yaérés, Logone) dues à pluies intenses soudaines + sols crustés.
    \item \textbf{Adamaoua/Reliefs :} Érosion hydrique (pentes, sols nus post-feux), Glissements de terrain (Mont Cameroun, Manengouba), Risque volcanique/limnique (Lac Nyos, Monoun).
\end{itemize}
\textbf{B. Déséquilibres démographiques et enclavement : le « Cameroun utile » vs « Cameroun profond ».}
\begin{itemize}
    \item \textbf{Contraste Nord/Sud inversé :} Extrême-Nord (125 hab/km²) surpeuplé, pression foncière, conflits agro-pasteurs, insécurité (Boko Haram) $\rightarrow$ Déplacés internes. Sud/Est (15-20 hab/km²) « vides », enclavés, ressources non valorisées localement (export brut).
    \item \textbf{Urbanisation macrocéphalique :} Douala (économique, port) + Yaoundé (politique) = 40\% population urbaine. Pompage des ressources/vitalité des régions.
    \item \textbf{Enclavement Est/Adamaoua :} Routes dégradées (N10, N15), chemin de fer unique (Transcamerounais) saturé, coût transport élevé $\rightarrow$ compétitivité faible.
\end{itemize}

\textbf{III. L'aménagement du territoire : politiques d'intégration et de valorisation durable}
\textbf{A. Infrastructures structurantes : relier les ensembles pour un marché national unifié.}
\begin{itemize}
    \item \textbf{Transport :} Achèvement autoroute Yaoundé-Douala, Route N1 (Yaoundé-Maroua) bitumée, Programme routes régionales (Ring-Road NW, N10 Est), Port en eau profonde de Kribi (désenclavement Sud/Est/CAE), Chemin de fer (réhabilitation + extension Ngaoundéré-Ngaoundal).
    \item \textbf{Énergie :} Barrages (Nachtigal 420 MW, Lom Pangar régulation, Memve'ele, Song Dong) $\rightarrow$ Objectif 3000 MW (2030). Interconnexion réseau Nord-Sud (ligne 225/400 kV) pour valoriser solaire Nord.
\end{itemize}
\textbf{B. Politiques sectorielles différenciées et développement local.}
\begin{itemize}
    \item \textbf{Forêt :} Zonage (UFA, Parcs, Communautaires), Certification FSC/PEFC, Transformation locale 2\textsuperscript{ème}/3\textsuperscript{ème} transformation (interdiction grumes brutes).
    \item \textbf{Agriculture :} PDC (Plantes de rente), PDA (Alimentaires), Mécanisation (Bénoué), Irrigation (SEMRY, Maga, Logone), Intrants subventionnés.
    \item \textbf{Environnement :} « Grande Muraille Verte » (reboisement Nord), Restauration paysages forestiers (AFR100), Gestion intégrée bassins versants (Sanaga, Logone, Bénoué).
    \item \textbf{Décentralisation :} Transferts compétences/ressources aux Communes/Régions (PNDP, FEICOM) pour développement endogène.
\end{itemize}

\textbf{CONCLUSION}
\textbf{Bilan :} La diversité biogéographique camerounaise est un \textbf{capital naturel majeur} (ressources complémentaires Nord/Sud/Est/Ouest) mais sa valorisation est entravée par la \textbf{fragilité écologique} (déforestation, désertification) et des \textbf{déséquilibres spatiaux hérités} (enclavement, concentration côtière).
\textbf{Ouverture :} L'émergence à l'horizon 2035 (SND30) suppose une \textbf{transition écologique} (économie verte, adaptation climat) et une \textbf{justice territoriale} (aménagement polycentrique, corridors de développement Nord-Sud, Est-Ouest) transformant la diversité en résilience nationale.
\end{exoresolu}

\begin{methode}[Analyser un document statistique / graphique (Évolution, Parts, Taux)]
\textbf{Objectif :} Exploiter un tableau de données ou un graphique (évolution production, pluviométrie, population) pour en tirer une analyse géographique.
\begin{enumerate}
    \item \textbf{Lire les métadonnées :} Titre, Source, Année, Unités, Périmètre géographique.
    \item \textbf{Calculer les indicateurs clés :}
    \begin{itemize}
        \item \textbf{Taux de variation (TV) :} $TV = \frac{V_{final} - V_{initial}}{V_{initial}} \times 100$ (en \%).
        \item \textbf{Taux moyen annuel de variation (TMAV) :} $TMAV = \left( \frac{V_{final}}{V_{initial}} \right)^{\frac{1}{n}} - 1$ (si $n$ années).
        \item \textbf{Part en \% :} $\% = \frac{\text{Partie}}{\text{Total}} \times 100$.
        \item \textbf{Contribution à la croissance :} $\frac{\Delta \text{Partie}}{\Delta \text{Total}} \times 100$.
    \end{itemize}
    \item \textbf{Construire le graphique adapté :} Histogramme (comparaison), Courbe (évolution temporelle), Camembert (structure en \%), Diagramme en barres empilées (évolution structure).
    \item \textbf{Commenter (DESC - EXPLIC - INTERPRÉT) :}
    \begin{itemize}
        \item \textbf{Description :} Tendances (hausse/baisse/stable), seuils, ruptures, valeurs extrêmes, écarts.
        \item \textbf{Explication :} Facteurs climatiques (sécheresse), politiques (subventions, SODECOTON/SEMRY), économiques (cours mondiaux), sécuritaires.
        \item \textbf{Interprétation :} Conséquences sur l'aménagement, la sécurité alimentaire, l'environnement.
    \end{itemize}
\end{enumerate}
\cle{Piège} : Ne pas confondre « variation en valeur absolue » et « variation en \% ». Une faible variation absolue sur une grande base = faible \%. Une forte variation sur une petite base = fort \% mais faible impact total.
\end{methode}

\begin{exoresolu}[Analyse statistique : Évolution de la production cotonnière et rizicole (2010-2023)]
\textbf{Énoncé :} Le tableau ci-dessous présente l'évolution de la production de coton-graine et de riz paddy au Cameroun (en tonnes).
\begin{center}
\begin{tabular}{|c|c|c|}\hline
\textbf{Année} & \textbf{Coton-graine (t)} & \textbf{Riz paddy (t)} \\ \hline
2010 & 120 000 & 130 000 \\ \hline
2015 & 250 000 & 180 000 \\ \hline
2020 & 300 000 & 300 000 \\ \hline
2023 & 280 000 & 350 000 \\ \hline
\end{tabular}
\end{center}
\textbf{Travail demandé :}
\begin{enumerate}
    \item Calculez le taux de variation de la production cotonnière entre 2010 et 2023. Interprétez.
    \item Calculez le Taux Moyen Annuel de Variation (TMAV) de la production de riz entre 2010 et 2023.
    \item En 2023, quelle part (en \%) représente le riz dans la somme des deux productions ? Que signifie cette évolution structurelle ?
    \item Proposez un graphique adapté pour représenter l'évolution de ces deux productions sur la période. Justifiez votre choix.
\end{enumerate}
\tcblower
\textbf{Solution détaillée :}

\textbf{1. Taux de variation Coton (2010-2023)}
$V_{2010} = 120\,000$ ; $V_{2023} = 280\,000$.
$TV = \frac{280\,000 - 120\,000}{120\,000} \times 100 = \frac{160\,000}{120\,000} \times 100 = \mathbf{+133,3\%}$.
\textbf{Interprétation :} La production a plus que doublé (+133\%) en 13 ans. Hausse forte 2010-2015 (réforme SODECOTON, meilleurs prix, intrants), stagnation/baisse récente 2020-2023 (insécurité Extrême-Nord/Bénoué, coûts intrants, cours mondial bas, aléas climatiques).

\textbf{2. TMAV Riz (2010-2023)}
$n = 2023 - 2010 = 13$ ans.
$V_{2010} = 130\,000$ ; $V_{2023} = 350\,000$.
$TMAV = \left( \frac{350\,000}{130\,000} \right)^{\frac{1}{13}} - 1 = (2,6923)^{0,0769} - 1$.
$\ln(2,6923) \approx 0,9906$.
$0,0769 \times 0,9906 \approx 0,0762$.
$e^{0,0762} \approx 1,0792$.
$TMAV \approx 0,0792 = \mathbf{7,9\%\ par\ an}$.
\textbf{Interprétation :} Croissance soutenue et régulière (~8\%/an), plus rapide que le coton. Liée à la politique de substitution aux importations (Stratégie Nationale de Développement de la Riziculture - SNDR), extension périmètres irrigués (SEMRY, Maga, Ndop, Santchou), variétés améliorées (NERICA).

\textbf{3. Part du Riz en 2023}
Total 2023 = $280\,000 + 350\,000 = 630\,000$ t.
$\% \text{Riz} = \frac{350\,000}{630\,000} \times 100 \approx \mathbf{55,6\%}$.
$\% \text{Coton} = 44,4\%$.
\textbf{Sens :} Basculement structurel. En 2010, Coton (48\%) > Riz (52\% $\approx$). En 2023, Riz devient la première production en volume parmi les deux. Témoigne de la priorité donnée à la \textbf{sécurité alimentaire} (riz = aliment de base urbain) vs culture de rente d'exportation (coton) soumise aux cours mondiaux et à l'insécurité.

\textbf{4. Choix graphique}
\textbf{Graphique :} \textbf{Courbes d'évolution temporelle (double courbe) sur un même repère orthonormé.}
\textbf{Justification :}
\begin{itemize}
    \item Permet de \textbf{comparer les dynamiques} (pentes, ruptures) des deux filières sur la même période.
    \item Met en évidence le \textbf{croisement} (vers 2018-2019) où le riz dépasse le coton.
    \item Lisible pour 4 dates seulement (histogramme groupé possible mais moins lisible pour voir la tendance continue).
    \item Axes : Abscisses = Années (2010, 2015, 2020, 2023). Ordonnées = Production (tonnes), échelle commune (0 à 400 000 par pas de 50 000). Légende : Courbe continue Coton, Courbe pointillée Riz. Titre : « Évolution contrastée des productions cotonnière et rizicole au Cameroun (2010-2023) ». Source : MINADER/INS.
\end{itemize}

\textbf{Schémas des courbes (représentation mentale) :}
\begin{popfigure}
\begin{tikzpicture}
\begin{axis}[
    width=\linewidth, height=8cm,
    xlabel={Année}, ylabel={Production (milliers de tonnes)},
    xtick={2010,2015,2020,2023}, xticklabels={2010,2015,2020,2023},
    ymin=0, ymax=400, ytick={0,100,200,300,400},
    legend pos=north west, grid=major, grid style={popline},
    title={Évolution Coton vs Riz (2010-2023)}
]
\addplot[popblue, thick, mark=*] coordinates {(2010,120) (2015,250) (2020,300) (2023,280)};
\addlegendentry{Coton-graine}
\addplot[poporange, thick, mark=square*, dashed] coordinates {(2010,130) (2015,180) (2020,300) (2023,350)};
\addlegendentry{Riz paddy}
\end{axis}
\end{tikzpicture}
\legende{Double courbe d'évolution. Croisement vers 2019. Hausse régulière du riz, pic coton en 2020 puis baisse.}
\end{popfigure}
\end{exoresolu}

\begin{methode}[Réaliser et commenter un croquis de coupe transversale (Profil topographique)]
\textbf{Objectif :} Visualiser la structure en « marches d'escalier » du relief camerounais et ses conséquences climatiques/végétales.
\begin{enumerate}
    \item \textbf{Tracer l'axe de coupe :} Généralement Sud-Nord (ex: Kribi $\rightarrow$ Yaoundé $\rightarrow$ Ngaoundéré $\rightarrow$ Garoua $\rightarrow$ Maroua $\rightarrow$ Lac Tchad) ou Ouest-Est (Mont Cameroun $\rightarrow$ Bénoué).
    \item \textbf{Repérer les points caractéristiques (Abscisse / Altitude) :}
    \begin{itemize}
        \item Côte (0 m), Plaine côtière (< 200 m).
        \item Plateau Sud-Camerounais (600-800 m) : Yaoundé ~750 m.
        \item Massif de l'Adamaoua (1000-1500 m, pics > 2000 m) : Ngaoundéré ~1100 m, Ligne de partage des eaux.
        \item Plaine de la Bénoué (200-300 m) : Garoua ~230 m. Faille tectonique.
        \item Massifs volcaniques Nord (Mandara, 1000-1400 m) : Mokolo ~800 m.
        \item Plaine du Lac Tchad (200-300 m) : Maroua ~400 m, Kousséri ~280 m.
    \end{itemize}
    \item \textbf{Dessiner le profil :} Échelle verticale exagérée (ex: 1 cm pour 200 m) vs horizontale (1 cm pour 100 km) pour voir le relief (EV $\approx$ 500).
    \item \textbf{Annoter :} Domaines biogéographiques, vents dominants (Mousson SW, Harmattan NE), cours d'eau, villes, activités.
\end{enumerate}
\cle{Exagération verticale} : Indispensable (EV $\approx$ 50 à 500) pour visualiser le relief « plat » à l'échelle horizontale.
\end{methode}

\begin{exoresolu}[Croquis de coupe Sud-Nord : Structure en marches du relief camerounais]
\textbf{Énoncé :} Réalisez un croquis de coupe topographique suivant l'axe Yaoundé - Ngaoundéré - Garoua - Maroua. Montrez l'organisation en marches du relief, les domaines climatiques/végétaux associés et les grands axes de communication.
\tcblower
\textbf{Solution détaillée :}

\textbf{1. Données altimétriques clés (Profil Sud-Nord approximatif)}
\begin{center}
\begin{tabular}{|l|c|c|}\hline
\textbf{Repère} & \textbf{Distance cumulative (km)} & \textbf{Altitude (m)} \\ \hline
Yaoundé (Pl. Sud) & 0 & 750 \\ \hline
Pied Adamaoua (Nditam) & $\sim$ 300 & 600 \\ \hline
Ngaoundéré (Haut Adamaoua) & $\sim$ 450 & 1100 \\ \hline
Col de la Ngaoundéré (Partage eaux) & $\sim$ 500 & 1300 \\ \hline
Pied Nord Adamaoua (Tignère) & $\sim$ 600 & 800 \\ \hline
Garoua (Plaine Bénoué) & $\sim$ 750 & 230 \\ \hline
Maroua (Piémont Mandara) & $\sim$ 950 & 400 \\ \hline
Kousséri (Lac Tchad) & $\sim$ 1100 & 280 \\ \hline
\end{tabular}
\end{center}

\textbf{2. Réalisation du croquis (TikZ - Profil vertical exagéré)}

\begin{popfigure}
\begin{tikzpicture}[xscale=0.01, yscale=0.005] % 1 unit x = 10km (0.1cm), 1 unit y = 1m (0.005cm) -> EV = 10000/0.5 = 200? Wait.
% Let's use explicit coordinates in cm for clarity.
% Horiz: 1100 km -> 11 cm (1cm = 100km). Vert: 1300 m -> 6.5 cm (1cm = 200m). EV = 500.
% Coordinates (x_cm, y_cm = alt/200)
\end{tikzpicture}
\end{popfigure}

% Re-drawing with explicit cm coordinates for reliability
\begin{popfigure}
\begin{tikzpicture}[scale=1.0]
% Horizontal: 1cm = 100km. Vertical: 1cm = 200m. EV = 500.
% Points: (dist_km/100, alt_m/200)
\draw[popdark, very thick] 
(0, 3.75) -- 
(3, 3.0) -- 
(4.5, 5.5) -- 
(5, 6.5) -- 
(6, 4.0) -- 
(7.5, 1.15) -- 
(9.5, 2.0) -- 
(11, 1.4);
% Remplissage domaines (couleurs transparentes sous la courbe)
% Sud (Forêt) : 0 à 3.5cm (environ)
\fill[popgreen!20] (0,0) -- (0,3.75) -- (3,3.0) -- (3.5,3.2) -- (3.5,0) -- cycle;
% Transition (Adamaoua) : 3.5 à 6.5cm
\fill[popgreen!10!popgold!20] (3.5,0) -- (3.5,3.2) -- (6,4.0) -- (6.5,2.5) -- (6.5,0) -- cycle;
% Hachures Adamaoua haut
\draw[pattern=north east lines, pattern color=popgreen!50] (4, 3.5) rectangle (5.5, 6.5);
% Nord (Bénoué) : 6.5 à 9.5cm
\fill[popgold!30] (6.5,0) -- (6.5,2.5) -- (9.5, 2.0) -- (9.5,0) -- cycle;
% Extrême-Nord (Sahélien) : 9.5 à 11cm
\fill[poporange!30] (9.5,0) -- (9.5,2.0) -- (11, 1.4) -- (11,0) -- cycle;
% Villes (points sur le profil)
\node[circle, fill=popdark, inner sep=1.5pt, label=above:\tiny Yaoundé (750m)] at (0, 3.75) {};
\node[circle, fill=popdark, inner sep=1.5pt, label=above:\tiny Ngaoundéré (1100m)] at (4.5, 5.5) {};
\node[circle, fill=popdark, inner sep=1.5pt, label=above:\tiny Garoua (230m)] at (7.5, 1.15) {};
\node[circle, fill=popdark, inner sep=1.5pt, label=above:\tiny Maroua (400m)] at (9.5, 2.0) {};
% Fleches vents
\draw[->, thick, popblue] (1, 7) -- (1, 5.5) node[left, pos=0.5] {\small Mousson SW};
\draw[->, thick, poporange] (10, 7) -- (10, 5) node[right, pos=0.5] {\small Harmattan NE};
% Echelle horizontale
\draw[|-|, thick] (0, -0.8) -- (11, -0.8) node[midway, below] {~1100 km (Yaoundé - Lac Tchad)};
% Echelle verticale
\draw[|-|, thick] (11.5, 0) -- (11.5, 6.5) node[midway, right] {1300 m (EV $\approx$ 500)};
% Legende domaines sur le graphique
\node[align=center, popgreen!80!popdark] at (1.5, 4.5) {Domaine\\Équatorial\\Forêt dense};
\node[align=center, popgreen!60!popdark] at (5, 6.8) {Adamaoua\\Transition\\Savane altitude};
\node[align=center, popgold!80!popdark] at (8, 1.5) {Bénoué\\Soudano-\\sahélien};
\node[align=center, poporange!80!popdark] at (10.2, 1.0) {Extrême-Nord\\Sahélien};
\end{tikzpicture}
\legende{Profil topographique Sud-Nord (Yaoundé - Lac Tchad). Exagération verticale ~500x. Vert : Forêt équatoriale. Vert/Jaune hachuré : Adamaoua (partage des eaux Sanaga/Bénoué/Logone). Jaune : Plaine de la Bénoué (faille). Orange : Plaine sahélienne (Lac Tchad). Flèches bleues = Mousson (pluies). Flèches oranges = Harmattan (sec).}
\end{popfigure}

\textbf{3. Commentaire structuré (Type Bac - 15-20 lignes)}
\textbf{Titre :} \emph{Profil transversal Sud-Nord : l'architecture en marches du relief camerounais, moteur de la diversité biogéographique.}

Le profil Sud-Nord révèle une structure en \textbf{quatre marches successives} qui commande l'organisation climatique et végétale du pays.
\begin{enumerate}
    \item \textbf{La marche côtière et le plateau Sud (0-300 km, 600-800 m) :} Large plaine côtière (bas-fonds, mangroves) puis plateau du Sud-Cameroun (Yaoundé 750 m). Climat équatorial (pluies abondantes, 2 saisons sèches courtes). Végétation : Forêt dense humide. Obstacle modéré aux flux.
    \item \textbf{La marche majeure : Le Massif de l'Adamaoua (300-600 km, 1000-1500 m) :} Véritable \textbf{charnière climatique et hydrographique} (ligne de partage des eaux : Sanaga/Atlantique au Sud ; Bénoué/Niger et Logone/Tchad au Nord). Altitude $\rightarrow$ températures plus fraîches, pluviométrie orographique. Végétation : Savanes boisées (gallery forests). Barrière historique aux échanges Nord-Sud (franchi par rail/route N1).
    \item \textbf{La marche de la Bénoué (600-800 km, 200-300 m) :} Large plaine d'effondrement (faille tectonique), traversée par la Bénoué navigable (saisonnière). Climat tropical soudanien (saison sèche 5-6 mois). Sols alluvionnaires/colluvionnaires fertiles. Végétation : Savane arborée (karité). Grenier coton/céréales/élevage.
    \item \textbf{La marche sahélienne (800-1100 km, 200-400 m) :} Piémont des Mandara (volcanique) puis plaine du Lac Tchad (ancien lac quaternaire). Climat sahélien (aride, < 800 mm). Végétation : Steppe, yaérés (bas-fonds inondables fertiles). Activités : Riziculture irriguée, pêche, élevage extensif.
\end{enumerate}
\textbf{Synthèse :} Le relief agit comme un \textbf{filtre} : il force la mousson à déverser ses pluies au Sud et sur l'Adamaoua (effet Foehn asséchant au Nord), créant le gradient pluviométrique (2000 mm $\rightarrow$ 700 mm) qui structure les 4 ensembles biogéographiques. L'aménagement (Route N1, Chemin de fer, Barrage Lom Pangar) doit franchir ces marches pour intégrer l'espace national.
\end{exoresolu}

\begin{methode}[Évaluer les potentialités et contraintes d'un ensemble pour un aménagement donné (Analyse SWOT simplifiée)]
\textbf{Objectif :} Pour un ensemble biogéographique donné (ex: Bassin de la Bénoué, Massif forestier Sud, Yaérés Extrême-Nord), lister Forces, Faiblesses, Opportunités, Menaces pour un projet d'aménagement (ex: Pôle agro-industriel, Corridor routier, Parc national).
\begin{enumerate}
    \item \textbf{Délimiter l'espace d'étude} (frontières physiques/administratives).
    \item \textbf{Remplir la matrice SWOT (FFOM) en géographie :}
    \begin{center}
    \begin{tabular}{|c|c|}\hline
    \textbf{Forces (Internes +)} & \textbf{Faiblesses (Internes -)} \\ \hline
    Ressources (sol, eau, minerai, main d'œuvre), & Enclavement, pauvreté sols, maladies, \\
    Infrastructures existantes, Compétences locales. & faible densité, insécurité, dégradation. \\ \hline
    \textbf{Opportunités (Externes +)} & \textbf{Menaces (Externes -)} \\ \hline
    Marchés (sous-région, monde), Politiques publiques, & Changement climatique, Cours matières premières, \\
    Investisseurs, Coopération, Nouvelles tech. & Conflits, Exode rural, Pollution. \\ \hline
    \end{tabular}
    \end{center}
    \item \textbf{Croiser les cases (Stratégies) :}
    \begin{itemize}
        \item \textbf{FO (Force-Opportunité) : }Offensive** (Ex: Sols fertiles Bénoué + Demande coton/monde $\rightarrow$ Extension périmètres SODECOTON).
        \item \textbf{FA (Force-Menace) : }Protection** (Ex: Forêt dense + Marché carbone/REDD+ $\rightarrow$ Conservation rémunérée).
        \item \textbf{FaO (Faiblesse-Opportunité) : }Renforcement** (Ex: Enclavement Est + Route N10/Kribi $\rightarrow$ Valorisation mines/bois).
        \item \textbf{FaM (Faiblesse-Menace) : }Défense/Urgence** (Ex: Yaérés + Sécheresses/Inondations + Boko Haram $\rightarrow$ Résilience hydraulique, Sécurité, Déplacés).
    \end{itemize}
    \item \textbf{Conclure par des recommandations d'aménagement priorisées.}
\end{enumerate}
\cle{Application Bac} : Souvent demandé dans le sujet de composition (Partie III) ou étude de document « Proposez des mesures pour... ».
\end{methode}

\begin{exoresolu}[Analyse SWOT : Le bassin de la Bénoué (Région du Nord) face au développement agro-industriel]
\textbf{Énoncé :} Le bassin de la Bénoué (Région du Nord) est identifié comme un pôle de croissance agricole majeur (PDR, SND30). Réalisez une analyse SWOT (FFOM) de cet espace pour le développement d'un pôle agro-industriel intégré (Coton, Maïs, Élevage, Transformation locale).
\tcblower
\textbf{Solution détaillée :}

\textbf{MATRICE SWOT DU BASSIN DE LA BÉNOUÉ}

\begin{center}
\begin{tabularx}{\linewidth}{|X|X|}\hline
\rowcolor{popblueL} \textbf{FORCES (Atouts internes)} & \textbf{FAIBLESSES (Contraintes internes)} \\ \hline
\begin{itemize}
    \item \textbf{Sols :} Alluvions/colluvions riches (vertisols, sols ferrugineux peu lessivés), plats $\rightarrow$ mécanisable.
    \item \textbf{Eau :} Fleuve Bénoué (navigable 6 mois), nappes phréatiques accessibles, potentialité irrigation (barrages Léré, Lagdo amont).
    \item \textbf{Climat :} Soudanien (1000-1200 mm), saison sèche marquée $\rightarrow$ séchage récoltes, 2 spéculations/an possible (irrigué).
    \item \textbf{Expertise :} 50 ans SODECOTON (encadrement, intrants, pistes, usines égrenage), filière coton structurée.
    \item \textbf{Foncier :} Disponibilité relative (densité 38 hab/km²), chefferies foncières organisées (Lamidats).
    \item \textbf{Main d'œuvre :} Jeune, rurale, expérimentée cultures annuelles.
\end{itemize} & \begin{itemize}
    \item \textbf{Enclavement :} Distance ports (Douala 1000 km, Kribi 1100 km). Route N1 unique, dégradée par poids lourds. Rail vétuste (Transcamerounais).
    \item \textbf{Énergie :} Déficit électrique (usines égrenage/gins tournent groupes électrogènes). Coût énergie élevé.
    \item \textbf{Dégradation sols :} Tassement (passage machines), perte Matière Organique, érosion hydrique (cultures sur pente piémonts), salinisation débuts (irrigation mal drainée).
    \item \textbf{Sanitaire :} Paludisme endémique, Trypanosomose (mouche tsé-tsé résiduelle bas-fonds), Bilharziose (plans d'eau).
    \item \textbf{Transformation locale faible :} Export brut (fibre coton, graines) $\rightarrow$ faible valeur ajoutée, chômage jeunes.
    \item \textbf{Gouvernance foncière :} Conflits agriculteurs/éleveurs (couloirs transhumance bloqués), pression foncière peri-urbaine (Garoua).
\end{itemize} \\ \hline
\rowcolor{poporangeL} \textbf{OPPORTUNITÉS (Facteurs externes favorables)} & \textbf{MENACES (Risques externes)} \\ \hline
\begin{itemize}
    \item \textbf{Marchés :} Demande croissante céréales (Nigéria, Tchad, CEMAC), Coton bio/équitable, Huile de coton, Tourteaux (élevage).
    \item \textbf{Politiques :} SND30 (Pôles de croissance), PDR (Programme Développement Rural), ZLECAf (marché unique africain), Port Kribi (nouvelle sortie).
    \item \textbf{Infrastructures :} Projet route Garoua-Ngaoundéré (bitume), Réhabilitation rail, Barrage de Lom Pangar (régulation Bénoué, électricité), Solaire (irradiation forte).
    \item \textbf{Technologies :} Semences améliorées (IITA, IRAD), Agriculture de précision, Solaire photovoltaïque décentralisé, Biotechnologies (variétés tolérantes sécheresse).
    \item \textbf{Financement :} Banques développement (BAD, BDEAC), Fonds climat (Adaptation), Partenariats Public-Privé (PPP).
\end{itemize} & \begin{itemize}
    \item \textbf{Climatiques :} Variabilité accrue (sécheresses précoces, pluies violentes $\rightarrow$ inondations Bénoué), Réchauffement (+2°C horizon 2050 $\rightarrow$ raccourcissement cycle cultural).
    \item \textbf{Sécuritaires :} Insécurité Boko Haram (Extrême-Nord), Coupeurs de route (Adamaoua/Nord), Vol bétail $\rightarrow$ découragement investissement, déplacés.
    \item \textbf{Économiques :} Cours mondiaux coton volatils (subventions US/UE), Concurrence riz/coton asiatique bon marché, Inflation intrants (engrais, fuel).
    \item \textbf{Environnementales :} Envahissement \emph{Typha} (bénoué), Perte biodiversité (parcs braconnage), Pollution pesticides (cotton).
    \item \textbf{Démographiques :} Pression foncière croissante (densité x2 en 20 ans), Exode jeunes vers villes/Orpaillage/Étranger.
\end{itemize} \\ \hline
\end{tabularx}
\end{center}

\textbf{STRATÉGIES D'AMÉNAGEMENT PRIORITAIRES (Croisement SWOT)}

\begin{enumerate}
    \item \textbf{Stratégie OFFENSIVE (Forces + Opportunités) : \emph{« Pôle Agro-Industriel Vert »}.}
    \begin{itemize}
        \item Valoriser sols/eau/expertise coton + Marché régional/ZLECAf + Énergie solaire/Lom Pangar.
        \item \textbf{Actions :} Usines transformation locale (Huilerie, Filature, Aliments bétail), Périmètres irrigués solaires (goutte-à-goutte maïs/légumes), Certification coton durable (Better Cotton), Plateforme logistique multimodale (Route/Rail/Fleuve) à Garoua.
    \end{itemize}
    \item \textbf{Stratégie de RENFORCEMENT (Faiblesses + Opportunités) : \emph{« Désenclavement \& Énergie »}.}
    \begin{itemize}
        \item Corriger enclavement/énergie + Programmes infrastructures/Solaire.
        \item \textbf{Actions :} Réhabilitation prioritaire Rail Ngaoundéré-Garoua (fret coton/engrais), Centrale solaire 50 MW Garoua (autoconsommation usines), Pistes rurales structurantes (couloirs transhumance), Formation professionnelle (maintenance, transformation).
    \end{itemize}
    \item \textbf{Stratégie de PROTECTION (Forces + Menaces) : \emph{« Résilience Climatique \& Sécuritaire »}.}
    \begin{itemize}
        \item Protéger capital sol/eau/biodiversité + Risques Climat/Insécurité.
        \item \textbf{Actions :} Agriculture de conservation (SCV - Semis Direct sous Couvert Végétal), Agroforesterie (Faïdherbia, Karité), Gestion intégrée bassin Bénoué (CBLB), Sécurisation couloirs transhumance (commissions mixtes), Assurance indicielle climatique (index weather).
    \end{itemize}
    \item \textbf{Stratégie de DÉFENSE (Faiblesses + Menaces) : \emph{« Gestion Conflits \& Adaptation »}.}
    \begin{itemize}
        \item Réduire vulnérabilité foncière/sanitaire + Pression démographique/Insécurité.
        \textbf{Actions :} Plans d'occupation sols (POS) communaux, Délimitation/balisage couloirs/pâturages, Lutte anti-vectorielle (tsé-tsé, moustiques), Intégration déplacés (formations métiers agro), Filets sociaux (transferts monétaires).
    \end{itemize}
\end{enumerate}

\textbf{Conclusion :} Le bassin de la Bénoué possède les \textbf{fondamentaux physiques et humains} (Forces) pour devenir le « Grenier de la CEMAC », mais son avenir dépend de la \textbf{levée des verrous infrastructuraux (énergie, transport)} et de l'\textbf{adaptation au changement climatique/sécuritaire} via une gouvernance foncière apaisée et une industrialisation locale (Transformation).
\end{exoresolu}

\begin{attention}[Les 5 erreurs qui coûtent des points au Bac]
\begin{enumerate}
    \item \textbf{Confondre « Domaine biogéographique » et « Région administrative ».} \textit{Exemple :} Dire « La région de l'Adamaoua est un domaine sahélien » (Faux : c'est la transition). Le domaine sahélien correspond à l'Extrême-Nord. \textbf{Reflexe :} Utiliser les critères physiques (climat, végé, sol, altitude) pour nommer l'ensemble, pas le nom de la région.
    \item \textbf{Oublier l'altitude comme facteur explicatif majeur (Effet Foehn / Étage montagnard).} \textit{Exemple :} Expliquer la savane à Ngaoundéré (1100 m) uniquement par la latitude. \textbf{Reflexe :} Citer systématiquement : « L'altitude adoucit les températures et crée une barrière orographique (partage des eaux) qui structure le climat de transition. »
    \item \textbf{Traiter la diversité comme une simple énumération (Inventaire) sans analyse de relations (Système).} \textit{Exemple :} Lister « Forêt au Sud, Savane au Nord » sans dire « Le gradient pluviométrique Sud-Nord, modulé par le relief de l'Adamaoua, commande cette zonation ». \textbf{Reflexe :} Verbes d'analyse : \textbf{Conditionne, Structure, Explique, Détermine, Oppose, Relie}.
    \item \textbf{Erreurs de calcul / Unités dans l'exercice statistique.} \textit{Exemple :} Calculer un taux de variation sans $\times 100$ (donner 1,33 au lieu de 133\%), ou mélanger tonnes et kg, ou oublier l'échelle de temps pour le TMAV. \textbf{Reflexe :} Écrire la formule, remplacer les chiffres \textbf{avec unités}, calculer, \textbf{écrire le résultat avec \% et unité}, phrase d'interprétation.
    \item \textbf{Croquis / Coupe non conformes aux normes cartographiques (0 point figure).} \textit{Exemple :} Pas de titre, pas de légende organisée, pas d'échelle, pas de Nord, couleurs inversées (bleu pour forêt), traits tremblants, pas de repères (villes, fleuves). \textbf{Reflexe :} \textbf{T.O.L.E.S.} : \textbf{T}itre, \textbf{O}rientation (Nord), \textbf{L}égende (hiérarchisée), \textbf{E}chelle, \textbf{S}ource. Soin graphique (règle, crayons de couleur).
\end{enumerate}
\end{attention}

\newpage

\coursec{Exercices}

\courssub{Je vérifie mes connaissances}

\begin{exercice}[QCM : Le relief et le climat]
Parmi les affirmations suivantes, une seule est exacte. La recopier en la complétant.
\begin{enumerate}
\item Le point culminant du Cameroun est le mont \textbf{...} (altitude \textbf{...} m).
\item Le climat équatorial de type guinéen se caractérise par \textbf{...} saisons des pluies et \textbf{...} saisons sèches.
\item La région de l'Extrême-Nord relève du domaine climatique \textbf{...}.
\item Le fleuve \textbf{...} prend sa source sur l'Adamaoua et se jette dans le lac Tchad.
\end{enumerate}
\end{exercice}

\begin{exercice}[Vrai / Faux justifié]
Dire si les affirmations sont vraies ou fausses. Justifier brièvement chaque réponse.
\begin{enumerate}
\item Les sols ferrallitiques sont riches en humus et très fertiles pour l'agriculture intensive.
\item L'Adamaoua constitue une barrière climatique entre le Nord et le Sud du Cameroun.
\item Les populations « Kirdi » désignent un groupe ethnique unique et homogène du Nord.
\item La végétation de la forêt dense humide couvre la totalité du territoire au sud du 5\textsuperscript{e} parallèle.
\end{enumerate}
\end{exercice}

\begin{exercice}[Définitions et concepts]
Définir précisément les termes suivants en géographie biogéographique :
\begin{enumerate}
\item Ensemble biogéographique
\item Écotone
\item Ferrallitisation
\item Groupe humain (ou peuple) au sens démographique et culturel
\end{enumerate}
\end{exercice}

\begin{exercice}[Calcul direct : Densité et taux d'urbanisation]
Le Cameroun compte environ 27 millions d'habitants pour une superficie de 475 442 km\textsuperscript{2}.
\begin{enumerate}
\item Calculer la densité moyenne de population du pays (arrondir à l'entier).
\item Si la population urbaine est estimée à 58 \%, calculer le nombre d'urbains et de ruraux (en millions, arrondi au dixième).
\end{enumerate}
\end{exercice}

\courssub{Je m'entraîne}

\begin{exercice}[Analyse de tableaux pluviométriques]
Le tableau ci-dessous donne les précipitations mensuelles moyennes (en mm) pour trois stations.
\begin{center}
\begin{tabular}{|l|*{12}{c|}}\hline
Mois & J & F & M & A & M & J & J & A & S & O & N & D \\ \hline
Douala & 30 & 50 & 150 & 250 & 350 & 700 & 800 & 700 & 500 & 400 & 150 & 40 \\ \hline
Ngaoundéré & 0 & 0 & 10 & 50 & 150 & 200 & 250 & 300 & 250 & 100 & 10 & 0 \\ \hline
Maroua & 0 & 0 & 0 & 5 & 30 & 80 & 180 & 220 & 150 & 30 & 0 & 0 \\ \hline
\end{tabular}
\end{center}
\begin{enumerate}
\item Calculer le total annuel des précipitations pour chaque station.
\item Identifier le type de climat de chaque station (guinéen, soudanien, sahélien) en justifiant par deux arguments chiffrés.
\item Tracer l'ombrothermique simplifiée de Maroua (courbe des pluies seulement) sur un papier millimétré.
\end{enumerate}
\end{exercice}

\begin{exercice}[Correspondance végétation – sol – climat]
Compléter le tableau de correspondance suivant pour les trois grands ensembles biogéographiques du Cameroun.
\begin{center}
\begin{tabularx}{\linewidth}{|l|X|X|X|}\hline
\textbf{Ensemble} & \textbf{Climat dominant} & \textbf{Végétation climax} & \textbf{Type de sol dominant} \\ \hline
Grand Sud &  &  &  \\ \hline
Hautes terres de l'Ouest / Adamaoua &  &  &  \\ \hline
Grand Nord &  &  &  \\ \hline
\end{tabularx}
\end{center}
\end{exercice}

\begin{exercice}[Analyse de photographies : Milieu et activités]
On observe deux photographies :
\begin{itemize}
\item \textbf{Photo A :} Paysage de plaine alluviale du Wouri (Douala) : mangroves, palétuviers, pirogues, habitations sur pilotis.
\item \textbf{Photo B :} Paysage de l'Adamaoua (Ngaoundéré) : savane herbeuse parsemée d'arbres (dattiers, nérés), troupeau de bovins (race Goudali), case ronde à toit conique.
\end{itemize}
Pour chaque photo, dégager \textbf{deux relations} entre le milieu physique (climat, sol, végétation) et les activités humaines observées (pêche, agriculture, élevage, habitat).
\end{exercice}

\begin{exercice}[Croquis des grands groupes humains]
À partir de la carte muette du Cameroun fournie en annexe (ou tracée à main levée), réaliser un croquis schématique de la répartition des \textbf{quatre grands groupes linguistiques/ethniques} (Soudanais, Bantous, Semi-Bantous, Arabes-Choa/Peuls).
\begin{enumerate}
\item Tracer les limites approximatives des quatre aires.
\item Utiliser quatre couleurs ou hachures différentes.
\item Compléter la légende, le titre, l'échelle et l'orientation.
\end{enumerate}
\end{exercice}

\courssub{Je me prépare au Bac}

\begin{exercice}[Sujet type Bac : Étude de documents – Climat, Végétation et Activités]
\textbf{Documents :}
\begin{itemize}
\item \textbf{Doc 1 :} Carte schématique des climats du Cameroun (Équatorial guinéen, Équatorial camerounais, Tropical soudanien, Tropical sahélien).
\item \textbf{Doc 2 :} Tableau des caractéristiques climatiques de trois stations (voir exercice « Je m'entraîne » n°1).
\item \textbf{Doc 3 :} Photographie aérienne de la région de Maroua (plaine de la Bénoué) montrant des champs de mil/sorgho, des parcelles de coton, des mares asséchées en saison sèche, habitat dispersé.
\end{itemize}
\textbf{Questions :}
\begin{enumerate}
\item \textbf{Doc 1 et 2 :} À l'aide des documents 1 et 2, identifier le type de climat de Douala, Ngaoundéré et Maroua. Justifier chaque réponse par deux éléments climatiques chiffrés prélevés dans le tableau. (6 pts)
\item \textbf{Doc 1 et 3 :} En vous appuyant sur le document 3 et vos connaissances, expliquer l'adaptation des systèmes agraires (cultures, élevage) aux contraintes du climat sahélien dans la plaine de la Bénoué. (5 pts)
\item \textbf{Synthèse :} « La diversité climatique est le facteur premier de la diversité des paysages végétaux et des activités agricoles au Cameroun. » Discuter cette affirmation en structurant votre réponse en deux parties argumentées. (9 pts)
\end{enumerate}
\end{exercice}

\begin{exercice}[Sujet type Bac : Dissertation géographique]
\textbf{Sujet :} \og Montrez que le Cameroun mérite l'appellation d'\textbf{Afrique en miniature}.\fg
\textbf{Consignes :} Rédiger une dissertation géographique structurée comprenant :
\begin{enumerate}
\item Une \textbf{introduction} (accroche, définition du sujet, délimitation spatiale, problématique, annonce du plan).
\item Un \textbf{développement} en \textbf{deux parties équilibrées} (I. Une diversité physique exceptionnelle ; II. Une mosaïque humaine et culturelle riche), chacune comprenant deux paragraphes argumentés avec des exemples précis (reliefs, climats, végétations, groupes humains, langues, religions).
\item Une \textbf{conclusion} (bilan, ouverture sur les défis de l'unité nationale / intégration régionale).
\end{enumerate}
\end{exercice}

\begin{exercice}[Sujet type Probatoire Technique : Construction graphique et Analyse démographique]
\textbf{Données :} Densité de population par région (hab/km\textsuperscript{2}) – Estimations récentes.
\begin{center}
\begin{tabularx}{\linewidth}{|l|X|}\hline
\textbf{Région} & \textbf{Densité} \\ \hline
Ouest & 150 \\ \hline
Littoral (hors Douala) & 95 \\ \hline
Nord-Ouest & 110 \\ \hline
Centre & 35 \\ \hline
Adamaoua & 15 \\ \hline
Nord & 30 \\ \hline
Extrême-Nord & 60 \\ \hline
Est & 8 \\ \hline
Sud & 12 \\ \hline
Sud-Ouest & 80 \\ \hline
\end{tabularx}
\end{center}
\begin{enumerate}
\item Construire un \textbf{histogramme} représentant la densité de population de ces 10 régions (classer les régions par ordre décroissant de densité sur l'axe des abscisses). Soigner le titre, les axes, l'échelle, la légende. (6 pts)
\item Analyser le contraste saisissant entre les \textbf{Hautes Terres de l'Ouest} (Ouest, Nord-Ouest) et l'\textbf{Adamaoua}. Mobiliser au moins \textbf{quatre facteurs explicatifs} (physiques, historiques, économiques). (7 pts)
\item Proposer deux mesures d'aménagement du territoire pour réduire les déséquilibres spatiaux de peuplement mis en évidence. (2 pts)
\end{enumerate}
\end{exercice}

\newpage

\coursec{Corrigés des exercices}

\coursec{Leçon 1 : La diversité des ensembles biogéographiques — Banques d'exercices corrigés}

\begin{corrige}[Exercice 1 — QCM : Le relief et le climat]
\begin{enumerate}
\item Le point culminant du Cameroun est le mont \textbf{Cameroun} (altitude \textbf{4\,095 m}), volcan actif situé dans la région du Sud-Ouest, près de Buéa.
\item Le climat équatorial de type guinéen se caractérise par \textbf{deux} saisons des pluies et \textbf{deux} saisons sèches (petite saison sèche et grande saison sèche alternent avec deux maxima pluviométriques).
\item La région de l'Extrême-Nord relève du domaine climatique \textbf{tropical sahélien} (saison des pluies courte, de juin à septembre, et longue saison sèche).
\item Le fleuve \textbf{Logone} prend sa source sur l'Adamaoua et se jette dans le lac Tchad. (On accepte aussi la \textbf{Bénoué} pour la source sur l'Adamaoua, mais celle-ci se jette dans le Niger ; le Logone est la réponse attendue.)
\end{enumerate}
\textbf{Point de vigilance :} ne pas confondre le mont Cameroun (4\,095 m) avec le mont Oku (3\,011 m), deuxième sommet du pays.
\end{corrige}

\begin{corrige}[Exercice 2 — Vrai / Faux justifié]
\begin{enumerate}
\item \textbf{Faux.} Les sols ferrallitiques sont pauvres en humus : la litière est rapidement décomposée et lessivée sous climat chaud et humide. Ils sont riches en hydroxydes de fer et d'aluminium (d'où leur couleur rouge), acides et peu fertiles ; ils exigent jachère ou fertilisation.
\item \textbf{Vrai.} Le plateau de l'Adamaoua (altitude moyenne 1\,000--1\,100 m) arrête les masses d'air humides venues du sud : il marque la transition entre le domaine équatorial (Sud) et le domaine tropical (Nord), d'où son nom de \og charnière \fg ou \og barrière \fg climatique.
\item \textbf{Faux.} Le terme \og Kirdi \fg (païens, en langue peule) est un terme péjoratif désignant l'ensemble des populations non islamisées du Nord (Mafa, Mofu, Toupouri, Kapsiki, etc.) : ce n'est ni un groupe unique ni homogène.
\item \textbf{Faux.} Le sud du 5\textsuperscript{e} parallèle comprend aussi des mangroves littorales, des forêts-galeries, des savanes (par exemple la savane du plateau Batéké prolongée vers le Sud) et des zones dégradées par l'agriculture ; la forêt dense humide n'y couvre pas la totalité du territoire.
\end{enumerate}
\end{corrige}

\begin{corrige}[Exercice 3 — Définitions et concepts]
\begin{enumerate}
\item \textbf{Ensemble biogéographique :} portion de l'espace terrestre caractérisée par une combinaison originale et cohérente de facteurs physiques (relief, climat, sol, végétation, faune) formant un milieu relativement homogène, distinct des ensembles voisins.
\item \textbf{Écotone :} zone de transition entre deux ensembles biogéographiques, où se mélangent les espèces des deux milieux (exemple : la savane arborée soudanienne, écotone entre la forêt équatoriale et la steppe sahélienne).
\item \textbf{Ferrallitisation :} processus d'altération des roches sous climat chaud et humide : lessivage des éléments solubles (silice, bases) et concentration relative des hydroxydes de fer et d'aluminium, donnant des sols rouges, acides et peu fertiles.
\item \textbf{Groupe humain (peuple) :} ensemble de populations partageant une langue (ou des langues proches), une histoire, des coutumes et des institutions communes, et ayant conscience d'une identité collective (exemples : les Bamiléké, les Bassa, les Peuls).
\end{enumerate}
\end{corrige}

\begin{corrige}[Exercice 4 — Calcul direct : Densité et taux d'urbanisation]
\begin{enumerate}
\item Densité moyenne :
\[ D = \frac{\text{Population}}{\text{Superficie}} = \frac{\num{27000000}}{\SI{475442}{km^2}} \approx 56,8 \text{ hab/km}^2 \]
soit environ \textbf{57 habitants au km\textsuperscript{2}} (arrondi à l'entier).
\item Population urbaine : $\num{27} \times \num{0,58} = \num{15,66} \approx \textbf{15,7 millions d'urbains}$.\\
Population rurale : $\num{27} - \num{15,7} = \textbf{11,3 millions de ruraux}$ (ou $\num{27} \times \num{0,42} = \num{11,34} \approx 11,3$).
\end{enumerate}
\textbf{Vérification :} $15,7 + 11,3 = 27,0$ million : le total est respecté.
\end{corrige}

\begin{corrige}[Exercice 5 — Analyse de tableaux pluviométriques]
\begin{enumerate}
\item Totaux annuels :
\begin{itemize}
\item Douala : $30+50+150+250+350+700+800+700+500+400+150+40 = \textbf{4\,170 mm}$.
\item Ngaoundéré : $0+0+10+50+150+200+250+300+250+100+10+0 = \textbf{1\,320 mm}$.
\item Maroua : $0+0+0+5+30+80+180+220+150+30+0+0 = \textbf{695 mm}$.
\end{itemize}
\item Identification des climats :
\begin{itemize}
\item \textbf{Douala : climat équatorial guinéen} (type côtier). Pluies abondantes toute l'année (4\,170 mm $>$ 2\,000 mm), aucun mois réellement sec (minimum 30 mm en janvier), deux maxima pluviométriques.
\item \textbf{Ngaoundéré : climat tropical soudanien}. Pluies modérées (1\,320 mm, entre 1\,000 et 1\,500 mm), une seule saison des pluies (avril à octobre, maximum en août : 300 mm) et une saison sèche marquée de novembre à mars (0 mm).
\item \textbf{Maroua : climat tropical sahélien}. Pluies faibles (695 mm $<$ 800 mm), saison des pluies très courte (juin à septembre, 4 mois) et longue saison sèche de 8 mois sans pluie significative.
\end{itemize}
\item Courbe ombrothermique de Maroua : en abscisse les 12 mois, en ordonnée les pluies (échelle conseillée : 1 cm pour 50 mm). La courbe reste à 0 de novembre à mars, monte progressivement (5, 30, 80 mm), culmine en août (220 mm), puis redescend (150, 30 mm) et retombe à 0. La forme en cloche étroite traduit l'unicité et la brièveté de la saison des pluies.
\end{enumerate}
\textbf{Point de vigilance :} un climat équatorial n'a jamais de mois sec strict ; un climat tropical présente une saison sèche nette dont la durée croît vers le nord.
\end{corrige}

\begin{corrige}[Exercice 6 — Correspondance végétation – sol – climat]
\begin{center}
\begin{tabularx}{\linewidth}{|l|X|X|X|}\hline
\rowcolor{popblueL}\textbf{Ensemble} & \textbf{Climat dominant} & \textbf{Végétation climax} & \textbf{Type de sol dominant} \\ \hline
Grand Sud & Équatorial (guinéen et camerounais), pluies abondantes toute l'année & Forêt dense humide sempervirente (+ mangrove sur le littoral) & Sols ferrallitiques rouges ou jaunes \\ \hline
Hautes terres de l'Ouest / Adamaoua & Équatorial camerounais d'altitude et tropical soudanien humide & Savane arborée et herbeuse d'altitude, galeries forestières & Sols ferrugineux tropicaux et sols volcaniques fertiles (Ouest) \\ \hline
Grand Nord & Tropical soudanien sec à sahélien & Savane herbeuse à steppique, steppe épineuse (Extrême-Nord) & Sols ferrugineux peu évolués, sols bruns subarides, vertisols (plaines) \\ \hline
\end{tabularx}
\end{center}
\textbf{Point de vigilance :} ne pas écrire \og sols ferrallitiques \fg pour le Nord : la ferrallitisation exige un climat très humide ; au Nord dominent les sols ferrugineux et les vertisols des plaines argileuses.
\end{corrige}

\begin{corrige}[Exercice 7 — Analyse de photographies : Milieu et activités]
\textbf{Photo A (plaine alluviale du Wouri) :}
\begin{enumerate}
\item Le climat équatorial très humide et les eaux saumâtres de l'estuaire permettent la \textbf{mangrove à palétuviers}, dont les poissons et crevettes nourrissent une \textbf{pêche artisanale en pirogue}.
\item Les \textbf{inondations fréquentes} de la plaine alluviale expliquent l'\textbf{habitat sur pilotis}, adapté à la submersion saisonnière du sol.
\end{enumerate}
\textbf{Photo B (Adamaoua) :}
\begin{enumerate}
\item La \textbf{savane herbeuse} (climat tropical soudanien humide, saison sèche qui réduit les mouches tsé-tsé) offre de vastes \textbf{pâturages} favorables à l'\textbf{élevage extensif de bovins} (race Goudali, adaptée au milieu).
\item Les \textbf{matériaux locaux} (terre, paille des graminées de savane) et la saison sèche longue expliquent l'\textbf{habitat en cases rondes à toit conique de chaume}, frais et facile à construire.
\end{enumerate}
\textbf{Point de vigilance :} chaque relation doit relier explicitement un fait physique à une activité humaine ; une simple description des photos ne suffit pas.
\end{corrige}

\begin{corrige}[Exercice 8 — Croquis des grands groupes humains]
\textbf{Fond de carte :} contour simplifié du Cameroun (forme de triangle allongé nord-sud), flèche du nord, échelle indicative (1 cm $\approx$ 100 km), titre : \og Les grands groupes humains du Cameroun \fg.

\textbf{Répartition attendue (quatre couleurs ou hachures) :}
\begin{itemize}
\item \textbf{Soudanais} (hachures horizontales) : Grand Nord (Nord et Extrême-Nord) — Toupouri, Massa, Mousgoum, Kotoko, populations dites \og Kirdi \fg.
\item \textbf{Arabes-Choa et Peuls} (hachures verticales) : frange septentrionale (rives du lac Tchad, Logone-et-Chari) et disséminés dans tout le Nord (Peuls aussi sur l'Adamaoua).
\item \textbf{Semi-Bantous (Grassfields)} (points) : hautes terres de l'Ouest et du Nord-Ouest (Bamiléké, Bamoun, Bamenda, Tikars) et Adamaoua.
\item \textbf{Bantous} (grisé) : tout le Sud forestier et le littoral (Fang-Beti, Ewondo, Douala, Bassa, Bulu, Bakoko, etc.).
\end{itemize}
\textbf{Légende :} quatre cases colorées nommées, plus les villes repères (Yaoundé, Douala, Garoua, Maroua, Ngaoundéré, Bafoussam).

\textbf{Barème indicatif :} fond de carte et orientation (2 pts) ; localisation correcte des quatre aires (4 pts) ; légende, titre, échelle (2 pts) ; propreté (2 pts).
\textbf{Point de vigilance :} les limites sont des zones de transition (fronts pionniers, brassages), non des frontières nettes ; les Peuls débordent vers l'Adamaoua.
\end{corrige}

\begin{corrige}[Exercice 9 — Sujet type Bac : Étude de documents]
\textbf{Question 1 (6 pts — 2 pts par station : 1 pt identification + 1 pt justification chiffrée).}
\begin{itemize}
\item \textbf{Douala : climat équatorial guinéen.} Total annuel 4\,170 mm ($>$ 2\,000 mm) ; pluies tous les mois de l'année (minimum 30 mm en janvier) avec deux maxima.
\item \textbf{Ngaoundéré : climat tropical soudanien.} Total 1\,320 mm ; une seule saison des pluies d'avril à octobre (maximum 300 mm en août) et 5 mois secs (novembre--mars, 0 à 10 mm).
\item \textbf{Maroua : climat tropical sahélien.} Total faible : 695 mm ; saison des pluies réduite à 4 mois (juin--septembre) et 8 mois de sécheresse.
\end{itemize}
\textbf{Question 2 (5 pts).} Adaptations au climat sahélien de la plaine de la Bénoué :
\begin{itemize}
\item Choix de \textbf{cultures à cycle court et résistantes à la sécheresse} : mil, sorgho, niébé, arachide (1,5 pt) ;
\item \textbf{Coton} cultivé en saison des pluies, adapté aux vertisols et à la saison sèche longue qui facilite la récolte (1 pt) ;
\item \textbf{Élevage extensif et transhumance} : les troupeaux se déplacent vers les points d'eau (mares, Bénoué) en saison sèche (1,5 pt) ;
\item \textbf{Calendrier agricole calé sur la courte saison des pluies} (juin--septembre) et habitat dispersé près des champs et des points d'eau (1 pt).
\end{itemize}
\textbf{Question 3 (9 pts) — Synthèse en deux parties.}
\begin{itemize}
\item \textbf{I. Oui, le climat commande les paysages et les activités (4,5 pts).} Du sud au nord, la diminution et la concentration des pluies déterminent la succession forêt dense $\rightarrow$ savane arborée $\rightarrow$ savane herbeuse $\rightarrow$ steppe ; à chaque étage correspond un système agraire : cultures pérennes industrielles (cacao, café, hévéa, bananes) au Sud humide ; cultures vivrières et élevage sur l'Adamaoua ; mil, sorgho, coton et transhumance au Nord sahélien.
\item \textbf{II. Mais d'autres facteurs interviennent (4,5 pts).} Le relief (altitude de l'Adamaoua et des hautes terres créant des climats d'altitude), la nature des sols (sols volcaniques fertiles de l'Ouest, vertisols du Nord), l'histoire et les sociétés (densités Bamiléké liées à des facteurs culturels et historiques, islamisation du Nord), les politiques d'aménagement modifient aussi paysages et activités. Le climat est donc un facteur premier mais non exclusif.
\end{itemize}
\textbf{Points de vigilance :} citer les chiffres du tableau (exigence du sujet) ; structurer la synthèse en deux parties équilibrées ; conclure par un jugement nuancé.
\end{corrige}

\begin{corrige}[Exercice 10 — Sujet type Bac : Dissertation géographique]
\textbf{Introduction (attendue).} Accroche : le Cameroun est souvent présenté comme l'\og Afrique en miniature \fg. Définition : cette formule signifie que le pays résume, sur ses \SI{475442}{km^2}, l'essentiel des milieux et des peuples du continent. Délimitation : le Cameroun, de 2° à 13° de latitude nord. Problématique : en quoi la diversité physique et humaine du Cameroun reproduit-elle celle de l'Afrique entière ? Annonce du plan : nous montrerons d'abord une diversité physique exceptionnelle, puis une mosaïque humaine et culturelle remarquable.

\textbf{I. Une diversité physique exceptionnelle.}
\begin{itemize}
\item \textbf{§1 — Reliefs et climats :} du niveau de la mer au mont Cameroun (4\,095 m), plaines côtières, plateaux du Sud, hautes terres volcaniques de l'Ouest, plateau de l'Adamaoua, plaines du Nord et monts Mandara ; quatre types de climats s'échelonnent de l'équatorial guinéen (4\,170 mm à Douala) au tropical sahélien (695 mm à Maroua), comme du golfe de Guinée au Sahel à l'échelle du continent.
\item \textbf{§2 — Végétations et sols :} mangroves, forêt dense humide, savanes arborées puis herbeuses, steppes épineuses du lac Tchad : on retrouve tous les grands paysages africains ; sols ferrallitiques au Sud, ferrugineux et vertisols au Nord.
\end{itemize}
\textbf{II. Une mosaïque humaine et culturelle riche.}
\begin{itemize}
\item \textbf{§1 — Les grands groupes humains :} Bantous du Sud forestier, Semi-Bantous des Grassfields, Soudanais et Peuls du Nord, Arabes-Choa des rives du Tchad : plus de 250 ethnies, à l'image du brassage continental.
\item \textbf{§2 — Langues et religions :} deux langues officielles (français et anglais, héritage colonial triple : allemand, français, britannique), plus de 250 langues nationales ; christianisme au Sud, islam au Nord, religions traditionnelles partout.
\end{itemize}
\textbf{Conclusion (attendue).} Bilan : par ses milieux comme par ses hommes, le Cameroun condense l'Afrique. Ouverture : cette diversité est un atout mais pose le défi de l'unité nationale et de l'intégration régionale (CEMAC).

\textbf{Barème indicatif :} introduction (3 pts) ; développement I (6 pts) ; développement II (6 pts) ; conclusion (3 pts) ; langue et présentation (2 pts).
\textbf{Points de vigilance :} exemples précis et chiffrés exigés ; deux parties équilibrées ; éviter le catalogue sans argumentation.
\end{corrige}

\begin{corrige}[Exercice 11 — Sujet type Probatoire Technique : Construction graphique et analyse démographique]
\textbf{Question 1 (6 pts).} Classement décroissant : Ouest (150), Nord-Ouest (110), Littoral (95), Sud-Ouest (80), Extrême-Nord (60), Centre (35), Nord (30), Adamaoua (15), Sud (12), Est (8). Histogramme : abscisses = régions dans cet ordre ; ordonnées = densité (échelle conseillée : 1 cm pour 20 hab/km\textsuperscript{2}, maximum 150) ; barres de largeur égale, jointives ; titre : \og Densité de population par région du Cameroun (hab/km\textsuperscript{2}) \fg ; axes nommés et gradués.

\begin{popfigure}
\begin{center}
\begin{tikzpicture}[x=0.95cm,y=0.035cm]
\draw[->] (0,0) -- (10.8,0) node[below left, font=\scriptsize]{Régions};
\draw[->] (0,0) -- (0,165) node[above left, font=\scriptsize]{hab/km$^2$};
\foreach \y in {0,30,60,90,120,150}{\draw (0,\y) -- (-0.08,\y) node[left, font=\scriptsize]{\y};}
\foreach \x/\h/\n in {0.5/150/Ouest, 1.5/110/N.-Ouest, 2.5/95/Littoral, 3.5/80/S.-Ouest, 4.5/60/Extr.-Nord, 5.5/35/Centre, 6.5/30/Nord, 7.5/15/Adamaoua, 8.5/12/Sud, 9.5/8/Est}{
\fill[popblue!60] (\x-0.4,0) rectangle (\x+0.4,\h);
\node[below, font=\tiny, rotate=40, anchor=east] at (\x,0) {\n};
\node[above, font=\tiny] at (\x,\h) {\h};}
\end{tikzpicture}
\end{center}
\legende{Histogramme des densités régionales (hab/km\textsuperscript{2}), classées par ordre décroissant. Barème : classement correct (1 pt), échelle et axes (2 pts), barres exactes (2 pts), titre et présentation (1 pt).}
\end{popfigure}

\textbf{Question 2 (7 pts).} Contraste : Ouest 150 hab/km\textsuperscript{2} et Nord-Ouest 110 contre Adamaoua 15, soit un rapport d'environ 10 à 1 ($150/15 = 10$). Facteurs explicatifs :
\begin{enumerate}
\item \textbf{Facteur physique — sols :} sols volcaniques très fertiles des hautes terres de l'Ouest ; sols ferrugineux pauvres et latéritiques de l'Adamaoua (1,5 pt).
\item \textbf{Facteur physique — climat et santé :} climat d'altitude sain à l'Ouest ; Adamaoua longtemps infestée par la mouche tsé-tsé et les endémies (1,5 pt).
\item \textbf{Facteur historique :} ancienneté du peuplement et systèmes fonciers denses chez les Bamiléké ; l'Adamaoua, zone de conquête peule tardive et de raids, est restée peu peuplée (2 pts).
\item \textbf{Facteur économique :} agriculture commerciale intensive (café, maraîchage), petites propriétés, proximité de Douala et des marchés à l'Ouest ; élevage extensif et enclavement (éloignement des voies de communication) sur l'Adamaoua (2 pts).
\end{enumerate}
\textbf{Question 3 (2 pts).} Deux mesures possibles (1 pt chacune) :
\begin{itemize}
\item \textbf{Désenclaver} les régions faiblement peuplées : achever et entretenir les routes et le chemin de fer Transcamerounais vers l'Adamaoua, le Nord et l'Est ;
\item \textbf{Décentraliser} les activités et services (universités, hôpitaux, zones industrielles) vers les villes moyennes (Ngaoundéré, Bertoua) pour fixer les populations et freiner l'exode vers Douala et Yaoundé.
\end{itemize}
\textbf{Points de vigilance :} les barres d'un histogramme sont jointives ; tout facteur doit être relié explicitement à la densité ; citer le rapport chiffré entre les deux espaces.
\end{corrige}

\newpage

\coursec{Fiche bilan}

\begin{popfigure}
\begin{tikzpicture}[
  mindmap,
  concept color=popblue,
  level 1 concept/.append style={level distance=4.5cm, sibling angle=60},
  level 2 concept/.append style={level distance=3cm, sibling angle=45},
  every node/.style={font=\small, text=white},
  grow cyclic,
  text width=2.8cm,
  align=center
]
\node[concept, concept color=popblue, align=center]{Diversité des ensembles\\biogéographiques\\Cameroun}
  child[concept color=poporange] {node[concept] {Reliefs \& Climats}
    child {node[concept, concept color=poporange!70!popblue, text width=2.5cm] {Plaines côtières, Plateau Sud, Hauts-Plateaux, Monts, Plaines du Nord}}
    child {node[concept, concept color=poporange!70!popblue, text width=2.5cm] {Climats : équatorial, tropical, soudano-sahélien}}}
  child[concept color=popgreen] {node[concept] {Végétation \& Sols}
    child {node[concept, concept color=popgreen!70!popblue, text width=2.5cm] {Forêt dense, savane boisée, steppe, mangrove}}
    child {node[concept, concept color=popgreen!70!popblue, text width=2.5cm] {Sols : ferrallitiques, ferrugineux, hydromorphes, vertisols}}}
  child[concept color=poppurple] {node[concept] {Ensembles biogéo.}
    child {node[concept, concept color=poppurple!70!popblue, text width=2.5cm] {Domaine forestier (Sud)}}
    child {node[concept, concept color=poppurple!70!popblue, text width=2.5cm] {Domaine savanicole (Centre-Nord)}}
    child {node[concept, concept color=poppurple!70!popblue, text width=2.5cm] {Domaine steppique (Extrême-Nord)}}}
  child[concept color=poppink] {node[concept] {Groupes humains}
    child {node[concept, concept color=poppink!70!popblue, text width=2.5cm] {Bantous (Sud, Est, Centre)}}
    child {node[concept, concept color=poppink!70!popblue, text width=2.5cm] {Semi-Bantous (Ouest, Nord-Ouest)}}
    child {node[concept, concept color=poppink!70!popblue, text width=2.5cm] {Soudanais \& Sahéliens (Nord, Extrême-Nord)}}}
  child[concept color=popgold] {node[concept] {Unité dans diversité}
    child {node[concept, concept color=popgold!70!popblue, text width=2.5cm] {Atouts : ressources, tourisme, culture}}
    child {node[concept, concept color=popgold!70!popblue, text width=2.5cm] {Défis : aménagement, intégration, désertification}}};
\end{tikzpicture}
\legende{Carte mentale : La diversité des ensembles biogéographiques du Cameroun}
\end{popfigure}

\begin{bilan}
\textbf{1. Définitions clés}
\begin{itemize}
  \item \cle{Afrique en miniature} : Surnom du Cameroun justifié par la présence de presque tous les types de climats, reliefs, végétations, sols et groupes humains du continent africain sur son territoire (\SI{475 442}{km^2}).
  \item \cle{Ensemble biogéographique} : Unité spatiale où le climat, le relief, la végétation, les sols et l'occupation humaine forment un tout cohérent et interdépendant.
  \item \cle{Domaine bioclimatique} : Zone définie par la combinaison température/précipitations (ex: domaine équatorial guinéen, domaine tropical soudanien, domaine sahélien).
\end{itemize}

\textbf{2. Données de référence à connaître (tableau de synthèse)}
\begin{center}
\begin{tabularx}{\linewidth}{|l|X|X|X|}\hline
\rowcolor{popblueL}
\textbf{Domaine} & \textbf{Climat / Relief} & \textbf{Végétation / Sols} & \textbf{Groupes humains dominants} \\ \hline
\textbf{Forestier (Sud, Est, Centre, Littoral, Sud-Ouest)} & Équatorial (4 saisons / 2 saisons), Plaines côtières, Plateau Sud, Massifs volcaniques & Forêt dense humide, mangrove ; Sols ferrallitiques (lessivés, pauvres en bases) & Bantous (Bassa, Beti, Boulou, Fang, Douala, Bassa) \\ \hline
\textbf{Savanicole / Hauts-Plateaux (Ouest, Nord-Ouest, Adamaoua, Nord)} & Tropical de transition / d'altitude (2 saisons), Hauts-Plateaux (\SI{1000-2500}{m}), Adamaoua & Savane boisée, galerie forestière ; Sols ferrugineux lessivés, ferrallitiques lessivés & Semi-Bantous (Bamiléké, Bamoun, Tikar), Groupes de l'Adamaoua (Fulbe, Mbum) \\ \hline
\textbf{Steppique / Sahélien (Extrême-Nord, partie Nord Bénoué)} & Soudano-sahélien / Sahélien (1 saison des pluies courte), Plaines du Nord (Diamaré, Logone, Bénoué) & Steppe à \emph{Acacia}, savane arbustive ; Sols ferrugineux tropicaux, vertisols (argile gonflante), sols hydromorphes (yaérés) & Peuls (Fulbe), Mousgoum, Kotoko, Arabes Choa, Kanouri \\ \hline
\end{tabularx}
\end{center}

\textbf{3. Méthodes express}
\begin{enumerate}
  \item \textbf{Caractériser un ensemble biogéographique} : Citer le domaine $\rightarrow$ Climat (type, pluviométrie, T°) $\rightarrow$ Relief $\rightarrow$ Végétation + Sol type $\rightarrow$ Activités humaines / Groupes ethniques.
  \item \textbf{Justifier « Afrique en miniature »} : Lister au moins 3 types de climats, 3 grands types de reliefs, 3 types de végétation, 3 types de sols, 3 familles linguistiques.
  \item \textbf{Analyser l'unité dans la diversité} : Identifier les flux (humains, marchandises) $\rightarrow$ Citer les infrastructures structurantes (Transcam, routes, ports) $\rightarrow$ Citer les institutions (bilingue, décentralisation) $\rightarrow$ Citer les défis (insécurité Nord/NO/SO, pression foncière, désertification).
\end{enumerate}
\end{bilan}

\begin{aretenir}[Checklist avant le Bac]
\begin{itemize}
  \item $\square$ Savoir définir « Afrique en miniature » avec 3 arguments physiques et 2 arguments humains précis.
  \item $\square$ Localiser sur carte muette : 5 grands reliefs, 3 domaines climatiques, 3 ensembles biogéographiques, 4 grands groupes humains.
  \item $\square$ Associer chaque domaine biogéographique à son couple Climat/Végétation/Sol dominant (tableau ci-dessus).
  \item $\square$ Expliquer la répartition des groupes humains (Bantous, Semi-Bantous, Soudanais, Sahéliens) en lien avec l'histoire migratoire et le milieu.
  \item $\square$ Citer 2 atouts économiques par ensemble biogéographique (ex: forêt/bois/cacao ; savane/élevage/coton ; steppe/riz/élevage/onions).
  \item $\square$ Identifier 2 défis majeurs par ensemble (ex: déforestation/érosion ; feux de brousse/exode rural ; désertification/insécurité alimentaire).
  \item $\square$ Expliquer le rôle des axes de communication (Transcam, Route Nationale 1, Port de Douala/Kribi) dans l'intégration nationale.
  \item $\square$ Différencier sol ferrallitique (Sud, lessivé, acide) et sol ferrugineux (Nord, moins lessivé, plus fertile) et vertisol (plaines alluviales, argile gonflante).
  \item $\square$ Justifier la concentration de la population au Sud et dans les Hauts-Plateaux (climat favorable, sols, histoire, urbanisation).
  \item $\square$ Rédiger une conclusion structurée : « Le Cameroun, par sa diversité..., fait face au défi de l'unité nationale par... »
\end{itemize}
\end{aretenir}

\findecours

\end{document}
